% Étude de fonctions — Mathématiques Tle A — Cours signé OnBuch+
% Programme officiel MINESEC. Compiler avec : tectonic MA08-etude-de-fonctions-a.tex
\newcommand{\DOCMATIERE}{Mathématiques}
\newcommand{\DOCNIVEAU}{Tle A}
\newcommand{\DOCMODULE}{Relations et opérations fondamentales dans l'ensemble des nombres réels}
\newcommand{\DOCLECON}{A08}
\newcommand{\DOCTITRE}{Étude de fonctions}
% OnBuch+ — préambule « cours signés » ENRICHI (compilé avec tectonic / XeLaTeX)
% Système visuel « Pop » d'OnBuch+ : fond crème (jamais blanc), bordures encre
% épaisses, ombres « dures » décalées, titres Archivo Black en capitales.
% Version MATHÉMATIQUES : graphes (pgfplots/tikz), boîtes pédagogiques
% (définition, propriété, méthode, attention, le savais-tu, exercice résolu,
% exercice, TP, bilan), page de garde et signature OnBuch+ ancrée en bas.
%
% Macros attendues dans main.tex AVANT \input{preamble} :
%   \DOCMATIERE  \DOCNIVEAU  \DOCMODULE  \DOCLECON (numéro)  \DOCTITRE
\documentclass[11pt]{article}

\usepackage{fontspec}
\usepackage[french]{babel}
\usepackage[a4paper,margin=2cm,top=3.4cm,bottom=2.8cm,headheight=1.4cm,footskip=1.3cm]{geometry}
\usepackage{amsmath,amssymb,amsfonts}
\usepackage{mathtools} % \xmapsto, \coloneqq... souvent utilisés par les modèles
\usepackage{cancel} % simplifications barrées
% \abs{z} (module, valeur absolue) et \norm{u} (norme), utilisés par les modèles
\providecommand{\abs}{}\let\abs\relax\DeclarePairedDelimiter{\abs}{\lvert}{\rvert}
\providecommand{\norm}{}\let\norm\relax\DeclarePairedDelimiter{\norm}{\lVert}{\rVert}
\usepackage[table]{xcolor} % [table] : fournit aussi \rowcolors (alternance de lignes)
\usepackage{pagecolor}
\usepackage{tikz}
\usetikzlibrary{arrows.meta,positioning,calc,shapes.geometric,shapes.misc,shapes.arrows,decorations.pathmorphing,decorations.markings,patterns,shadows,mindmap,trees,fit,backgrounds,matrix,intersections}
\usepackage{pgfplots}
\usepackage{pgf-pie} % diagrammes circulaires \pie (statistiques)
\pgfplotsset{compat=1.18}
\usepgfplotslibrary{fillbetween,groupplots} % aires entre courbes (calcul intégral)
\usepackage{enumitem}
\usepackage{fancyhdr}
% [most] charge toutes les bibliothèques tcolorbox (listings, minted, raster,
% documentation...) et gonfle inutilement la mémoire TeX sur un document long
% (~300 boîtes) : on ne charge que ce qui est réellement utilisé.
\usepackage[skins,breakable]{tcolorbox}
\usepackage{graphicx}
\usepackage{colortbl}
\usepackage{booktabs}
\usepackage{tabularx}
\usepackage{diagbox} % case d'en-tête diagonale des tableaux à double entrée
% Colonnes extensibles centrée / gauche / droite (C, L, R), souvent utilisées par les modèles
\newcolumntype{C}{>{\centering\arraybackslash}X}
\newcolumntype{L}{>{\raggedright\arraybackslash}X}
\newcolumntype{R}{>{\raggedleft\arraybackslash}X}
\usepackage{array}
\usepackage{multicol}
\usepackage{siunitx}
% parse-numbers=false : accepte aussi \SI{2,5 \times 10^{-3}}{...}
\sisetup{locale=FR,output-decimal-marker={,},group-minimum-digits=4,parse-numbers=false}
\DeclareSIUnit\ohm{\ensuremath{\symup{\Omega}}} % U+2126 absent de la police
\usepackage{qrcode}
% \degree : commande fréquemment utilisée par les modèles pour un angle ou une
% température en dehors de \SI{}{} (non fournie par les paquets chargés).
\providecommand{\degree}{^{\circ}}
% \qed (fin de démonstration), utilisé par les modèles sans amsthm
\providecommand{\qed}{\hfill$\square$}
% \barre : double barre des valeurs interdites dans un tableau de variations (tkz-tab)
\providecommand{\barre}{\ensuremath{\|}}
% environnement proof (amsthm non chargé) utilisé par les modèles
\ifdefined\proof\else\newenvironment{proof}[1][Démonstration]{\par\noindent\textit{#1.}\ }{\qed\par}\fi
\usepackage{lastpage}
\usepackage{hyperref}
\hypersetup{hidelinks}

% ============================================================
% Palette « Pop » OnBuch+ — mêmes hex que la classe P (plus_ui.dart)
% ============================================================
\definecolor{poporange}{HTML}{F59321}
\definecolor{popdark}{HTML}{E07A0C}
\definecolor{popcream}{HTML}{FFF6E8}
\definecolor{popink}{HTML}{1C1714}
\definecolor{poppurple}{HTML}{7A5AE0}
\definecolor{popgreen}{HTML}{2E9E6A}
\definecolor{poppink}{HTML}{E0457B}
\definecolor{popblue}{HTML}{2F7DE1}
\definecolor{popgold}{HTML}{C98A12}
\definecolor{popmuted}{HTML}{8A7F76}
\definecolor{popline}{HTML}{EADFD2}
\definecolor{popgray}{HTML}{8A7F76}
\colorlet{popgrey}{popgray}
% Alias tolérants (le modèle nomme parfois les couleurs différemment)
\colorlet{popwhite}{white}
\colorlet{popblack}{popink}
\colorlet{popred}{poppink}
\colorlet{popyellow}{popgold}
% Teintes claires pour fonds de tableaux / zones
\colorlet{popblueL}{popblue!10!white}
\colorlet{poporangeL}{poporange!14!white}
\colorlet{popgreenL}{popgreen!12!white}
\colorlet{poppurpleL}{poppurple!12!white}
\colorlet{poppinkL}{poppink!12!white}
\colorlet{popgoldL}{popgold!14!white}

\pagecolor{popcream}

% ============================================================
% Typographie
% ============================================================
\defaultfontfeatures{Ligatures=TeX}
\setmainfont{PlusJakartaSans-Medium.ttf}[Path=../../../fonts/,BoldFont=PlusJakartaSans-Bold.ttf]
% Maths en sans-serif (Fira Math) avec chiffres et lettres droites de Plus
% Jakarta, pour que les formules se fondent dans le texte.
\usepackage[math-style=ISO,bold-style=ISO]{unicode-math}
\setmathfont{FiraMath-Regular.otf}[Path=../../../fonts/]
\setmathfont{PlusJakartaSans-Medium.ttf}[Path=../../../fonts/,range={up/num,up/{Latin,latin},\mathup}]
\usepackage{icomma} % virgule décimale sans espace en mode math
% Caractères mathématiques Unicode absents des polices de texte (Plus Jakarta,
% Archivo Black) : rendus par leur équivalent mathématique, en texte comme en
% formule — sinon ils s'affichent en carré vide (ex. « Arithmétique dans ℤ »).
\usepackage{newunicodechar}
\newunicodechar{ℝ}{\ensuremath{\mathbb{R}}}
\newunicodechar{ℤ}{\ensuremath{\mathbb{Z}}}
\newunicodechar{ℂ}{\ensuremath{\mathbb{C}}}
\newunicodechar{ℕ}{\ensuremath{\mathbb{N}}}
\newunicodechar{ℚ}{\ensuremath{\mathbb{Q}}}
\newunicodechar{𝔻}{\ensuremath{\mathbb{D}}}
\newunicodechar{α}{\ensuremath{\alpha}}
\newunicodechar{β}{\ensuremath{\beta}}
\newunicodechar{γ}{\ensuremath{\gamma}}
\newunicodechar{δ}{\ensuremath{\delta}}
\newunicodechar{ε}{\ensuremath{\varepsilon}}
\newunicodechar{θ}{\ensuremath{\theta}}
\newunicodechar{λ}{\ensuremath{\lambda}}
\newunicodechar{μ}{\ensuremath{\mu}}
\newunicodechar{σ}{\ensuremath{\sigma}}
\newunicodechar{τ}{\ensuremath{\tau}}
\newunicodechar{φ}{\ensuremath{\varphi}}
\newunicodechar{ψ}{\ensuremath{\psi}}
\newunicodechar{ω}{\ensuremath{\omega}}
\newunicodechar{Σ}{\ensuremath{\Sigma}}
% majuscules produites par \MakeUppercase dans les titres (α-aminés -> Α)
\newunicodechar{Α}{\ensuremath{\alpha}}
\newunicodechar{Β}{\ensuremath{\beta}}
\newunicodechar{Γ}{\ensuremath{\Gamma}}
\newunicodechar{Δ}{\ensuremath{\Delta}}
\newunicodechar{Ε}{\ensuremath{\varepsilon}}
\newunicodechar{Θ}{\ensuremath{\Theta}}
\newunicodechar{Λ}{\ensuremath{\Lambda}}
\newunicodechar{Μ}{\ensuremath{\mu}}
\newunicodechar{Τ}{\ensuremath{\tau}}
\newunicodechar{Φ}{\ensuremath{\Phi}}
\newunicodechar{Ψ}{\ensuremath{\Psi}}
\newunicodechar{Ω}{\ensuremath{\Omega}}
\newunicodechar{↦}{\ensuremath{\mapsto}}
\newunicodechar{∈}{\ensuremath{\in}}
\newunicodechar{∉}{\ensuremath{\notin}}
\newunicodechar{∧}{\ensuremath{\wedge}}
\newunicodechar{⇔}{\ensuremath{\Leftrightarrow}}
\newunicodechar{⟺}{\ensuremath{\Longleftrightarrow}}
\newunicodechar{⇒}{\ensuremath{\Rightarrow}}
\newunicodechar{⟹}{\ensuremath{\Longrightarrow}}
\newunicodechar{∘}{\ensuremath{\circ}}
\newunicodechar{∪}{\ensuremath{\cup}}
\newunicodechar{∩}{\ensuremath{\cap}}
\newunicodechar{≡}{\ensuremath{\equiv}}
\newunicodechar{⊂}{\ensuremath{\subset}}
\newunicodechar{⊥}{\ensuremath{\perp}}
\newunicodechar{∃}{\ensuremath{\exists}}
\newunicodechar{∀}{\ensuremath{\forall}}
\newunicodechar{∛}{\ensuremath{\sqrt[3]{\,}}}
\newunicodechar{∜}{\ensuremath{\sqrt[4]{\,}}}
\newunicodechar{⃗}{\ensuremath{{}^{\to}}}
\newunicodechar{ⁿ}{\ifmmode^{n}\else\textsuperscript{\ensuremath{n}}\fi}
\newunicodechar{ˣ}{\ifmmode^{x}\else\textsuperscript{\ensuremath{x}}\fi}
\newunicodechar{ᵇ}{\ifmmode^{b}\else\textsuperscript{\ensuremath{b}}\fi}
\newunicodechar{ʳ}{\ifmmode^{r}\else\textsuperscript{\ensuremath{r}}\fi}
\newunicodechar{ᵘ}{\ifmmode^{u}\else\textsuperscript{\ensuremath{u}}\fi}
\newunicodechar{ʸ}{\ifmmode^{y}\else\textsuperscript{\ensuremath{y}}\fi}
\newunicodechar{ᵃ}{\ifmmode^{a}\else\textsuperscript{\ensuremath{a}}\fi}
\newunicodechar{ᵉ}{\ifmmode^{e}\else\textsuperscript{\ensuremath{e}}\fi}
\newunicodechar{ᵏ}{\ifmmode^{k}\else\textsuperscript{\ensuremath{k}}\fi}
\newunicodechar{ᵖ}{\ifmmode^{p}\else\textsuperscript{\ensuremath{p}}\fi}
\newunicodechar{ᴺ}{\ifmmode^{N}\else\textsuperscript{\ensuremath{N}}\fi}
\newunicodechar{ᶜ}{\ifmmode^{c}\else\textsuperscript{\ensuremath{c}}\fi}
\newunicodechar{ʰ}{\ifmmode^{h}\else\textsuperscript{\ensuremath{h}}\fi}
\newunicodechar{ᵠ}{\ifmmode^{q}\else\textsuperscript{\ensuremath{q}}\fi}
\newunicodechar{ₙ}{\ifmmode_{n}\else\textsubscript{\ensuremath{n}}\fi}
\newunicodechar{ₐ}{\ifmmode_{a}\else\textsubscript{\ensuremath{a}}\fi}
\newunicodechar{ᵢ}{\ifmmode_{i}\else\textsubscript{\ensuremath{i}}\fi}
\newunicodechar{ᵣ}{\ifmmode_{r}\else\textsubscript{\ensuremath{r}}\fi}
\newunicodechar{ₖ}{\ifmmode_{k}\else\textsubscript{\ensuremath{k}}\fi}
\newunicodechar{ᵦ}{\ifmmode_{\beta}\else\textsubscript{\ensuremath{\beta}}\fi}
\newunicodechar{ₓ}{\ifmmode_{x}\else\textsubscript{\ensuremath{x}}\fi}
\newfontfamily\popdisplay{ArchivoBlack-Regular.ttf}[Path=../../../fonts/]
\newfontfamily\poplogo{SpaceGrotesk-Bold.ttf}[Path=../../../fonts/]
\newfontfamily\popsemi{PlusJakartaSans-SemiBold.ttf}[Path=../../../fonts/]
\newfontfamily\popxbold{PlusJakartaSans-ExtraBold.ttf}[Path=../../../fonts/]
\newfontfamily\popxboldls{PlusJakartaSans-ExtraBold.ttf}[Path=../../../fonts/,LetterSpace=3.0]

\setlist[enumerate]{leftmargin=1.7em,itemsep=5pt,topsep=4pt}
\setlist[itemize]{leftmargin=1.7em,itemsep=4pt,topsep=4pt,label={\color{poporange}\textbullet}}
\setlength{\parindent}{0pt}
\setlength{\parskip}{5pt}
\renewcommand{\arraystretch}{1.25}

% pgfplots : style Pop par défaut
\pgfplotsset{
  every axis/.append style={axis line style={popink,line width=1pt},tick style={popink},
    grid style={popline,line width=0.5pt},label style={font=\small},tick label style={font=\footnotesize}},
  popcurve/.style={poporange,line width=1.8pt,no markers,smooth},
}
% Même style disponible dans un \draw TikZ simple (hors pgfplots)
\tikzset{popcurve/.style={poporange,line width=1.8pt}}

% ============================================================
% Pill / squiggle
% ============================================================
\newcommand{\poppill}[3]{%
  \tikz[baseline=-0.55ex]\node[draw=popink,line width=1.2pt,rounded corners=2.6pt,%
    fill=#2,inner xsep=6.5pt,inner ysep=3pt]{%
    \popxboldls\fontsize{7}{7}\selectfont\color{#3}\MakeUppercase{#1}};%
}
\newcommand{\popsquiggle}[1]{%
  \tikz[baseline=-0.4ex]{
    \draw[#1,line width=2.6pt,line cap=round]
      plot [smooth] coordinates {(0,0.09) (0.34,0.01) (0.68,0.21) (1.02,0.01) (1.36,0.10)};
  }%
}
% Mot-clé surligné (marqueur orange sous le texte)
\newcommand{\cle}[1]{{\popxbold\color{popdark}#1}}

% ============================================================
% En-tête / pied
% ============================================================
\pagestyle{fancy}
\fancyhf{}
\renewcommand{\headrulewidth}{0pt}
\renewcommand{\footrulewidth}{0pt}
\lhead{\raisebox{-0.15ex}{\poplogo\fontsize{13}{13}\selectfont\color{popink}OnBuch\color{poporange}+}}
\rhead{\poppill{\DOCMATIERE\ \textbullet\ \DOCNIVEAU\ \textbullet\ Leçon \DOCLECON}{white}{popink}}
\lfoot{\popxbold\fontsize{7.6}{7.6}\selectfont\color{popmuted}\MakeUppercase{Cours signé OnBuch+}\ \textbullet\ {\popsemi\DOCTITRE}}
\rfoot{\tikz[baseline=-0.6ex]\node[rounded corners=8.5pt,fill=popink,minimum height=17pt,minimum width=17pt,inner xsep=5pt,inner sep=0]{\popxbold\fontsize{8}{8}\selectfont\color{popcream}\thepage\,/\,\pageref*{LastPage}};}

% ============================================================
% Titres
% ============================================================
\newcommand{\chaptitre}[2]{%
  \begin{tcolorbox}[enhanced,colback=poporange,colframe=popink,boxrule=2.6pt,arc=11pt,
      drop shadow=popink,left=15pt,right=15pt,top=12pt,bottom=11pt,before skip=3pt,after skip=18pt]
    \poppill{#2}{popink}{popcream}\par\vspace{9pt}
    {\raggedright\hyphenpenalty=10000\exhyphenpenalty=10000\popdisplay\fontsize{22}{24}\selectfont\color{popcream}\MakeUppercase{#1}\par}
  \end{tcolorbox}
}
% Section numérotée : pastille orange avec le numéro + titre Archivo
\newcounter{popsec}
\newcommand{\coursec}[1]{%
  \stepcounter{popsec}\setcounter{popsub}{0}%
  \par\vspace{16pt}\noindent
  \begin{minipage}{\linewidth}
  \tikz[baseline=-0.7ex]\node[draw=popink,line width=1.6pt,fill=poporange,rounded corners=4pt,minimum size=22pt,inner sep=0]{\popdisplay\fontsize{12}{12}\selectfont\color{popcream}\thepopsec};%
  \hspace{8pt}{\popdisplay\fontsize{15}{17}\selectfont\color{popink}\MakeUppercase{#1}}\par
  \vspace{4pt}\noindent{\color{poporange}\rule{36pt}{3.6pt}}
  \end{minipage}\par\nopagebreak\vspace{8pt}
}
\newcounter{popsub}
\newcommand{\courssub}[1]{%
  \stepcounter{popsub}%
  \par\vspace{9pt}\noindent{\popxbold\fontsize{11.5}{13}\selectfont\color{popink}{\color{poporange}\thepopsec.\thepopsub}\hspace{6pt}#1}\par\nopagebreak\vspace{4pt}
}

% ============================================================
% Boîtes pédagogiques — même recette : fond blanc, bordure épaisse colorée,
% ombre dure encre, badge plein. Titre optionnel [#1].
% ============================================================
\tcbset{popbase/.style={breakable,enhanced,colback=white,boxrule=2.3pt,arc=9pt,
  shadow={3pt}{-3pt}{0pt}{popink},left=14pt,right=14pt,top=11pt,bottom=11pt,before skip=11pt,after skip=15pt,
  fontupper=\popsemi\fontsize{10.6}{14.5}\selectfont\color{popink},
  fontlower=\popsemi\fontsize{10.6}{14.5}\selectfont\color{popink}}}
% Coche dessinée (le glyphe ✓ n'existe pas dans les polices du document)
\newcommand{\popcheck}[1]{\tikz[baseline=-0.5ex]\draw[#1,line width=1.6pt,line cap=round,line join=round](-0.13,0.0)--(-0.04,-0.09)--(0.13,0.1);}
\newcommand{\popboxhead}[3]{\poppill{#1}{#2}{white}\ifx\relax#3\relax\else\hspace{6pt}{\popxbold\fontsize{10.6}{12}\selectfont\color{popink}#3}\fi\par\vspace{7pt}}

\newtcolorbox{aretenir}[1][]{popbase,colframe=popgreen,before upper={\popboxhead{À retenir}{popgreen}{#1}}}
\newtcolorbox{exemplebox}[1][]{popbase,colframe=poporange,before upper={\popboxhead{Exemple}{poporange}{#1}}}
\newtcolorbox{definition}[1][]{popbase,colframe=popblue,before upper={\popboxhead{Définition}{popblue}{#1}}}
\newtcolorbox{propriete}[1][]{popbase,colframe=poppurple,before upper={\popboxhead{Propriété}{poppurple}{#1}}}
\newtcolorbox{methode}[1][]{popbase,colframe=popgold,before upper={\popboxhead{Méthode}{popgold}{#1}}}
\newtcolorbox{attention}[1][]{popbase,colframe=poppink,before upper={\popboxhead{Attention}{poppink}{#1}}}
\newtcolorbox{experience}[1][]{popbase,colframe=popblue,colback=popblueL,before upper={\popboxhead{Expérience}{popblue}{#1}}}
% « Le savais-tu ? » : carte violette pleine (contexte, culture, Cameroun)
\newtcolorbox{savaistu}[1][]{popbase,colback=poppurple,colframe=popink,
  fontupper=\popsemi\fontsize{10.4}{14.2}\selectfont\color{white},
  before upper={\poppill{Le savais-tu\,?}{white}{poppurple}\ifx\relax#1\relax\else\hspace{6pt}{\popxbold\color{white}#1}\fi\par\vspace{7pt}}}
% Exercice résolu : énoncé en haut, \tcblower puis la solution sur fond crème
\newcounter{popexo}\newcounter{popexr}
\newtcolorbox{exoresolu}[1][]{popbase,colframe=poporange,colbacklower=popcream,
  segmentation style={popink,line width=1.2pt,dashed},
  before upper={\stepcounter{popexr}\popboxhead{Exercice résolu \thepopexr}{poporange}{#1}},
  before lower={\poppill{Solution}{popink}{popcream}\par\vspace{6pt}}}
% Exercice (non résolu en place) — la correction est dans la partie Corrigés
\newtcolorbox{exercice}[1][]{popbase,colframe=popink,
  before upper={\stepcounter{popexo}\popboxhead{Exercice \thepopexo}{popink}{#1}}}
% Corrigé
\newtcolorbox{corrige}[1][]{popbase,colframe=popgreen,colback=popgreenL,
  before upper={\popboxhead{Corrigé}{popgreen}{#1}}}
% Objectifs / prérequis / bilan
\newtcolorbox{objectifs}{popbase,colframe=popink,colback=poporangeL,
  before upper={\popboxhead{Objectifs de la leçon}{popink}{}}}
\newtcolorbox{prerequis}{popbase,colframe=popink,colback=popblueL,
  before upper={\popboxhead{Prérequis}{popblue}{}}}
\newtcolorbox{bilan}{popbase,colframe=popink,colback=white,boxrule=2.8pt,
  before upper={\popboxhead{Fiche bilan}{poporange}{L'essentiel à connaître pour l'examen}}}
% Figure encadrée avec légende
\newtcolorbox{popfigure}[1][]{enhanced,breakable=false,colback=white,colframe=popline,boxrule=1.4pt,arc=8pt,
  left=10pt,right=10pt,top=10pt,bottom=8pt,before skip=10pt,after skip=12pt,halign=center,
  lower separated=false,halign lower=center,
  fontlower=\popsemi\fontsize{9}{11.5}\selectfont\color{popmuted},
  #1}
\newcommand{\legende}[1]{\par\vspace{6pt}{\popsemi\fontsize{9}{11.5}\selectfont\color{popmuted}#1}}

% ============================================================
% Signature OnBuch+
% ============================================================
\newcommand{\poplogoinline}[1]{{\poplogo\fontsize{#1}{#1}\selectfont\color{popink}OnBuch\color{poporange}+}}
\newcommand{\signatureonbuch}{%
  \begin{tcolorbox}[enhanced,colback=popink,colframe=popink,boxrule=2.2pt,arc=10pt,
      drop shadow=poporange,left=14pt,right=14pt,top=10pt,bottom=10pt,before skip=0pt,after skip=0pt]
    \tikz[baseline=-0.7ex]\node[circle,fill=popgreen,draw=popcream,line width=1.3pt,minimum size=22pt,inner sep=0]{\popcheck{white}};%
    \hspace{9pt}\begin{minipage}[c]{0.62\linewidth}
      {\popxbold\fontsize{12}{13}\selectfont\color{popcream}Signé\ \textbullet\ {\poplogo OnBuch\color{poporange}+}}\par\vspace{2pt}
      {\raggedright\popsemi\fontsize{8}{10.5}\selectfont\color{popline}\MakeUppercase{Équipe pédagogique OnBuch+}\\\MakeUppercase{Conforme au programme officiel MINESEC}\par}
    \end{minipage}\hfill
    {\poplogo\fontsize{22}{22}\selectfont\color{popcream}OnBuch\color{poporange}+}
  \end{tcolorbox}%
}

% ============================================================
% Page de garde
% #1 titre · #2 matière · #3 niveau · #4 module · #5 n° leçon · #6 durée ·
% #7 description / objectifs (texte ou itemize)
% ============================================================
\newcommand{\pagegarde}[7]{%
  \thispagestyle{empty}
  \newgeometry{a4paper,margin=1.7cm,top=1.6cm,bottom=1.4cm}%
  \begin{tikzpicture}[remember picture,overlay]
    \fill[poporange!18] ([xshift=-3.2cm,yshift=-3.6cm]current page.north east) circle (4.6cm);
    \fill[poppurple!14] ([xshift=2.4cm,yshift=7.6cm]current page.south west) circle (3.8cm);
  \end{tikzpicture}%
  {\poplogo\fontsize{30}{30}\selectfont\color{popink}OnBuch\color{poporange}+}\hfill
  \raisebox{0.6ex}{\poppill{Cours signé \textbullet\ Premium}{popink}{popcream}}\par
  \vspace{1pt}\popsquiggle{poppurple}\par
  \vspace{8pt}
  \begin{tcolorbox}[enhanced,colback=poporange,colframe=popink,boxrule=2.8pt,arc=13pt,
      drop shadow=popink,left=18pt,right=18pt,top=12pt,bottom=13pt,after skip=10pt]
    \poppill{#2 \textbullet\ #3}{popink}{popcream}\hspace{5pt}\poppill{#4}{white}{popink}\par\vspace{8pt}
    {\popdisplay\fontsize{14}{14}\selectfont\color{popink}LEÇON #5}\par\vspace{4pt}
    {\raggedright\hyphenpenalty=10000\exhyphenpenalty=10000\popdisplay\fontsize{25}{27}\selectfont\color{popcream}\MakeUppercase{#1}\par}
    \vspace{8pt}
    {\popxbold\fontsize{9.5}{9.5}\selectfont\color{popink}\MakeUppercase{Durée conseillée : #6}}
  \end{tcolorbox}
  \begin{tcolorbox}[enhanced,colback=white,colframe=popink,boxrule=2.2pt,arc=10pt,
      drop shadow=popink,left=16pt,right=16pt,top=10pt,bottom=10pt,after skip=8pt]
    \poppill{Dans cette leçon}{poppurple}{white}\par\vspace{5pt}
    \popsemi\fontsize{9.3}{12.6}\selectfont\color{popink}#7
  \end{tcolorbox}
  \noindent\begin{minipage}[t]{0.487\linewidth}
    \begin{tcolorbox}[enhanced,colback=popgreen,colframe=popink,boxrule=2.2pt,arc=10pt,
        drop shadow=popink,left=11pt,right=11pt,top=10pt,bottom=10pt,height=3.1cm,valign=center]
      \tcbox[colback=white,colframe=popink,boxrule=1.3pt,arc=5pt,boxsep=0pt,nobeforeafter,
        left=3pt,right=3pt,top=3pt,bottom=3pt,valign=center]{\qrcode[height=1.9cm]{https://play.google.com/store/apps/details?id=com.luvvix.onbuch&hl=fr}}%
      \hspace{8pt}\begin{minipage}[c]{0.5\linewidth}
        {\popdisplay\fontsize{11}{12.5}\selectfont\color{white}\MakeUppercase{Télécharge OnBuch}}\par\vspace{4pt}
        {\raggedright\popsemi\fontsize{8.4}{11}\selectfont\color{white}Gratuit \textbullet\ Google Play\\Tuteur IA, annales, concours, résultats\par}
      \end{minipage}
    \end{tcolorbox}
  \end{minipage}\hfill
  \begin{minipage}[t]{0.487\linewidth}
    \begin{tcolorbox}[enhanced,colback=poppurple,colframe=popink,boxrule=2.2pt,arc=10pt,
        drop shadow=popink,left=11pt,right=11pt,top=10pt,bottom=10pt,height=3.1cm,valign=center]
      \tikz[baseline=-0.7ex]\node[circle,fill=white,draw=popink,line width=1.3pt,minimum size=1.6cm,inner sep=0]{\popdisplay\fontsize{24}{24}\selectfont\color{poppurple}+};%
      \hspace{8pt}\begin{minipage}[c]{0.6\linewidth}
        {\popdisplay\fontsize{11}{12.5}\selectfont\color{white}\MakeUppercase{Passe à OnBuch+}}\par\vspace{4pt}
        {\raggedright\popsemi\fontsize{8.4}{11}\selectfont\color{white}Tous les cours signés, Tuteur IA illimité, fiches PDF — onglet OnBuch+ de l'app\par}
      \end{minipage}
    \end{tcolorbox}
  \end{minipage}
  \vspace{8pt}
  \popcontenu
  \vfill
  \signatureonbuch
  \restoregeometry
  \newpage
}

% Rangée « Ce cours contient » (page de garde)
\newcommand{\poptile}[3]{% #1 couleur #2 gros texte #3 légende
  \begin{tcolorbox}[enhanced,colback=#1,colframe=popink,boxrule=2pt,arc=9pt,shadow={2.5pt}{-2.5pt}{0pt}{popink},
      left=6pt,right=6pt,top=9pt,bottom=9pt,halign=center,valign=center,height=1.75cm,width=\linewidth,nobeforeafter]
    {\popdisplay\fontsize{12.5}{13}\selectfont\color{white}\MakeUppercase{#2}}\par\vspace{3pt}
    {\centering\popsemi\fontsize{7.8}{9.5}\selectfont\color{white}#3\par}
  \end{tcolorbox}}
\newcommand{\popcontenu}{%
  \par\noindent{\popxboldls\fontsize{7.5}{7.5}\selectfont\color{popmuted}CE COURS SIGNÉ CONTIENT}\par\vspace{7pt}
  \noindent\begin{minipage}[t]{0.235\linewidth}\poptile{popblue}{Cours}{complet, illustré, pas à pas}\end{minipage}\hfill
  \begin{minipage}[t]{0.235\linewidth}\poptile{poporange}{Méthodes}{et exercices résolus}\end{minipage}\hfill
  \begin{minipage}[t]{0.235\linewidth}\poptile{poppink}{Exercices}{type examen + corrigés}\end{minipage}\hfill
  \begin{minipage}[t]{0.235\linewidth}\poptile{popgreen}{Fiche bilan}{l'essentiel à retenir}\end{minipage}\par}

% Sommaire
\newtcolorbox{sommaire}{enhanced,colback=white,colframe=popink,boxrule=2.2pt,arc=10pt,
  drop shadow=popink,left=18pt,right=18pt,top=13pt,bottom=13pt,before skip=6pt,after skip=20pt,
  before upper={\poppill{Sommaire}{poppurple}{white}\par\vspace{9pt}\popsemi\fontsize{10.6}{14.5}\selectfont\color{popink}}}

% Signature de fin de cours — poussée tout en bas de la dernière page
\newcommand{\findecours}{%
  \par\vfill\nopagebreak
  \begin{center}\popsquiggle{poporange}\end{center}\vspace{4pt}
  \signatureonbuch
}
\AtBeginDocument{\def\llbracket{\mathopen{[\mkern-3.5mu[}}\def\rrbracket{\mathclose{]\mkern-3.5mu]}}} % intervalles d'entiers (glyphe ⟦ absent de la police)
% Glyphes absents de Fira Math : redéfinis à partir de glyphes disponibles
\AtBeginDocument{%
  \def\setminus{\mathbin{\backslash}}\let\smallsetminus\setminus
  \def\vdots{\mathord{\vbox{\baselineskip3.2pt\lineskiplimit0pt\kern4pt\hbox{.}\hbox{.}\hbox{.}}}}%
  \def\ddots{\mathinner{\mkern1mu\raise6.5pt\hbox{.}\mkern2mu\raise3.6pt\hbox{.}\mkern2mu\raise0.7pt\hbox{.}\mkern1mu}}%
  \def\triangle{\mathord{\scalebox{0.9}{$\symup{\Delta}$}}}%
  \def\nabla{\mathord{\raisebox{\depth}{\scalebox{1}[-1]{$\symup{\Delta}$}}}}%
  \def\checkmark{\popcheck{popdark}}%
  \def\star{\mathbin{\ast}}\def\bigstar{\mathord{\ast}}%
  \def\diamond{\mathbin{\scalebox{0.6}{$\Diamond$}}}%
}

%% FIGFIX-BEGIN v3 — correctifs globaux des figures (docs/onbuch-generation/tools/figfix)
\usepackage{adjustbox}\usepackage{etoolbox}
% toute figure TikZ/pgfplots plus large que la ligne est réduite à la largeur disponible
\BeforeBeginEnvironment{tikzpicture}{\begin{adjustbox}{max width=\linewidth}}
\AfterEndEnvironment{tikzpicture}{\end{adjustbox}}
% légendes pgfplots lisibles par-dessus les courbes
\pgfplotsset{every axis/.append style={legend style={fill=white,fill opacity=0.92,text opacity=1,draw=black!15,inner sep=3pt}}}
% étiquettes d'arêtes (syntaxe edge["..."]) avec fond blanc
\tikzset{every edge quotes/.append style={fill=white,inner sep=1.2pt}}
% bulles de cartes mentales : texte à la ligne dans la bulle, bulle un peu plus grande
\tikzset{every concept/.append style={text width=2.0cm,align=center,minimum size=2.3cm,inner sep=2pt}}
%% FIGFIX-END
\begin{document}

\pagegarde{Étude de fonctions}{Mathématiques}{Tle A}{Relations et opérations fondamentales dans l'ensemble des nombres réels}{A08}{7 h}{Cette leçon outille l'élève pour mener l'étude complète des quatre grandes familles de fonctions au programme de Terminale A : polynômes de degré au plus 3, fonctions rationnelles homographiques généralisées, logarithmes et exponentielles. Elle culmine dans la résolution de problèmes concrets de modélisation économique et sociale.
\vspace{4pt}
\begin{multicols}{2}\small
\begin{itemize}
\item À la fin de cette leçon, je sais déterminer l'ensemble de définition d'une fonction polynôme, rationnelle, logarithmique ou exponentielle
\item À la fin de cette leçon, je sais calculer les limites aux bornes de l'ensemble de définition de ces fonctions
\item À la fin de cette leçon, je sais calculer la dérivée de chacun de ces types de fonctions et étudier son signe
\item À la fin de cette leçon, je sais dresser le tableau de variations complet d'une fonction
\item À la fin de cette leçon, je sais déterminer les asymptotes verticales, horizontales et obliques d'une courbe
\item À la fin de cette leçon, je sais tracer la courbe représentative d'une fonction étudiée dans un repère
\end{itemize}\end{multicols}}

\begin{sommaire}
\begin{multicols}{2}
\begin{enumerate}
\item Étude des fonctions polynômes de degré ≤ 3
\item Étude des fonctions rationnelles du type (ax² + bx + c)/(dx + e)
\item Étude des fonctions du type ln(ax + b)
\item Étude des fonctions du type e\^{}(ax + b)
\item Synthèse : méthode générale d'étude et comparaison des croissances
\item Problèmes de modélisation et optimisation
\item Activité d'intégration
\item Méthodes et exercices résolus
\item Exercices
\item Corrigés des exercices
\item Fiche bilan
\end{enumerate}
\end{multicols}
\end{sommaire}

\begin{objectifs}
\begin{itemize}
\item À la fin de cette leçon, je sais déterminer l'ensemble de définition d'une fonction polynôme, rationnelle, logarithmique ou exponentielle
\item À la fin de cette leçon, je sais calculer les limites aux bornes de l'ensemble de définition de ces fonctions
\item À la fin de cette leçon, je sais calculer la dérivée de chacun de ces types de fonctions et étudier son signe
\item À la fin de cette leçon, je sais dresser le tableau de variations complet d'une fonction
\item À la fin de cette leçon, je sais déterminer les asymptotes verticales, horizontales et obliques d'une courbe
\item À la fin de cette leçon, je sais tracer la courbe représentative d'une fonction étudiée dans un repère
\item À la fin de cette leçon, je sais résoudre graphiquement et algébriquement des équations et inéquations liées à ces fonctions
\item À la fin de cette leçon, je sais modéliser et résoudre un problème concret (coût, bénéfice, croissance) à l'aide de l'une de ces fonctions
\end{itemize}
\end{objectifs}

\begin{prerequis}
\begin{itemize}
\item Dérivation des fonctions usuelles et opérations sur les dérivées (Première)
\item Limites de fonctions et notion d'asymptote (début de Terminale)
\item Fonctions logarithme népérien et exponentielle : définitions et propriétés algébriques
\item Résolution d'équations et d'inéquations du premier et du second degré
\item Lecture graphique et tableau de variations (Première)
\end{itemize}
\end{prerequis}

\newpage

\coursec{Étude des fonctions polynômes de degré ≤ 3}

\textbf{Situation d'accroche.} M. Nkoulou gère un cybercafé dans le quartier Essos à Yaoundé. Chaque mois, il imprime des centaines de documents pour les élèves et les étudiants du voisinage. En relevant soigneusement ses comptes, il a remarqué trois choses étonnantes : le \emph{coût total} de ses impressions (encre, papier, électricité, usure des machines) évolue comme une fonction \cle{polynôme du troisième degré} du nombre d'impressions ; le \emph{prix moyen} par impression suit une fonction \emph{rationnelle} ; et le \emph{nombre de ses clients}, lui, croît selon une fonction \emph{exponentielle}. Trois questions le taraudent : combien d'impressions mensuelles minimisent son coût moyen ? Où se situe le seuil de rentabilité ? Comment prévoir l'évolution de sa clientèle pour investir au bon moment ?

\textbf{Problématique.} Toutes ces questions, et bien d'autres issues de la vie économique camerounaise, trouvent leurs réponses dans un même outil : l'\cle{étude complète de fonctions}. Dans cette leçon, nous apprenons à mener cette étude de bout en bout : ensemble de définition, limites, dérivée, variations, extremums, tracé de la courbe. Nous commençons par le cas le plus fondamental : les fonctions polynômes de degré inférieur ou égal à 3.

\courssub{Fonctions polynômes : définitions et premières propriétés}

\textbf{Intuition.} Un polynôme, c'est une fonction construite uniquement avec des additions, des soustractions et des multiplications de $x$ par lui-même et par des nombres. Ce sont les fonctions les plus « douces » qui soient : elles sont définies partout, continues partout, dérivables partout. C'est pourquoi on les rencontre partout, de la trajectoire d'un ballon au coût de production de M. Nkoulou.

\begin{definition}[Fonction polynôme de degré $n$]
On appelle \cle{fonction polynôme de degré $n$} (avec $n \in \mathbb{N}$) toute fonction $f$ définie sur $\mathbb{R}$ par
\[ f(x) = a_n x^n + a_{n-1} x^{n-1} + \cdots + a_1 x + a_0, \]
où $a_0, a_1, \ldots, a_n$ sont des réels appelés \cle{coefficients}, avec $a_n \neq 0$. Le réel $a_n$ est le \cle{coefficient dominant} et $a_n x^n$ le \cle{monôme de plus haut degré}.
\end{definition}

\begin{definition}[Fonction polynôme de degré 3, ou cubique]
On appelle \cle{cubique} (ou fonction polynôme du troisième degré) toute fonction $f$ définie sur $\mathbb{R}$ par
\[ f(x) = ax^3 + bx^2 + cx + d, \quad \text{avec } a, b, c, d \in \mathbb{R} \text{ et } a \neq 0. \]
Son ensemble de définition est $D_f = \mathbb{R}$ : aucune valeur interdite, aucun dénominateur, aucune racine carrée.
\end{definition}

\begin{aretenir}[Les trois familles de polynômes de degré $\leq 3$]
\begin{itemize}
\item Degré 1 : $f(x) = cx + d$ ($c \neq 0$) : \cle{fonction affine}, courbe = droite.
\item Degré 2 : $f(x) = bx^2 + cx + d$ ($b \neq 0$) : \cle{trinôme du second degré}, courbe = parabole.
\item Degré 3 : $f(x) = ax^3 + bx^2 + cx + d$ ($a \neq 0$) : \cle{cubique}, courbe en « S » plus ou moins marqué.
\end{itemize}
\end{aretenir}

\courssub{Limites aux bornes : le règne du monôme dominant}

\textbf{Intuition.} Quand $x$ devient immense (pense à $x = 10^{6}$), que pèse le terme $-3x^2$ face à $x^3$ ? Presque rien : $x^3 = 10^{18}$ alors que $3x^2 = 3 \times 10^{12}$, soit un million de fois moins. À l'infini, c'est le \cle{monôme de plus haut degré} qui commande, et tous les autres termes deviennent négligeables. Formalisons cette idée.

\begin{propriete}[Limite d'un polynôme en $\pm\infty$]
En $+\infty$ et en $-\infty$, une fonction polynôme a la même limite que son monôme de plus haut degré. En particulier, pour $f(x) = ax^3 + bx^2 + cx + d$ avec $a \neq 0$ :
\[ \lim_{x \to +\infty} f(x) = \lim_{x \to +\infty} ax^3 \quad \text{et} \quad \lim_{x \to -\infty} f(x) = \lim_{x \to -\infty} ax^3. \]
\end{propriete}

\begin{experience}[Démonstration guidée : factorisons le monôme dominant]
Montrons que $\displaystyle\lim_{x \to +\infty} f(x) = \lim_{x \to +\infty} ax^3$ pour $f(x) = ax^3 + bx^2 + cx + d$.
\begin{enumerate}
\item \textbf{Factorisons} $ax^3$, le terme dominant, valable pour $x \neq 0$ :
\[ f(x) = ax^3\left(1 + \frac{b}{a}\cdot\frac{1}{x} + \frac{c}{a}\cdot\frac{1}{x^2} + \frac{d}{a}\cdot\frac{1}{x^3}\right). \]
\item \textbf{Étudions chaque morceau} de la parenthèse. Quand $x \to +\infty$ :
\[ \frac{1}{x} \to 0, \qquad \frac{1}{x^2} \to 0, \qquad \frac{1}{x^3} \to 0. \]
\item \textbf{Concluons sur la parenthèse} : par somme de limites, elle tend vers $1 + 0 + 0 + 0 = 1$.
\item \textbf{Concluons sur $f$} : par produit de limites, $f(x)$ se comporte comme $ax^3 \times 1$, c'est-à-dire comme $ax^3$. Le même raisonnement vaut en $-\infty$. $\blacksquare$
\end{enumerate}
\end{experience}

\begin{exemplebox}[Calcul direct de limites]
Soit $f(x) = -2x^3 + 5x^2 - x + 7$. Le monôme dominant est $-2x^3$, de coefficient $-2 < 0$ et de degré impair. Donc :
\[ \lim_{x \to +\infty} f(x) = \lim_{x \to +\infty} (-2x^3) = -\infty \quad \text{et} \quad \lim_{x \to -\infty} f(x) = \lim_{x \to -\infty} (-2x^3) = +\infty. \]
Retiens : pour une cubique, les deux limites aux bornes sont \emph{toujours} infinies et de signes \emph{opposés}, car $x^3$ change de signe avec $x$.
\end{exemplebox}

\courssub{Dérivée et sens de variation : les trois visages d'une cubique}

\begin{propriete}[Dérivée d'une cubique]
Si $f(x) = ax^3 + bx^2 + cx + d$, alors $f$ est dérivable sur $\mathbb{R}$ et
\[ f'(x) = 3ax^2 + 2bx + c. \]
La dérivée d'une cubique est donc un \cle{trinôme du second degré}.
\end{propriete}

\textbf{Intuition.} Voici la clé de toute la section : le sens de variation de $f$ est donné par le signe de $f'$, et $f'$ est un trinôme. Or tu sais depuis la Première étudier le signe d'un trinôme grâce à son discriminant. Tout se ramène donc à l'étude de
\[ \Delta = (2b)^2 - 4(3a)(c) = 4b^2 - 12ac = 4(b^2 - 3ac). \]
Comme le facteur $4$ est positif, le signe de $\Delta$ est celui de $b^2 - 3ac$. Trois cas se présentent.

\textbf{Cas 1 : $\Delta < 0$ (c'est-à-dire $b^2 - 3ac < 0$).} Le trinôme $f'$ ne s'annule jamais et garde un signe constant : celui de $3a$, donc celui de $a$. La fonction $f$ est \emph{strictement monotone} sur $\mathbb{R}$ (strictement croissante si $a > 0$, strictement décroissante si $a < 0$). Aucun extremum.

\textbf{Cas 2 : $\Delta = 0$.} Le trinôme $f'$ s'annule en une unique racine double $x_0 = -\dfrac{b}{3a}$ mais \emph{sans changer de signe}. La fonction $f$ reste strictement monotone sur $\mathbb{R}$ ; la courbe présente simplement une tangente horizontale au point d'abscisse $x_0$.

\textbf{Cas 3 : $\Delta > 0$.} Le trinôme $f'$ s'annule en deux racines distinctes
\[ x_1 = \frac{-b - \sqrt{b^2 - 3ac}}{3a} \quad \text{et} \quad x_2 = \frac{-b + \sqrt{b^2 - 3ac}}{3a}, \]
et change de signe en chacune d'elles. La fonction $f$ admet alors \cle{deux extremums locaux} : un \cle{maximum local} et un \cle{minimum local}.

\begin{attention}[Piège classique : « une cubique a toujours un maximum et un minimum »]
C'est \textbf{faux} ! Une cubique ne possède deux extremums locaux que si le discriminant de sa dérivée est \emph{strictement positif}. Si $\Delta \leq 0$, la fonction est strictement monotone sur $\mathbb{R}$ et n'a \emph{aucun} extremum. Exemple : $f(x) = x^3 + x$ vérifie $f'(x) = 3x^2 + 1 > 0$ pour tout $x$ : elle est strictement croissante, sans maximum ni minimum. Au Bac, calcule toujours $\Delta$ avant d'annoncer des extremums.
\end{attention}

\begin{popfigure}
\begin{center}
\begin{tabular}{|c|ccccccc|}
\hline
\rowcolor{popblueL} $x$ & $-\infty$ & & $x_1$ & & $x_2$ & & $+\infty$ \\
\hline
$f'(x)$ & & $+$ & $0$ & $-$ & $0$ & $+$ & \\
\hline
$f(x)$ & $-\infty$ & $\nearrow$ & $f(x_1)$ & $\searrow$ & $f(x_2)$ & $\nearrow$ & $+\infty$ \\
\hline
\end{tabular}

\vspace{0.4cm}
\textbf{Cas $a > 0$, $\Delta > 0$} : maximum local en $x_1$, minimum local en $x_2$.

\vspace{0.6cm}
\begin{tabular}{|c|ccccccc|}
\hline
\rowcolor{poporangeL} $x$ & $-\infty$ & & $x_1$ & & $x_2$ & & $+\infty$ \\
\hline
$f'(x)$ & & $-$ & $0$ & $+$ & $0$ & $-$ & \\
\hline
$f(x)$ & $+\infty$ & $\searrow$ & $f(x_1)$ & $\nearrow$ & $f(x_2)$ & $\searrow$ & $-\infty$ \\
\hline
\end{tabular}

\vspace{0.4cm}
\textbf{Cas $a < 0$, $\Delta > 0$} : minimum local en $x_1$, maximum local en $x_2$.

\vspace{0.6cm}
\begin{tabular}{|c|ccc|}
\hline
\rowcolor{popgreenL} $x$ & $-\infty$ & & $+\infty$ \\
\hline
$f'(x)$ & & $+$ & \\
\hline
$f(x)$ & $-\infty$ & $\nearrow$ & $+\infty$ \\
\hline
\end{tabular}

\vspace{0.4cm}
\textbf{Cas $a > 0$, $\Delta \leq 0$} : $f$ strictement croissante, aucun extremum (si $a<0$, tout s'inverse).
\end{center}
\legende{Les trois tableaux de variations possibles d'une cubique selon le signe de $\Delta$ et de $a$.}
\end{popfigure}

\begin{definition}[Extremum local]
Soit $f$ une fonction définie sur un intervalle $I$ et $x_0 \in I$. On dit que $f$ admet un \cle{maximum local} en $x_0$ s'il existe un intervalle ouvert $J$ contenant $x_0$ tel que $f(x) \leq f(x_0)$ pour tout $x \in J$. On définit de même un \cle{minimum local}. En pratique, pour une fonction dérivable, les extremums locaux se trouvent parmi les points où $f'$ \emph{s'annule en changeant de signe}.
\end{definition}

\courssub{Exemple chiffré : étude complète de $f(x) = x^3 - 3x^2 + 2$}

\begin{exoresolu}[Étude complète d'une cubique]
On considère la fonction $f$ définie sur $\mathbb{R}$ par $f(x) = x^3 - 3x^2 + 2$. Mener l'étude complète de $f$ : limites, dérivée, variations, extremums, points remarquables.
\tcblower
\textbf{1. Ensemble de définition.} $f$ est un polynôme : $D_f = \mathbb{R}$.

\textbf{2. Limites aux bornes.} Le monôme dominant est $x^3$ :
\[ \lim_{x \to +\infty} f(x) = \lim_{x \to +\infty} x^3 = +\infty \quad ; \quad \lim_{x \to -\infty} f(x) = \lim_{x \to -\infty} x^3 = -\infty. \]

\textbf{3. Dérivée.} $f'(x) = 3x^2 - 6x = 3x(x - 2)$.

\textbf{4. Signe de $f'$.} $f'$ s'annule en $x = 0$ et $x = 2$. C'est un trinôme avec $a' = 3 > 0$ : il est positif à l'extérieur des racines, négatif entre elles. Donc $f'(x) > 0$ sur $]-\infty ; 0[ \cup ]2 ; +\infty[$ et $f'(x) < 0$ sur $]0 ; 2[$.

\textbf{5. Tableau de variations.}
\begin{center}
\begin{tabular}{|c|ccccccc|}
\hline
\rowcolor{popblueL} $x$ & $-\infty$ & & $0$ & & $2$ & & $+\infty$ \\
\hline
$f'(x)$ & & $+$ & $0$ & $-$ & $0$ & $+$ & \\
\hline
$f(x)$ & $-\infty$ & $\nearrow$ & $2$ & $\searrow$ & $-2$ & $\nearrow$ & $+\infty$ \\
\hline
\end{tabular}
\end{center}

\textbf{6. Extremums locaux.} $f(0) = 2$ : \cle{maximum local} au point $A(0\,;\,2)$. $f(2) = 8 - 12 + 2 = -2$ : \cle{minimum local} au point $B(2\,;\,-2)$.

\textbf{7. Points remarquables.} Intersection avec l'axe des ordonnées : $f(0) = 2$, point $(0\,;\,2)$. Intersections avec l'axe des abscisses : on résout $x^3 - 3x^2 + 2 = 0$ ; $x = 1$ est racine évidente ($1 - 3 + 2 = 0$), et la factorisation donne $f(x) = (x-1)(x^2 - 2x - 2)$, d'où $x = 1 - \sqrt{3} \approx -0{,}73$ et $x = 1 + \sqrt{3} \approx 2{,}73$.

\textbf{8. Tangente en un point.} Par exemple en $x = 1$ : $f(1) = 0$, $f'(1) = -3$, d'où la tangente $y = -3(x - 1)$, soit $y = -3x + 3$.
\end{exoresolu}

\courssub{Le point d'inflexion : le « centre de symétrie » de la cubique}

\textbf{Intuition.} Regarde la courbe d'une cubique : elle « tourne » d'abord dans un sens, puis dans l'autre. Le point où elle change de sens de courbure s'appelle le \cle{point d'inflexion}. Fait remarquable : c'est aussi le \emph{centre de symétrie} de la courbe.

\begin{propriete}[Point d'inflexion d'une cubique — complément utile au Bac]
La courbe de $f(x) = ax^3 + bx^2 + cx + d$ admet un unique point d'inflexion en
\[ x_I = -\frac{b}{3a}, \]
abscisse où la dérivée seconde $f''(x) = 6ax + 2b$ s'annule en changeant de signe. Ce point $I\left(-\dfrac{b}{3a}\,;\, f\left(-\dfrac{b}{3a}\right)\right)$ est le \cle{centre de symétrie} de la courbe. Ce résultat n'est pas exigible, mais il est précieux pour vérifier un tracé ou accélérer certaines questions au Bac.
\end{propriete}

\begin{exemplebox}[Vérification sur notre exemple]
Pour $f(x) = x^3 - 3x^2 + 2$ : $x_I = -\dfrac{-3}{3 \times 1} = 1$ et $f(1) = 0$. Le point d'inflexion est $I(1\,;\,0)$. Vérifions la symétrie : le maximum $A(0\,;\,2)$ et le minimum $B(2\,;\,-2)$ ont bien pour milieu $\left(\dfrac{0+2}{2}\,;\,\dfrac{2+(-2)}{2}\right) = (1\,;\,0) = I$. Les deux extremums sont symétriques par rapport au point d'inflexion : c'est toujours vrai pour une cubique !
\end{exemplebox}

\begin{popfigure}
\begin{center}
\begin{tikzpicture}
\begin{axis}[
  width=0.85\linewidth, height=7cm,
  axis lines=middle, xlabel=$x$, ylabel=$y$,
  xmin=-2.6, xmax=3.4, ymin=-3.5, ymax=4.5,
  xtick={-2,-1,0,1,2,3}, ytick={-3,-2,-1,1,2,3,4},
  grid=both, grid style={popline!40},
  legend style={at={(0.02,0.98)}, anchor=north west, font=\small}
]
\addplot[popcurve,domain=-2.15:3.1,samples=200]{x^3 - 3*x + 1};
\addlegendentry{$f(x)=x^3-3x+1$}
\addplot[poporange, thick, dashed, domain=-1:1, samples=2] {-3*x + 1};
\addlegendentry{tangente en $I$ : $y=-3x+1$}
\addplot[only marks, mark=*, mark size=2pt, poppink] coordinates {(-1,3) (1,-1) (0,1)};
\node[popdark, anchor=south west, font=\small] at (axis cs:-1,3) {max local $(-1\,;\,3)$};
\node[popdark, anchor=north west, font=\small] at (axis cs:1,-1) {min local $(1\,;\,-1)$};
\node[popdark, anchor=west, font=\small] at (axis cs:0.15,1.3) {inflexion $I(0\,;\,1)$};
\end{axis}
\end{tikzpicture}
\end{center}
\legende{Courbe de $f(x) = x^3 - 3x + 1$. Les extremums locaux sont en $(-1\,;\,3)$ (maximum) et $(1\,;\,-1)$ (minimum). Le point d'inflexion $I(0\,;\,1)$ est le centre de symétrie. La tangente en $I$ d'équation $y = -3x + 1$ (en pointillés orange) traverse la courbe.}
\end{popfigure}

\courssub{Cas particuliers : degré 2 et degré 1}

\textbf{Fonctions du second degré.} Pour $f(x) = bx^2 + cx + d$ ($b \neq 0$), tout se simplifie : $f'(x) = 2bx + c$ est affine et s'annule une seule fois, en $x_S = -\dfrac{c}{2b}$. La courbe est une \cle{parabole} de sommet $S\left(-\dfrac{c}{2b}\,;\, f\left(-\dfrac{c}{2b}\right)\right)$, tournée vers le haut si $b > 0$ (le sommet est un minimum), vers le bas si $b < 0$ (maximum). La droite d'équation $x = -\dfrac{c}{2b}$ est l'\cle{axe de symétrie} de la parabole. Les limites en $\pm\infty$ sont celles de $bx^2$ : toutes deux égales à $+\infty$ si $b > 0$, à $-\infty$ si $b < 0$.

\textbf{Fonctions affines.} Pour $f(x) = cx + d$ ($c \neq 0$), la courbe est une droite de coefficient directeur $c$ : strictement croissante si $c > 0$, strictement décroissante si $c < 0$. Les limites sont $\pm\infty$ selon le signe de $c$.

\begin{methode}[Plan type d'étude complète d'une fonction]
Voici la feuille de route que tu suivras pour \emph{toute} fonction de cette leçon :
\begin{enumerate}
\item \textbf{Ensemble de définition} $D_f$ : cherche les valeurs interdites (dénominateur nul, quantité sous un logarithme $\leq 0$, etc.). Pour un polynôme : $D_f = \mathbb{R}$.
\item \textbf{Limites aux bornes} de $D_f$ : pour un polynôme, utilise le monôme dominant.
\item \textbf{Dérivée} $f'$ : calcule-la soigneusement et factorise-la si possible.
\item \textbf{Signe de $f'$} : résous $f'(x) = 0$ et dresse le tableau de signes (pour une cubique, utilise $\Delta$).
\item \textbf{Tableau de variations} : reporte les limites, les zéros de $f'$ et les valeurs des extremums.
\item \textbf{Points remarquables} : intersections avec les axes, tangentes utiles, point d'inflexion éventuel.
\item \textbf{Tracé de la courbe} : place les extremums, les points remarquables, respecte le tableau de variations.
\end{enumerate}
\end{methode}

\begin{attention}[Erreurs fréquentes à bannir]
\begin{itemize}
\item \textbf{Dérivation fautive} : la dérivée de $bx^2$ est $2bx$, pas $bx$ ; celle de la constante $d$ est $0$. Écris $f'(x) = 3ax^2 + 2bx + c$ sans raccourci.
\item \textbf{Confusion extremum / valeur} : le maximum local est la \emph{valeur} $f(x_1)$, atteinte \emph{en} $x_1$. Au Bac, précise toujours les deux : « maximum local égal à $2$, atteint en $x = 0$ ».
\item \textbf{Tableau incohérent} : les valeurs placées dans la ligne de $f$ doivent être compatibles avec les flèches (une fonction croissante ne peut pas passer de $5$ à $2$).
\item \textbf{Oubli du cas $\Delta \leq 0$} : ne dessine jamais deux extremums « par habitude ».
\end{itemize}
\end{attention}

\begin{exercice}[À toi de jouer]
Soit $g$ la fonction définie sur $\mathbb{R}$ par $g(x) = -x^3 + 3x + 2$.
\begin{enumerate}
\item Déterminer les limites de $g$ en $-\infty$ et en $+\infty$.
\item Calculer $g'(x)$ et étudier son signe.
\item Dresser le tableau de variations de $g$ et préciser ses extremums locaux.
\item Déterminer l'équation de la tangente à la courbe au point d'abscisse $0$.
\end{enumerate}
\end{exercice}

\begin{corrige}[Exercice 1]
\textbf{1.} Monôme dominant $-x^3$ : $\displaystyle\lim_{x \to +\infty} g(x) = -\infty$ et $\displaystyle\lim_{x \to -\infty} g(x) = +\infty$.

\textbf{2.} $g'(x) = -3x^2 + 3 = -3(x^2 - 1) = -3(x-1)(x+1)$. Trinôme de racines $-1$ et $1$, de coefficient dominant $-3 < 0$ : $g'(x) < 0$ sur $]-\infty ; -1[ \cup ]1 ; +\infty[$ et $g'(x) > 0$ sur $]-1 ; 1[$.

\textbf{3.} $g$ est décroissante de $+\infty$ à $g(-1) = 1 - 3 + 2 = 0$ (minimum local $0$ en $x = -1$), puis croissante jusqu'à $g(1) = -1 + 3 + 2 = 4$ (maximum local $4$ en $x = 1$), puis décroissante vers $-\infty$.

\textbf{4.} $g(0) = 2$ et $g'(0) = 3$ : la tangente a pour équation $y = g'(0)(x - 0) + g(0)$, soit $y = 3x + 2$.
\end{corrige}

\begin{savaistu}[Les cubiques au service de l'économie]
En sciences économiques, le \emph{coût total} de production d'une entreprise est très souvent modélisé par une fonction polynôme du troisième degré $C(q) = aq^3 + bq^2 + cq + d$ : le terme constant $d$ représente les coûts fixes (loyer, machines), et le terme cubique traduit la hausse brutale des coûts quand on sature les capacités de production. Le \emph{coût marginal} — coût d'une unité supplémentaire — est alors exactement la dérivée $C'(q)$ ! C'est précisément l'outil qu'utiliserait M. Nkoulou pour piloter son cybercafé : en étudiant les variations de ses fonctions de coût, il repère les niveaux de production les plus rentables. Les mathématiques de Terminale sont déjà des mathématiques de gestionnaire.
\end{savaistu}

\begin{aretenir}[L'essentiel de la section]
\begin{itemize}
\item Un polynôme est défini, continu et dérivable sur $\mathbb{R}$ tout entier.
\item En $\pm\infty$, un polynôme a la même limite que son \cle{monôme de plus haut degré} (on le prouve en factorisant ce monôme).
\item La dérivée de $f(x) = ax^3 + bx^2 + cx + d$ est $f'(x) = 3ax^2 + 2bx + c$ : son signe se discute grâce à $\Delta = 4(b^2 - 3ac)$.
\item Si $\Delta > 0$ : deux extremums locaux ; si $\Delta \leq 0$ : fonction strictement monotone, aucun extremum.
\item Le point d'abscisse $-\dfrac{b}{3a}$ est le point d'inflexion, centre de symétrie de la courbe.
\item L'étude complète suit toujours le même plan : $D_f$, limites, dérivée, signe de $f'$, tableau, points remarquables, tracé.
\end{itemize}
\end{aretenir}

\coursec{Étude des fonctions rationnelles du type $\dfrac{ax^2+bx+c}{dx+e}$}

Dans la section précédente, nous avons étudié les fonctions polynômes : des fonctions \emph{définies partout}, aux courbes lisses et sans rupture. La vie réelle nous propose souvent des situations plus subtiles. Imagine un transporteur de Yaoundé qui doit répartir un coût fixe entre ses passagers : le coût par passager est un \emph{quotient}, et lorsque le nombre de passagers approche une certaine valeur critique, le coût peut « exploser ». Mathématiquement, cela se traduit par un \emph{dénominateur qui s'annule} : la courbe de la fonction présente alors une brisure, une asymptote verticale, que nous allons apprendre à maîtriser. Bienvenue dans l'étude des fonctions rationnelles, grand classique du Baccalauréat A.

\courssub{Fonction rationnelle et ensemble de définition}

\begin{definition}[Fonction rationnelle du type étudié]
On appelle \cle{fonction rationnelle} du type $\dfrac{ax^2+bx+c}{dx+e}$ toute fonction $f$ définie par
\[ f(x)=\dfrac{ax^2+bx+c}{dx+e} \quad \text{avec } a\neq 0 \text{ et } d\neq 0. \]
Le numérateur est un trinôme du second degré, le dénominateur est une fonction affine. La valeur $x_0=-\dfrac{e}{d}$, qui annule le dénominateur, est appelée \cle{valeur interdite}.
\end{definition}

L'ensemble de définition s'obtient en excluant la valeur interdite :
\[ D_f=\mathbb{R}\setminus\left\{-\dfrac{e}{d}\right\}=\left]-\infty;-\dfrac{e}{d}\right[\cup\left]-\dfrac{e}{d};+\infty\right[. \]

\begin{exemplebox}[Lecture directe]
Pour $f(x)=\dfrac{x^2+2x+2}{x+1}$, le dénominateur s'annule en $x=-1$. Donc $D_f=\mathbb{R}\setminus\{-1\}$.
\end{exemplebox}

\begin{attention}[La valeur interdite, partout !]
La valeur interdite $x_0=-\dfrac{e}{d}$ doit apparaître \textbf{trois fois} dans ta copie : dans l'ensemble de définition, dans le calcul des limites (à gauche et à droite de $x_0$), et dans le tableau de variations (double barre verticale). Oublier la double barre au tableau est l'erreur la plus fréquente — et la plus sanctionnée — au Baccalauréat.
\end{attention}

\begin{attention}[Cas particulier : simplification possible]
Si le numérateur s'annule \emph{aussi} en $x_0$ (i.e. $ax_0^2+bx_0+c=0$), la fraction se simplifie par $x-x_0$. La fonction se prolonge par continuité en $x_0$ (trou dans le graphique) et \textbf{n'a pas d'asymptote verticale}. Vérifie toujours si le numérateur s'annule en $x_0$ avant de conclure à l'asymptote verticale.
\end{attention}

\courssub{Limites à la valeur interdite et asymptote verticale}

Que se passe-t-il lorsque $x$ se rapproche de $x_0=-\dfrac{e}{d}$ ? Le numérateur $ax^2+bx+c$ tend vers un nombre $N=ax_0^2+bx_0+c$ (supposé non nul ici), tandis que le dénominateur tend vers $0$. Le quotient « explose » : la limite est infinie, et son signe dépend du signe du dénominateur, qui change selon que l'on approche $x_0$ par la gauche ou par la droite.

\begin{propriete}[Asymptote verticale (admise)]
Soit $f$ une fonction et $x_0$ un réel. Si
\[ \lim_{\substack{x\to x_0\\ x<x_0}} f(x)=\pm\infty \quad \text{ou} \quad \lim_{\substack{x\to x_0\\ x>x_0}} f(x)=\pm\infty, \]
alors la droite d'équation $x=x_0$ est \cle{asymptote verticale} à la courbe $\mathcal{C}_f$.
\end{propriete}

\begin{methode}[Déterminer le signe de la limite infinie]
\begin{enumerate}
\item Calculer la limite du numérateur en $x_0$ : on obtient un réel $N$ (supposé non nul) dont on détermine le signe.
\item Étudier le signe du dénominateur $dx+e$ : il est du signe de $d$ pour $x>x_0$ et du signe opposé pour $x<x_0$.
\item Appliquer la règle des signes au quotient $N/(dx+e)$ avec un dénominateur tendant vers $0$ : la limite est $+\infty$ ou $-\infty$.
\end{enumerate}
\end{methode}

\begin{exemplebox}[Limites en $-1$ pour $f(x)=\dfrac{x^2+2x+2}{x+1}$]
Le numérateur en $-1$ vaut $(-1)^2+2(-1)+2=1>0$. Le dénominateur $x+1$ est négatif si $x<-1$ et positif si $x>-1$. Donc :
\[ \lim_{\substack{x\to -1\\ x<-1}} f(x)=-\infty \qquad \text{et} \qquad \lim_{\substack{x\to -1\\ x>-1}} f(x)=+\infty. \]
La droite d'équation $x=-1$ est asymptote verticale à $\mathcal{C}_f$.
\end{exemplebox}

\courssub{Division euclidienne et asymptote oblique}

Aux voisinages de $+\infty$ et $-\infty$, le comportement de $f$ est dicté par les termes de plus haut degré :
\[ \lim_{x\to\pm\infty} f(x)=\lim_{x\to\pm\infty}\dfrac{ax^2}{dx}=\lim_{x\to\pm\infty}\dfrac{a}{d}\,x=\pm\infty \quad (\text{car } a/d \neq 0). \]
La limite est infinie : pas d'asymptote horizontale. Mais la courbe ne part pas n'importe comment vers l'infini : elle se rapproche d'une \emph{droite oblique}. Pour la découvrir, on effectue la division euclidienne du numérateur par le dénominateur.

\begin{propriete}[Écriture réduite et asymptote oblique]
Toute fonction $f(x)=\dfrac{ax^2+bx+c}{dx+e}$ ($a\neq0$, $d\neq0$) s'écrit de manière unique sous la forme
\[ f(x)=\alpha x+\beta+\dfrac{\gamma}{dx+e} \qquad \text{avec } \alpha=\dfrac{a}{d}\neq 0. \]
La droite $\Delta$ d'équation $y=\alpha x+\beta$ est alors \cle{asymptote oblique} à $\mathcal{C}_f$ en $+\infty$ et en $-\infty$, car
\[ \lim_{x\to\pm\infty}\bigl[f(x)-(\alpha x+\beta)\bigr]=\lim_{x\to\pm\infty}\dfrac{\gamma}{dx+e}=0. \]
\end{propriete}

\begin{experience}[Démonstration guidée : pourquoi cette écriture existe-t-elle ?]
La division euclidienne des polynômes affirme que pour tout polynôme $P$ et tout polynôme non nul $D$, il existe un unique couple $(Q,R)$ tel que $P=D\times Q+R$ avec $\deg R<\deg D$. Ici $P(x)=ax^2+bx+c$ et $D(x)=dx+e$ (degré 1), donc le reste $R$ est une constante $\gamma$ et le quotient $Q$ est de degré 1 : $Q(x)=\alpha x+\beta$. En divisant l'égalité $P=D\cdot Q+\gamma$ par $D$, on obtient exactement $f(x)=\alpha x+\beta+\dfrac{\gamma}{dx+e}$. La différence $f(x)-(\alpha x+\beta)=\dfrac{\gamma}{dx+e}$ tend vers $0$ quand $x\to\pm\infty$ : la courbe « colle » à la droite $\Delta$, ce qui justifie le mot \emph{asymptote}.
\end{experience}

\begin{methode}[Effectuer la division en 3 étapes]
Pour écrire $f(x)=\dfrac{x^2-x+1}{x-1}$ sous la forme $\alpha x+\beta+\dfrac{\gamma}{x-1}$ :
\begin{enumerate}
\item \textbf{Premier terme du quotient.} On divise le terme dominant du numérateur par celui du dénominateur : $\dfrac{x^2}{x}=x$. On calcule $x(x-1)=x^2-x$ et on soustrait : $(x^2-x+1)-(x^2-x)=1$.
\item \textbf{Reste.} Le reste $1$ est de degré $0<1$ : la division est terminée. Le quotient est $x$, le reste est $1$.
\item \textbf{Conclusion.} $x^2-x+1=(x-1)\cdot x+1$, donc $f(x)=x+\dfrac{1}{x-1}$. L'asymptote oblique est la droite d'équation $y=x$.
\end{enumerate}
\emph{Astuce de vérification :} réduis au même dénominateur $\alpha x+\beta+\dfrac{\gamma}{dx+e}$ et vérifie que tu retrouves le numérateur de départ.
\end{methode}

\begin{propriete}[Position relative de la courbe et de l'asymptote oblique]
La position de $\mathcal{C}_f$ par rapport à $\Delta : y=\alpha x+\beta$ est donnée par le signe de la différence
\[ f(x)-(\alpha x+\beta)=\dfrac{\gamma}{dx+e}. \]
\begin{itemize}
\item Si $\dfrac{\gamma}{dx+e}>0$, la courbe est \textbf{au-dessus} de $\Delta$ ;
\item si $\dfrac{\gamma}{dx+e}<0$, la courbe est \textbf{en dessous} de $\Delta$.
\end{itemize}
Le signe dépend donc de $\gamma$ et de la position de $x$ par rapport à la valeur interdite $-\dfrac{e}{d}$.
\end{propriete}

\begin{propriete}[Intersection avec l'asymptote oblique]
La courbe $\mathcal{C}_f$ coupe son asymptote oblique $\Delta$ si et seulement si l'équation $f(x)=\alpha x+\beta$ admet une solution, c'est-à-dire si $\gamma=0$. Si $\gamma\neq0$, la courbe ne rencontre jamais $\Delta$.
\end{propriete}

\begin{exemplebox}[Position et intersection pour $f(x)=\dfrac{x^2-x+1}{x-1}$]
On a $f(x)-x=\dfrac{1}{x-1}$. Donc :
\begin{itemize}
\item si $x>1$, $f(x)-x>0$ : la courbe est au-dessus de la droite $y=x$ ;
\item si $x<1$, $f(x)-x<0$ : la courbe est en dessous de la droite $y=x$.
\end{itemize}
Comme $\gamma=1\neq0$, la courbe ne traverse jamais son asymptote oblique.
\end{exemplebox}

\begin{popfigure}
\begin{center}
\begin{tikzpicture}
\begin{axis}[
  width=0.92\linewidth, height=8cm,
  axis lines=middle, axis line style={->},
  xlabel={$x$}, ylabel={$y$},
  xmin=-3.5, xmax=5.5, ymin=-5, ymax=7,
  xtick={-3,-2,-1,0,1,2,3,4,5}, ytick={-4,-2,2,4,6},
  tick label style={font=\small},
  clip=true]
  \addplot[popcurve,domain=-3.5:0.82,samples=150]{(x^2-x+1)/(x-1)};
  \addplot[popcurve,domain=1.18:5.5,samples=150]{(x^2-x+1)/(x-1)};
  \addplot[poporange,dashed,domain=-3.5:5.5,samples=2]{x};
  \draw[popgreen,dashed] (axis cs:1,-5) -- (axis cs:1,7);
  \node[poporange,font=\small,rotate=38] at (axis cs:4.1,3.4) {$y=x$};
  \node[popgreen,font=\small] at (axis cs:1.25,6.3) {$x=1$};
  \node[popblue,font=\small] at (axis cs:4.3,5.6) {$\mathcal{C}_f$};
\end{axis}
\end{tikzpicture}
\end{center}
\legende{Courbe de $f(x)=\dfrac{x^2-x+1}{x-1}$, avec son asymptote verticale $x=1$ (vert) et son asymptote oblique $y=x$ (orange).}
\end{popfigure}

\begin{popfigure}
\begin{center}
\begin{tikzpicture}[scale=1.0]
  \draw[->,popdark] (-4.5,0) -- (4.5,0) node[below left]{$x$};
  \draw[->,popdark] (0,-2.6) -- (0,3.2) node[below left]{$y$};
  \draw[poporange,dashed,thick] (-4.2,-1.9) -- (4.2,2.7) node[pos=0.92,below,sloped,font=\small]{$\Delta : y=\alpha x+\beta$};
  \draw[popgreen,dashed,thick] (0.9,-2.6) -- (0.9,3.2) node[above,font=\small]{$x=x_0$};
  \draw[popblue,very thick,domain=-4.2:0.35,samples=80] plot (\x,{0.55*\x+0.35-1.1/(\x-0.9)});
  \draw[popblue,very thick,domain=1.45:4.2,samples=80] plot (\x,{0.55*\x+0.35-1.1/(\x-0.9)});
  \node[popblue,font=\small] at (-3.1,-2.1) {courbe en dessous};
  \node[popblue,font=\small] at (3.4,3.0) {courbe au-dessus};
\end{tikzpicture}
\end{center}
\legende{Position relative : ici $\gamma<0$, donc $\mathcal{C}_f$ est en dessous de $\Delta$ pour $x<x_0$ et au-dessus pour $x>x_0$.}
\end{popfigure}

\courssub{Dérivée et sens de variation}

\begin{propriete}[Dérivée d'un quotient]
Si $f=\dfrac{u}{v}$ avec $v$ ne s'annulant pas sur un intervalle $I$, alors $f$ est dérivable sur $I$ et
\[ f'(x)=\dfrac{u'(x)\,v(x)-u(x)\,v'(x)}{\bigl[v(x)\bigr]^2}. \]
Pour $f(x)=\dfrac{ax^2+bx+c}{dx+e}$, le numérateur de $f'$ est un \cle{trinôme du second degré} : le signe de $f'$ s'étudie donc avec le discriminant, comme en Première.
\end{propriete}

\begin{attention}[Le carré au dénominateur]
Le dénominateur de $f'$ est $(dx+e)^2$, toujours strictement positif sur $D_f$ : il ne faut \textbf{jamais} l'inclure dans l'étude du signe. En revanche, la valeur interdite $-\dfrac{e}{d}$ doit figurer dans le tableau de variations avec une double barre, même si $f'$ n'y est pas définie.
\end{attention}

\courssub{Méthode générale d'étude complète}

\begin{methode}[Étude complète d'une fonction rationnelle $f(x)=\dfrac{ax^2+bx+c}{dx+e}$]
\begin{enumerate}
\item \textbf{Ensemble de définition :} $D_f=\mathbb{R}\setminus\{-e/d\}$. Vérifier si le numérateur s'annule aussi en $-e/d$ (simplification possible).
\item \textbf{Limites et asymptotes :}
      \begin{itemize}
      \item Limites en $\pm\infty$ (termes de plus haut degré).
      \item Limites à la valeur interdite (gauche/droite) $\to$ Asymptote verticale $x=-e/d$ (si numérateur non nul).
      \item Division euclidienne $\to$ Écriture $f(x)=\alpha x+\beta+\frac{\gamma}{dx+e}$ $\to$ Asymptote oblique $y=\alpha x+\beta$.
      \item Position par rapport à $\Delta$ : signe de $\frac{\gamma}{dx+e}$.
      \item Intersection avec $\Delta$ : résoudre $\gamma=0$.
      \end{itemize}
\item \textbf{Dérivée :} Calculer $f'(x)=\frac{u'v-uv'}{v^2}$, simplifier le numérateur (trinôme $N(x)$).
\item \textbf{Signe de $f'$ :} Étudier le signe de $N(x)$ (discriminant, racines) sur chaque intervalle de $D_f$.
\item \textbf{Tableau de variations :} Placer $-\infty$, les racines de $N(x)$, la valeur interdite (double barre $||$), $+\infty$. Indiquer variations ($\nearrow/\searrow$) et valeurs (images des extrema, limites aux bornes).
\item \textbf{Points particuliers :} Intersection avec l'axe des ordonnées ($f(0)$ si $0\in D_f$), intersection avec l'axe des abscisses (racines du numérateur).
\item \textbf{Tracé :} Repère, asymptotes (pointillés), points remarquables, branches respectant variations et position par rapport à $\Delta$.
\end{enumerate}
\end{methode}

\courssub{Exemple chiffré : étude complète}

\begin{exoresolu}[Étude complète de $f(x)=\dfrac{x^2+2x+2}{x+1}$]
Étudier $f$ et tracer sa courbe représentative $\mathcal{C}_f$.
\tcblower
\textbf{1. Ensemble de définition.} $x+1=0\iff x=-1$, donc $D_f=\mathbb{R}\setminus\{-1\}$. Le numérateur en $-1$ vaut $1\neq0$ : pas de simplification.

\textbf{2. Limites aux bornes.}
En $-1$ : le numérateur vaut $1>0$ et $x+1$ change de signe, d'où
\[ \lim_{\substack{x\to -1\\ x<-1}}f(x)=-\infty, \qquad \lim_{\substack{x\to -1\\ x>-1}}f(x)=+\infty. \]
La droite $x=-1$ est asymptote verticale.
En l'infini : $\displaystyle\lim_{x\to\pm\infty}f(x)=\lim_{x\to\pm\infty}\dfrac{x^2}{x}=\pm\infty$.

\textbf{3. Asymptote oblique.} Division : $x^2+2x+2=(x+1)(x+1)+1$, donc
\[ f(x)=x+1+\dfrac{1}{x+1}. \]
Comme $\displaystyle\lim_{x\to\pm\infty}\bigl[f(x)-(x+1)\bigr]=\lim_{x\to\pm\infty}\dfrac{1}{x+1}=0$, la droite $\Delta : y=x+1$ est asymptote oblique en $-\infty$ et en $+\infty$.
Position : $f(x)-(x+1)=\dfrac{1}{x+1}$, positif si $x>-1$ (courbe au-dessus de $\Delta$), négatif si $x<-1$ (courbe en dessous).
Intersection avec $\Delta$ : $\frac{1}{x+1}=0$ impossible, pas d'intersection.

\textbf{4. Dérivée.} Avec $u=x^2+2x+2$, $v=x+1$ :
\[ f'(x)=\dfrac{(2x+2)(x+1)-(x^2+2x+2)}{(x+1)^2}=\dfrac{x^2+2x}{(x+1)^2}=\dfrac{x(x+2)}{(x+1)^2}. \]
Le signe de $f'$ est celui de $x(x+2)$ : positif sur $]-\infty;-2[$ et $]0;+\infty[$, négatif sur $]-2;-1[$ et $]-1;0[$, nul en $-2$ et $0$.

\textbf{5. Valeurs remarquables.} $f(-2)=-2$ (maximum local), $f(0)=2$ (minimum local). $f(0)=2$ (intersection avec Oy). Numérateur $x^2+2x+2=0$ : $\Delta=4-8<0$, pas d'intersection avec Ox.

\textbf{6. Tableau de variations.}
\begin{center}
\begin{tabular}{|c|cccc|c||c|cccc|}
\hline
$x$ & $-\infty$ & & $-2$ & & $-\infty$ & $+\infty$ & & $0$ & & $+\infty$\\
\hline
$f'(x)$ & & $+$ & $0$ & $-$ & $||$ & $-$ & $0$ & $+$ & & \\
\hline
$f(x)$ & $-\infty$ & $\nearrow$ & $-2$ & $\searrow$ & $||$ & $+\infty$ & $\searrow$ & $2$ & $\nearrow$ & $+\infty$\\
\hline
\end{tabular}
\end{center}

\textbf{7. Tracé.} On place les asymptotes $x=-1$ et $y=x+1$, les points $(-2;-2)$ et $(0;2)$ à tangente horizontale, le point $(0;2)$ sur l'axe des ordonnées, puis on trace chaque branche en respectant la position par rapport à $\Delta$.
\end{exoresolu}

\courssub{Cas particulier : la fonction homographique}

\begin{definition}[Fonction homographique (rappel de Première)]
On appelle \cle{fonction homographique} toute fonction de la forme
\[ f(x)=\dfrac{ax+b}{cx+d} \quad \text{avec } c\neq 0 \text{ et } ad-bc\neq 0. \]
Sa courbe représentative est une \cle{hyperbole} de centre $\Omega\left(-\dfrac{d}{c}\,;\,\dfrac{a}{c}\right)$, admettant deux asymptotes : la droite verticale $x=-\dfrac{d}{c}$ et la droite horizontale $y=\dfrac{a}{c}$.
\end{definition}

La condition $ad-bc\neq0$ garantit que la fraction ne se simplifie pas en une constante. On peut voir l'homographique comme le cas « dégénéré » de notre étude : le numérateur étant de degré 1, la division euclidienne donne un quotient \emph{constant} $\dfrac{a}{c}$, d'où une asymptote \emph{horizontale} au lieu d'une asymptote oblique. Par exemple, $f(x)=\dfrac{2x+1}{x-1}=2+\dfrac{3}{x-1}$ : hyperbole de centre $\Omega(1;2)$.

\begin{savaistu}[Les hyperboles autour de nous]
La loi d'Ohm généralisée, la relation entre la dose d'engrais et le rendement d'un champ de maïs à l'approche de la saturation, ou encore le temps de trajet en fonction de la vitesse ($t=\dfrac{d}{v}$, une hyperbole !) sont des situations homographiques. Les fonctions rationnelles du type étudié dans cette section apparaissent dès que le numérateur croît plus vite : coût moyen de production, intensité dans un circuit avec résistance variable, etc.
\end{savaistu}

\begin{exercice}[Pour t'entraîner]
Soit $g(x)=\dfrac{x^2-3x+3}{x-2}$.
\begin{enumerate}
\item Déterminer l'ensemble de définition de $g$ et les limites aux bornes.
\item Montrer que $g(x)=x-1+\dfrac{1}{x-2}$ et en déduire les asymptotes de $\mathcal{C}_g$.
\item Étudier la position de $\mathcal{C}_g$ par rapport à son asymptote oblique et chercher une éventuelle intersection.
\item Calculer $g'(x)$, dresser le tableau de variations et tracer $\mathcal{C}_g$.
\end{enumerate}
\end{exercice}

\begin{corrige}[Exercice 1]
\textbf{1.} $D_g=\mathbb{R}\setminus\{2\}$. Numérateur en $2$ : $4-6+3=1>0$. Donc $\displaystyle\lim_{\substack{x\to2\\x<2}}g=-\infty$ et $\displaystyle\lim_{\substack{x\to2\\x>2}}g=+\infty$ : asymptote verticale $x=2$. En l'infini, $g(x)\sim x$ donc les limites sont $-\infty$ en $-\infty$ et $+\infty$ en $+\infty$.

\textbf{2.} $(x-2)(x-1)+1=x^2-3x+2+1=x^2-3x+3$ : l'égalité est vérifiée. Comme $\dfrac{1}{x-2}\to0$ en $\pm\infty$, la droite $y=x-1$ est asymptote oblique.

\textbf{3.} $g(x)-(x-1)=\dfrac{1}{x-2}$ : courbe au-dessus pour $x>2$, en dessous pour $x<2$. Pas d'intersection car $\gamma=1\neq0$.

\textbf{4.} $g'(x)=\dfrac{(2x-3)(x-2)-(x^2-3x+3)}{(x-2)^2}=\dfrac{x^2-4x+3}{(x-2)^2}=\dfrac{(x-1)(x-3)}{(x-2)^2}$. Signe : $g'$ positive sur $]-\infty;1]$ et $[3;+\infty[$, négative sur $[1;2[$ et $]2;3]$. Extrema : $g(1)=-1$ (maximum local), $g(3)=3$ (minimum local). Le tableau comporte une double barre en $x=2$, et le tracé s'appuie sur les asymptotes $x=2$ et $y=x-1$.
\end{corrige}

\begin{aretenir}[L'essentiel de la section]
\begin{itemize}
\item $f(x)=\dfrac{ax^2+bx+c}{dx+e}$ est définie sur $\mathbb{R}\setminus\left\{-\dfrac{e}{d}\right\}$ ; la droite $x=-\dfrac{e}{d}$ est asymptote verticale (sauf simplification).
\item La division euclidienne donne $f(x)=\alpha x+\beta+\dfrac{\gamma}{dx+e}$ : la droite $y=\alpha x+\beta$ est asymptote oblique, et le signe de $\dfrac{\gamma}{dx+e}$ donne la position de la courbe.
\item $f'(x)$ est un quotient dont le numérateur est un trinôme : son signe pilote les variations.
\item Le tableau de variations doit toujours contenir la double barre à la valeur interdite.
\item Intersections : avec Oy ($f(0)$), avec Ox (racines du numérateur), avec $\Delta$ ($\gamma=0$).
\end{itemize}
\end{aretenir}

\coursec{Étude des fonctions du type ln(ax + b)}

Dans la section précédente, nous avons appris à étudier les fonctions rationnelles du type $\dfrac{ax^2+bx+c}{dx+e}$ : recherche de la valeur interdite, limites, asymptotes verticale et oblique. Nous quittons maintenant le monde des quotients de polynômes pour entrer dans celui des \cle{fonctions transcendantes}. La première d'entre elles, la fonction logarithme népérien, est un outil formidable : elle transforme les produits en sommes, les puissances en produits, et c'est elle qui permet de mesurer l'intensité d'un séisme ou l'acidité d'une solution. Dans cette section, nous allons mener l'étude complète des fonctions du type $f(x)=\ln(ax+b)$, où $a$ et $b$ sont des réels avec $a\neq 0$.

\courssub{Rappels sur la fonction logarithme népérien}

\textbf{Intuition.} Lorsqu'une grandeur croît très vite (population, épidémie, capital avec intérêts composés), il est difficile de la représenter sur un graphique : les valeurs explosent. Le logarithme \og compresse \fg{} les grandes valeurs tout en respectant l'ordre : c'est une fonction croissante qui croît de plus en plus lentement.

\begin{definition}[Fonction logarithme népérien]
La fonction \cle{logarithme népérien}, notée $\ln$, est l'unique fonction définie et dérivable sur $]0\,;\,+\infty[$ vérifiant :
\[ \ln(1)=0 \qquad \text{et} \qquad \ln'(x)=\frac{1}{x} \quad \text{pour tout } x>0. \]
C'est la \emph{fonction réciproque} de la fonction exponentielle : pour tout $x>0$ et tout réel $y$,
\[ y=\ln(x) \iff x=e^{y}. \]
\end{definition}

\begin{propriete}[Propriétés algébriques du logarithme]
Pour tous réels $u>0$, $v>0$ et tout entier $n$ :
\[ \ln(uv)=\ln u+\ln v \qquad \ln\!\left(\frac{u}{v}\right)=\ln u-\ln v \qquad \ln(u^n)=n\ln u \qquad \ln\!\left(\frac{1}{u}\right)=-\ln u. \]
En particulier $\ln(e)=1$, et pour tout réel $x$ : $\ln(e^x)=x$ ; pour tout $x>0$ : $e^{\ln x}=x$.
\end{propriete}

\begin{aretenir}[Signe et valeurs remarquables de ln]
\begin{itemize}
\item $\ln(x)>0 \iff x>1$ ; $\ln(x)<0 \iff 0<x<1$ ; $\ln(1)=0$.
\item Limites de référence : $\displaystyle\lim_{x\to 0^+}\ln(x)=-\infty$ (asymptote verticale d'équation $x=0$) et $\displaystyle\lim_{x\to +\infty}\ln(x)=+\infty$.
\item La fonction $\ln$ est \textbf{strictement croissante} sur $]0\,;\,+\infty[$, car sa dérivée $\frac{1}{x}$ est strictement positive sur cet intervalle.
\end{itemize}
\end{aretenir}

\courssub{Ensemble de définition de f(x) = ln(ax + b)}

\textbf{Intuition.} La fonction $\ln$ n'accepte que des nombres \emph{strictement positifs}. Donc $f(x)=\ln(ax+b)$ n'existe que lorsque la quantité sous le logarithme, $ax+b$, est strictement positive. Tout le travail consiste à résoudre une inéquation du premier degré.

\begin{methode}[Déterminer l'ensemble de définition de ln(ax + b)]
\begin{enumerate}
\item Écrire la condition d'existence : $ax+b>0$ (inéquation \textbf{stricte}).
\item Isoler $x$ : $ax>-b$.
\item Diviser par $a$ \textbf{en tenant compte de son signe} :
\begin{itemize}
\item si $a>0$, le sens de l'inégalité est conservé : $x>-\dfrac{b}{a}$, donc $D_f=\left]-\dfrac{b}{a}\,;\,+\infty\right[$ ;
\item si $a<0$, le sens de l'inégalité est renversé : $x<-\dfrac{b}{a}$, donc $D_f=\left]-\infty\,;\,-\dfrac{b}{a}\right[$.
\end{itemize}
\item Conclure en précisant que la borne $-\dfrac{b}{a}$ est \emph{exclue} (crochet ouvert).
\end{enumerate}
\end{methode}

\begin{exemplebox}[Deux cas selon le signe de a]
\begin{itemize}
\item Pour $f(x)=\ln(3x-6)$ : $3x-6>0 \iff x>2$, donc $D_f=]2\,;\,+\infty[$.
\item Pour $g(x)=\ln(-2x+8)$ : $-2x+8>0 \iff x<4$ (on divise par $-2<0$, on change le sens de l'inégalité), donc $D_g=]-\infty\,;\,4[$.
\end{itemize}
\end{exemplebox}

\begin{attention}[Pièges sur l'ensemble de définition]
\begin{itemize}
\item N'écris jamais $ax+b\geq 0$ : le logarithme de $0$ n'existe pas. L'inéquation est \textbf{stricte}.
\item La valeur $x=-\dfrac{b}{a}$ n'est pas une \og valeur interdite \fg{} au sens des fractions : c'est une \textbf{borne de l'ensemble de définition}, et elle donnera une asymptote verticale.
\item N'oublie pas de renverser le sens de l'inégalité quand $a<0$.
\end{itemize}
\end{attention}

\courssub{Limites aux bornes et asymptote verticale}

Les bornes de $D_f$ sont $-\dfrac{b}{a}$ et l'infini. On utilise la composition des limites avec les limites de référence de $\ln$.

\begin{propriete}[Limites de ln(ax + b) aux bornes]
Soit $f(x)=\ln(ax+b)$ avec $a\neq 0$.
\begin{itemize}
\item Quand $x$ tend vers $-\dfrac{b}{a}$ (du côté où $f$ est définie), $ax+b$ tend vers $0$ par valeurs strictement positives, donc :
\[ \lim_{x\to -\frac{b}{a}} f(x)=-\infty. \]
La droite d'équation $x=-\dfrac{b}{a}$ est \cle{asymptote verticale} à la courbe de $f$.
\item À l'autre borne : si $a>0$, $\displaystyle\lim_{x\to+\infty}f(x)=+\infty$ ; si $a<0$, $\displaystyle\lim_{x\to-\infty}f(x)=+\infty$.
\end{itemize}
\end{propriete}

\begin{exemplebox}[Limites pour f(x) = ln(3x - 6)]
$D_f=]2\,;\,+\infty[$. Quand $x\to 2^+$, $3x-6\to 0^+$, donc $\displaystyle\lim_{x\to 2^+}\ln(3x-6)=-\infty$ : la droite d'équation $x=2$ est asymptote verticale. Quand $x\to+\infty$, $3x-6\to+\infty$, donc $\displaystyle\lim_{x\to+\infty}\ln(3x-6)=+\infty$.
\end{exemplebox}

\courssub{Dérivée et sens de variation}

\textbf{Intuition.} La fonction $f(x)=\ln(ax+b)$ est la \emph{composée} de la fonction affine $u(x)=ax+b$ suivie du logarithme. Pour dériver une telle composition, il faut une règle générale que nous établissons maintenant.

\begin{experience}[Démonstration guidée : la dérivée de ln(u)]
Soit $u$ une fonction dérivable et strictement positive sur un intervalle $I$, et $f=\ln\circ u$, c'est-à-dire $f(x)=\ln(u(x))$.
\begin{enumerate}
\item Rappelons la formule de dérivation des fonctions composées : $(v\circ u)'(x)=u'(x)\times v'(u(x))$.
\item Ici $v=\ln$, donc $v'(X)=\dfrac{1}{X}$ pour tout $X>0$.
\item Appliquons : $f'(x)=u'(x)\times \dfrac{1}{u(x)}$.
\item Conclusion : $\left(\ln u\right)'=\dfrac{u'}{u}$.
\end{enumerate}
Cette formule est exigible au Baccalauréat : sache la restituer et l'appliquer.
\end{experience}

\begin{propriete}[Dérivée et monotonie de ln(ax + b)]
Soit $f(x)=\ln(ax+b)$ avec $a\neq 0$, définie sur $D_f$. Alors $f$ est dérivable sur $D_f$ et :
\[ f'(x)=\frac{a}{ax+b}. \]
Comme $ax+b>0$ sur tout $D_f$, le signe de $f'(x)$ est \textbf{exactement le signe de $a$}. Par conséquent :
\begin{itemize}
\item si $a>0$, $f$ est \textbf{strictement croissante} sur $D_f$ ;
\item si $a<0$, $f$ est \textbf{strictement décroissante} sur $D_f$.
\end{itemize}
\end{propriete}

\begin{aretenir}[Résultat-clé]
Pour $f(x)=\ln(ax+b)$, il n'y a \textbf{jamais d'extremum} : la dérivée ne s'annule pas (son numérateur $a$ est une constante non nulle). La fonction est monotone sur tout son ensemble de définition, et son sens de variation se lit directement sur le signe de $a$.
\end{aretenir}

\textbf{Les deux tableaux de variations possibles.}

\begin{popfigure}
\begin{center}
\begin{tikzpicture}[scale=1]
% Cas a > 0
\node[font=\bfseries] at (3,2.6) {Cas $a>0$ : $D_f=\left]-\frac{b}{a}\,;\,+\infty\right[$};
\draw[thick] (0,0) rectangle (6,2);
\draw[thick] (1.4,0) -- (1.4,2);
\draw[thick] (0,1.4) -- (6,1.4);
\node at (0.7,1.7) {$x$};
\node at (2.2,1.7) {$-\frac{b}{a}$};
\node at (5.5,1.7) {$+\infty$};
\draw[thick,popink] (2.0,1.4) -- (2.0,0);
\draw[thick,popink] (2.4,1.4) -- (2.4,0);
\node at (0.7,0.7) {$f$};
\node[popink] at (2.6,0.25) {$-\infty$};
\node[popink] at (5.3,1.15) {$+\infty$};
\draw[->,very thick,popink] (2.75,0.35) -- (5.15,1.05);
% Cas a < 0
\node[font=\bfseries] at (10.5,2.6) {Cas $a<0$ : $D_f=\left]-\infty\,;\,-\frac{b}{a}\right[$};
\draw[thick] (7.5,0) rectangle (13.5,2);
\draw[thick] (8.9,0) -- (8.9,2);
\draw[thick] (7.5,1.4) -- (13.5,1.4);
\node at (8.2,1.7) {$x$};
\node at (9.6,1.7) {$-\infty$};
\node at (12.9,1.7) {$-\frac{b}{a}$};
\draw[thick,popblue] (12.6,1.4) -- (12.6,0);
\draw[thick,popblue] (13.0,1.4) -- (13.0,0);
\node at (8.2,0.7) {$f$};
\node[popblue] at (9.8,1.15) {$+\infty$};
\node[popblue] at (12.4,0.25) {$-\infty$};
\draw[->,very thick,popblue] (9.95,1.05) -- (12.25,0.35);
\end{tikzpicture}
\end{center}
\legende{Tableaux de variations de $f(x)=\ln(ax+b)$ selon le signe de $a$. La double barre matérialise la borne exclue $x=-\frac{b}{a}$ (asymptote verticale).}
\end{popfigure}

\courssub{Tracé de la courbe : transformations de la courbe de référence}

La courbe de $f(x)=\ln(ax+b)$ s'obtient à partir de celle de $y=\ln(x)$ par des transformations géométriques. Par exemple, pour $x>2$ :
\[ \ln(2x-4)=\ln\!\big(2(x-2)\big)=\ln 2+\ln(x-2). \]
La courbe de $y=\ln(2x-4)$ est donc l'image de celle de $y=\ln(x)$ par la translation de vecteur $2\vec{\imath}$ suivie de la translation de vecteur $(\ln 2)\,\vec{\jmath}$. Concrètement : l'asymptote verticale passe de $x=0$ à $x=2$, et la courbe garde la même allure, croissante et de plus en plus \og plate \fg{}.

\begin{popfigure}
\begin{center}
\begin{tikzpicture}
\begin{axis}[
  width=\linewidth, height=7.5cm,
  axis lines=middle,
  xlabel={$x$}, ylabel={$y$},
  xmin=-1.5, xmax=9, ymin=-4, ymax=3.5,
  xtick={1,2,3,4,5,6,7,8}, ytick={-3,-2,-1,1,2,3},
  grid=both, grid style={popline!40},
  legend style={at={(0.98,0.35)},anchor=east,font=\small}
]
\addplot[popcurve,popblue,domain=0.02:9,samples=200]{ln(x)};
\addlegendentry{$y=\ln(x)$}
\addplot[popcurve,popink,domain=2.005:9,samples=200]{ln(2*x-4)};
\addlegendentry{$y=\ln(2x-4)$}
\addplot[poporange,dashed,thick] coordinates {(2,-4) (2,3.5)};
\addlegendentry{asymptote $x=2$}
\addplot[only marks,mark=*,popblue] coordinates {(1,0)};
\addplot[only marks,mark=*,popink] coordinates {(2.5,0)};
\end{axis}
\end{tikzpicture}
\end{center}
\legende{Courbes de $y=\ln(x)$ (bleue, asymptote $x=0$, passe par $(1\,;\,0)$) et de $y=\ln(2x-4)$ (rouge, asymptote $x=2$, passe par $\left(\frac{5}{2}\,;\,0\right)$).}
\end{popfigure}

\begin{methode}[Plan d'étude complet de f(x) = ln(ax + b)]
\begin{enumerate}
\item \textbf{Ensemble de définition} : résoudre $ax+b>0$.
\item \textbf{Limites aux bornes} : en $-\frac{b}{a}$ (résultat $-\infty$, asymptote verticale) et à l'infini (résultat $+\infty$).
\item \textbf{Dérivée} : $f'(x)=\dfrac{a}{ax+b}$, du signe de $a$.
\item \textbf{Tableau de variations} : monotone, avec la double barre en $-\frac{b}{a}$.
\item \textbf{Points remarquables} : résoudre $f(x)=0 \iff ax+b=1$, soit $x=\dfrac{1-b}{a}$ (intersection avec l'axe des abscisses, à condition que cette valeur appartienne à $D_f$, ce qui est toujours le cas car $a\times\frac{1-b}{a}+b=1>0$).
\item \textbf{Tracé} : asymptote verticale, point d'intersection avec l'axe des abscisses, allure de logarithme.
\end{enumerate}
\end{methode}

\begin{exoresolu}[Étude complète de f(x) = ln(3x - 6) et inéquation]
\textbf{Énoncé.} Soit $f$ la fonction définie par $f(x)=\ln(3x-6)$ et $\mathcal{C}_f$ sa courbe représentative dans un repère orthogonal.
\begin{enumerate}
\item Déterminer l'ensemble de définition de $f$.
\item Calculer les limites de $f$ aux bornes de $D_f$ et interpréter graphiquement.
\item Étudier les variations de $f$ et dresser son tableau de variations.
\item Déterminer les coordonnées du point d'intersection de $\mathcal{C}_f$ avec l'axe des abscisses.
\item Résoudre dans $\mathbb{R}$ l'inéquation $\ln(3x-6)\geq 0$.
\end{enumerate}
\tcblower
\textbf{Solution détaillée.}
\begin{enumerate}
\item $f(x)$ existe si et seulement si $3x-6>0 \iff x>2$. Donc $D_f=]2\,;\,+\infty[$.
\item Quand $x\to 2^+$ : $3x-6\to 0^+$, donc $\displaystyle\lim_{x\to 2^+}f(x)=-\infty$. La droite d'équation $x=2$ est \textbf{asymptote verticale} à $\mathcal{C}_f$. Quand $x\to+\infty$ : $3x-6\to+\infty$, donc $\displaystyle\lim_{x\to+\infty}f(x)=+\infty$.
\item $f$ est de la forme $\ln(u)$ avec $u(x)=3x-6$, $u'(x)=3$. Donc :
\[ f'(x)=\frac{3}{3x-6}. \]
Sur $D_f$, $3x-6>0$, donc $f'(x)>0$ : $f$ est \textbf{strictement croissante} sur $]2\,;\,+\infty[$.
\begin{center}
\begin{tabular}{|c|ccc|}
\hline
$x$ & $2$ & & $+\infty$ \\
\hline
$f'(x)$ & \multicolumn{3}{c|}{$+$} \\
\hline
$f$ & $-\infty$ & $\nearrow$ & $+\infty$ \\
\hline
\end{tabular}
\end{center}
(avec une double barre en $x=2$).
\item $f(x)=0 \iff \ln(3x-6)=0 \iff 3x-6=e^0=1 \iff x=\dfrac{7}{3}$. Comme $\dfrac{7}{3}>2$, ce point appartient bien à $D_f$. La courbe coupe l'axe des abscisses en $A\left(\dfrac{7}{3}\,;\,0\right)$.
\item \textbf{Condition d'existence} : $x>2$. Sous cette condition :
\[ \ln(3x-6)\geq 0 \iff 3x-6\geq e^0=1 \iff x\geq \frac{7}{3}, \]
car la fonction exponentielle est strictement croissante (ou, de façon équivalente, parce que $\ln$ est croissante et $\ln(1)=0$). Comme $\frac{7}{3}>2$, la condition d'existence est automatiquement satisfaite et l'ensemble des solutions est :
\[ \mathcal{S}=\left[\frac{7}{3}\,;\,+\infty\right[. \]
\end{enumerate}
\end{exoresolu}

\courssub{Équations et inéquations avec ln(ax + b)}

\begin{propriete}[Résolution des équations et inéquations logarithmiques]
Pour résoudre une équation ou une inéquation faisant intervenir $\ln(ax+b)$ :
\begin{enumerate}
\item on commence \textbf{toujours} par poser la condition d'existence $ax+b>0$ ;
\item on utilise ensuite la stricte croissance des fonctions $\ln$ et $\exp$ :
\[ \ln(ax+b)=k \iff ax+b=e^{k}, \qquad \ln(ax+b)\geq k \iff ax+b\geq e^{k} \quad (\text{sens conservé}) ; \]
\item on vérifie enfin que les solutions trouvées appartiennent à l'ensemble de définition.
\end{enumerate}
\end{propriete}

\begin{exemplebox}[Équation avec second membre non nul]
Résolvons $\ln(2x+1)=3$.
Condition : $2x+1>0 \iff x>-\dfrac{1}{2}$. Alors :
\[ \ln(2x+1)=3 \iff 2x+1=e^{3} \iff x=\frac{e^{3}-1}{2}\approx 9{,}54. \]
Cette valeur est bien supérieure à $-\frac{1}{2}$, donc $\mathcal{S}=\left\{\dfrac{e^{3}-1}{2}\right\}$.
\end{exemplebox}

\begin{attention}[Erreurs fréquentes avec le logarithme]
\begin{itemize}
\item \textbf{Le logarithme ne transforme pas les sommes en produits !} Il est \emph{faux} d'écrire $\ln(ax+b)=\ln(ax)+\ln(b)$. La formule $\ln(uv)=\ln u+\ln v$ ne s'applique qu'à un \textbf{produit}. Par exemple $\ln(2x+3)$ ne se décompose pas ; en revanche $\ln(6x)=\ln 6+\ln x$ est correct (pour $x>0$).
\item Il n'existe aucune formule pour $\ln(u+v)$.
\item N'oublie pas la condition d'existence \emph{avant} de résoudre : une \og solution \fg{} hors de $D_f$ doit être rejetée.
\item Dans une inéquation, le sens est conservé lorsqu'on applique $\exp$ ou $\ln$ car ces fonctions sont \emph{strictement croissantes} ; ne change le sens que lorsque tu divises par un nombre strictement négatif.
\end{itemize}
\end{attention}

\begin{exercice}[À toi de jouer]
Soit $g(x)=\ln(-x+4)$.
\begin{enumerate}
\item Déterminer l'ensemble de définition de $g$.
\item Étudier les limites de $g$ aux bornes de $D_g$, sa dérivée et ses variations.
\item Résoudre dans $\mathbb{R}$ l'inéquation $g(x)\geq 2$.
\end{enumerate}
\end{exercice}

\begin{corrige}[Exercice : g(x) = ln(-x + 4)]
\begin{enumerate}
\item $-x+4>0 \iff x<4$, donc $D_g=]-\infty\,;\,4[$.
\item Quand $x\to 4^-$ : $-x+4\to 0^+$, donc $\displaystyle\lim_{x\to 4^-}g(x)=-\infty$ (asymptote verticale d'équation $x=4$). Quand $x\to-\infty$ : $-x+4\to+\infty$, donc $\displaystyle\lim_{x\to-\infty}g(x)=+\infty$. Dérivée : $g'(x)=\dfrac{-1}{-x+4}=\dfrac{1}{x-4}<0$ sur $D_g$ : $g$ est strictement décroissante (ici $a=-1<0$, conforme au cours).
\item Condition : $x<4$. Sous cette condition :
\[ g(x)\geq 2 \iff -x+4\geq e^{2} \iff x\leq 4-e^{2}\approx -3{,}39. \]
Comme $4-e^2<4$, la condition d'existence est satisfaite et $\mathcal{S}=\left]-\infty\,;\,4-e^{2}\right]$.
\end{enumerate}
\end{corrige}

\begin{savaistu}[Le logarithme mesure le monde : Richter et pH]
L'\textbf{échelle de Richter}, utilisée pour mesurer la magnitude des séismes, est logarithmique : un séisme de magnitude $6$ a une amplitude environ $10$ fois plus grande qu'un séisme de magnitude $5$, car chaque unité correspond à une multiplication de l'amplitude par $10$. Le logarithme permet d'afficher sur une même échelle des séismes imperceptibles et des catastrophes majeures. En chimie, le \textbf{pH} d'une solution est défini par $\mathrm{pH}=-\log[\mathrm{H}^+]$, où $\log$ est le logarithme décimal : l'eau pure a un pH de $7$, le jus de citron environ $2$. Sans le logarithme, il faudrait manipuler des nombres allant de $1$ à $10^{-14}$ ! Les sismologues qui surveillent l'activité du mont Cameroun, volcan actif, utilisent quotidiennement ces échelles logarithmiques.
\end{savaistu}

\begin{aretenir}[Fiche récapitulative de la section]
Pour $f(x)=\ln(ax+b)$, $a\neq 0$ :
\begin{itemize}
\item $D_f$ : résoudre $ax+b>0$ (attention au signe de $a$) ;
\item limite en $-\frac{b}{a}$ : $-\infty$ (asymptote verticale) ; limite à l'autre borne (infinie) : $+\infty$ ;
\item $f'(x)=\dfrac{a}{ax+b}$, du signe de $a$ : croissante si $a>0$, décroissante si $a<0$, jamais d'extremum ;
\item $\ln(ax+b)=k \iff ax+b=e^k$, toujours sous condition $ax+b>0$ ;
\item $\ln(ax+b)$ ne se décompose pas : pas de formule pour le logarithme d'une somme.
\end{itemize}
\end{aretenir}

\coursec{Étude des fonctions du type $e^{ax+b}$}

Dans la section précédente, nous avons étudié les fonctions du type $x \longmapsto \ln(ax+b)$. Nous avons vu que le logarithme népérien « transforme les produits en sommes » et qu'il croît très lentement vers $+\infty$. Nous allons maintenant rencontrer sa partenaire inséparable : la \cle{fonction exponentielle}. Si le logarithme « défait » ce que l'exponentielle « fait », alors ces deux fonctions sont les deux faces d'une même pièce. Cette dualité n'est pas qu'une jolie image : c'est l'outil qui nous permettra de résoudre toutes les équations et inéquations de cette section.

L'exponentielle modélise les phénomènes de \emph{croissance rapide} : propagation d'une épidémie, intérêts composés d'un compte épargne à la microfinance, croissance d'une population de bactéries dans un laboratoire de l'Université de Yaoundé, ou encore décroissance radioactive. La maîtriser, c'est disposer d'un des outils les plus puissants du programme de Terminale.

\courssub{La fonction exponentielle : réciproque du logarithme népérien}

\textbf{Intuition.} Rappelons ce que nous savons du logarithme népérien : la fonction $\ln$ est une bijection strictement croissante de $]0\,;+\infty[$ sur $\mathbb{R}$. Dire qu'une fonction est une bijection, c'est dire que chaque valeur de l'ensemble d'arrivée est atteinte une et une seule fois. On peut donc « remonter » : à tout réel $y$, on associe l'unique réel strictement positif dont le logarithme vaut $y$. Cette « remontée » définit une nouvelle fonction : l'exponentielle.

\begin{definition}[Fonction exponentielle]
La fonction logarithme népérien $\ln$ étant une bijection de $]0\,;+\infty[$ sur $\mathbb{R}$, elle admet une \cle{bijection réciproque}, appelée \cle{fonction exponentielle} et notée $\exp$, définie de $\mathbb{R}$ dans $]0\,;+\infty[$. On note :
\[ \exp(x) = e^x. \]
Par définition d'une bijection réciproque, on a les deux relations fondamentales :
\[ \boxed{\ \ln(e^x) = x \ \text{pour tout } x \in \mathbb{R} \quad\text{et}\quad e^{\ln x} = x \ \text{pour tout } x > 0.\ } \]
En particulier $e^0 = 1$ (car $\ln 1 = 0$) et $e^1 = e \approx \num{2,718}$ (car $\ln e = 1$).
\end{definition}

\begin{aretenir}[Propriétés algébriques de l'exponentielle]
L'exponentielle « transforme les sommes en produits » (c'est le miroir exact du logarithme). Pour tous réels $a$ et $b$ :
\[ e^{a+b} = e^a \times e^b \ ; \qquad e^{a-b} = \dfrac{e^a}{e^b} \ ; \qquad e^{-a} = \dfrac{1}{e^a} \ ; \qquad e^{na} = \left(e^a\right)^n \ (n \in \mathbb{Z}). \]
Ces formules se démontrent en appliquant $\ln$ aux deux membres et en utilisant les propriétés algébriques du logarithme.
\end{aretenir}

\begin{exemplebox}[Calculs avec les propriétés algébriques]
Simplifions $A = \dfrac{e^{2x+1} \times e^{3-x}}{e^{x+2}}$.
\begin{align*}
A &= \dfrac{e^{(2x+1)+(3-x)}}{e^{x+2}} = \dfrac{e^{x+4}}{e^{x+2}} = e^{(x+4)-(x+2)} = e^2.
\end{align*}
Le résultat est constant, indépendant de $x$ : $A = e^2 \approx \num{7,39}$.
\end{exemplebox}

\begin{propriete}[Stricte positivité de l'exponentielle — admise]
Pour tout réel $x$, on a $\boxed{e^x > 0}$.
\end{propriete}

Cette propriété est \emph{admise} au programme de Terminale, mais elle est très naturelle : l'exponentielle est définie comme réciproque de $\ln$, donc elle prend ses valeurs dans l'ensemble de départ de $\ln$, c'est-à-dire $]0\,;+\infty[$. Une exponentielle ne s'annule \textbf{jamais} et n'est \textbf{jamais} négative. C'est le point crucial pour l'étude du signe des dérivées dans toute cette section.

\begin{savaistu}[Le nombre $e$, un nombre célèbre]
Le nombre $e \approx \num{2,71828}$ porte le nom du mathématicien suisse Leonhard Euler (1707--1783), bien qu'il ait été introduit implicitement par John Napier et Jacob Bernoulli en étudiant les intérêts composés. Comme $\pi$, c'est un nombre irrationnel : son écriture décimale est infinie et non périodique. Il apparaît dès qu'un phénomène croît « proportionnellement à lui-même » : si un capital placé à la microfinance croît continuellement à un taux de $100\,\%$ l'an, un franc devient exactement $e$ francs au bout d'un an.
\end{savaistu}

\courssub{Symétrie des courbes de $\ln$ et de $\exp$}

Puisque $\exp$ est la réciproque de $\ln$, leurs courbes sont \cle{symétriques par rapport à la droite d'équation $y = x$} (la première bissectrice). En effet, dire que le point $(a\,;b)$ appartient à la courbe de $\exp$ signifie $b = e^a$, c'est-à-dire $a = \ln b$, donc le point $(b\,;a)$ appartient à la courbe de $\ln$ : on a échangé les coordonnées, ce qui est exactement la symétrie par rapport à $y = x$.

\begin{popfigure}
\begin{center}
\begin{tikzpicture}[scale=0.85, >=Stealth]
\draw[->, popmuted] (-3.2,0) -- (4.2,0) node[below left] {$x$};
\draw[->, popmuted] (0,-1.8) -- (0,4.2) node[below left] {$y$};
\draw[popmuted, dashed] (-1.8,-1.8) -- (4.2,4.2) node[below right, pos=0.92] {$y=x$};
\draw[popblue, very thick, domain=-3.0:1.45, samples=150] plot (\x,{exp(\x)});
\node[popblue] at (-2.2,0.35) {$y=e^x$};
\draw[poporange, very thick, domain=0.06:4.0, samples=150] plot (\x,{ln(\x)});
\node[poporange] at (3.3,-0.9) {$y=\ln x$};
\fill[popdark] (0,1) circle (1.6pt) node[above left=1pt] {$(0\,;1)$};
\fill[popdark] (1,0) circle (1.6pt) node[below right=1pt] {$(1\,;0)$};
\fill[popdark] (1,2.718) circle (1.6pt);
\fill[popdark] (2.718,1) circle (1.6pt);
\draw[popmuted, dotted] (2.718,0) node[below] {$e$} -- (2.718,1);
\draw[popmuted, dotted] (0,2.718) node[left] {$e$} -- (1,2.718);
\end{tikzpicture}
\end{center}
\legende{Les courbes de $\exp$ et de $\ln$ sont symétriques par rapport à la droite $y = x$. Les points $(0\,;1)$ et $(1\,;0)$, $(1\,;e)$ et $(e\,;1)$ se correspondent.}
\end{popfigure}

\courssub{Ensemble de définition et limites de $f(x) = e^{ax+b}$}

Contrairement aux fonctions $\ln(ax+b)$ de la section précédente, dont l'ensemble de définition exigeait la condition $ax+b > 0$, la situation est ici d'une simplicité réjouissante.

\begin{propriete}[Ensemble de définition]
Soient $a$ et $b$ deux réels avec $a \neq 0$. La fonction $f : x \longmapsto e^{ax+b}$ est définie sur $\mathbb{R}$ tout entier :
\[ \mathcal{D}_f = \mathbb{R} = ]-\infty\,;+\infty[. \]
En effet, $ax+b$ est défini pour tout réel $x$, et l'exponentielle accepte \emph{tout} réel en entrée.
\end{propriete}

\begin{attention}[Ne pas confondre les conditions d'existence]
Pour $\ln(ax+b)$, il faut résoudre $ax+b > 0$ : l'ensemble de définition est une demi-droite. Pour $e^{ax+b}$, il n'y a \textbf{aucune condition} : $\mathcal{D}_f = \mathbb{R}$. Écrire « $f$ est définie si $ax+b > 0$ » pour une exponentielle est une erreur classique qui coûte des points au Baccalauréat.
\end{attention}

\textbf{Limites.} On connaît les limites de référence de l'exponentielle :
\[ \lim_{X \to +\infty} e^X = +\infty \qquad\text{et}\qquad \lim_{X \to -\infty} e^X = 0. \]
Pour $f(x) = e^{ax+b}$, on pose $X = ax+b$ et on applique le théorème de composition des limites. Tout dépend alors du \cle{signe de $a$} :

\begin{center}
\begin{tabularx}{\linewidth}{|l|X|X|}
\hline
\rowcolor{popblueL} & \textbf{Si $a > 0$} & \textbf{Si $a < 0$} \\
\hline
En $+\infty$ & $ax+b \to +\infty$, donc $\lim\limits_{x\to+\infty} e^{ax+b} = +\infty$ & $ax+b \to -\infty$, donc $\lim\limits_{x\to+\infty} e^{ax+b} = 0$ \\
\hline
En $-\infty$ & $ax+b \to -\infty$, donc $\lim\limits_{x\to-\infty} e^{ax+b} = 0$ & $ax+b \to +\infty$, donc $\lim\limits_{x\to-\infty} e^{ax+b} = +\infty$ \\
\hline
\end{tabularx}
\end{center}

\begin{propriete}[Asymptote horizontale]
Si $a > 0$, la droite d'équation $y = 0$ (l'axe des abscisses) est \cle{asymptote horizontale} à la courbe de $f$ en $-\infty$. Si $a < 0$, c'est en $+\infty$ que la droite $y = 0$ est asymptote horizontale. Dans tous les cas, la courbe « colle » à l'axe des abscisses d'un côté, sans jamais le toucher (puisque $e^{ax+b} > 0$).
\end{propriete}

\courssub{Dérivée et sens de variation}

\begin{experience}[Démonstration guidée : la dérivée de $e^u$]
\textbf{Objectif :} montrer que si $u$ est une fonction dérivable, alors $\left(e^u\right)' = u' \times e^u$.

\textbf{Étape 1 — Rappel.} La fonction $\exp$ est dérivable sur $\mathbb{R}$ et $\exp' = \exp$, c'est-à-dire $\left(e^x\right)' = e^x$. (Cela se démontre à partir de la dérivabilité de $\ln$ : la dérivée d'une bijection réciproque $g$ de $f$ est $g' = \dfrac{1}{f' \circ g}$, et comme $\ln' = \dfrac{1}{x}$, on obtient $\exp'(x) = \dfrac{1}{\frac{1}{\exp(x)}} = \exp(x)$.)

\textbf{Étape 2 — Dérivation d'une composée.} Soit $u$ une fonction dérivable sur un intervalle $I$. La fonction $e^u$ est la composée $x \mapsto \exp(u(x))$. La formule de dérivation des fonctions composées, $(v \circ u)' = u' \times (v' \circ u)$, appliquée avec $v = \exp$, donne :
\[ \left(e^{u(x)}\right)' = u'(x) \times \exp'(u(x)) = u'(x) \times e^{u(x)}. \]

\textbf{Étape 3 — Application à $u(x) = ax+b$.} Ici $u'(x) = a$, donc :
\[ \boxed{\ \left(e^{ax+b}\right)' = a\,e^{ax+b}.\ } \]
\end{experience}

\begin{propriete}[Dérivée et monotonie de $f(x) = e^{ax+b}$]
Soit $f(x) = e^{ax+b}$ avec $a \neq 0$. La fonction $f$ est dérivable sur $\mathbb{R}$ et :
\[ f'(x) = a\,e^{ax+b}. \]
Comme $e^{ax+b} > 0$ pour tout réel $x$, le signe de $f'(x)$ est \textbf{exactement le signe de $a$}. Par conséquent :
\begin{itemize}
\item si $a > 0$, alors $f$ est \cle{strictement croissante} sur $\mathbb{R}$ ;
\item si $a < 0$, alors $f$ est \cle{strictement décroissante} sur $\mathbb{R}$.
\end{itemize}
\end{propriete}

C'est un luxe rare en mathématiques : l'étude du signe de la dérivée est immédiate, sans tableau de signes, sans racines à chercher. L'exponentielle étant toujours strictement positive, elle ne « perturbe » jamais le signe.

\textbf{Tableaux de variations.} On en déduit les deux tableaux types.

\begin{center}
\textbf{Cas $a > 0$}\qquad
\begin{tabular}{|c|ccc|}
\hline
$x$ & $-\infty$ & & $+\infty$ \\
\hline
$f'(x) = a e^{ax+b}$ & & $+$ & \\
\hline
$f(x)$ & $0$ & $\nearrow$ & $+\infty$ \\
\hline
\end{tabular}
\qquad
\textbf{Cas $a < 0$}\qquad
\begin{tabular}{|c|ccc|}
\hline
$x$ & $-\infty$ & & $+\infty$ \\
\hline
$f'(x) = a e^{ax+b}$ & & $-$ & \\
\hline
$f(x)$ & $+\infty$ & $\searrow$ & $0$ \\
\hline
\end{tabular}
\end{center}

\courssub{Tracé de la courbe}

Pour tracer la courbe de $f(x) = e^{ax+b}$, on rassemble les informations : monotonie, limites, asymptote horizontale $y = 0$, et quelques points remarquables. Le point d'abscisse $0$ est précieux : $f(0) = e^b$, et la tangente en ce point a pour équation :
\[ y = f'(0)(x-0) + f(0) = a e^b\, x + e^b. \]
La courbe est toujours située \emph{strictement au-dessus} de l'axe des abscisses.

\begin{popfigure}
\begin{center}
\begin{tikzpicture}
\begin{axis}[
  width=0.95\linewidth, height=6.2cm,
  axis lines=middle, xlabel={$x$}, ylabel={$y$},
  xmin=-3, xmax=2.6, ymin=-0.5, ymax=6,
  xtick={-2,-1,1,2}, ytick={1,2,3,4,5},
  legend style={at={(0.02,0.97)}, anchor=north west, draw=popline, font=\small},
  grid=both, grid style={popline!40, very thin},
]
\addplot[popcurve, domain=-3:1.75, samples=200] {exp(x)};
\addlegendentry{$y=e^x$}
\addplot[poporange, very thick, domain=-3:0.85, samples=200] {exp(2*x)};
\addlegendentry{$y=e^{2x}$}
\addplot[popgreen, very thick, domain=-1.75:2.6, samples=200] {exp(-x)};
\addlegendentry{$y=e^{-x}$}
\draw[popmuted, dashed, thick] (axis cs:-3,0) -- (axis cs:2.6,0);
\node[popmuted, font=\small] at (axis cs:2.3,0.28) {asymptote $y=0$};
\end{axis}
\end{tikzpicture}
\end{center}
\legende{Comparaison de $y=e^x$, $y=e^{2x}$ (croissance plus rapide, $a=2>0$) et $y=e^{-x}$ (décroissante, $a=-1<0$). Toutes trois passent par $(0\,;1)$ et admettent $y=0$ comme asymptote horizontale d'un côté.}
\end{popfigure}

Observons : plus $a > 0$ est grand, plus la courbe « monte » vite vers $+\infty$ et « colle » vite à l'axe en $-\infty$. Pour $a < 0$, la courbe est le miroir : elle descend de $+\infty$ vers $0$. Remarquons aussi que $e^{-x}$ est la symétrique de $e^x$ par rapport à l'axe des ordonnées.

\courssub{Équations et inéquations avec l'exponentielle}

Voici où le lien exponentielle-logarithme devient un outil opérationnel. L'idée directrice : pour « descendre » l'inconnue $x$ coincée dans l'exposant, on applique $\ln$, qui « défait » l'exponentielle grâce à $\ln(e^u) = u$.

\begin{methode}[Résoudre $e^{ax+b} = k$]
\begin{enumerate}
\item \textbf{Vérifier le signe de $k$.} Comme une exponentielle est strictement positive :
   \begin{itemize}
   \item si $k \leq 0$ : l'équation n'a \textbf{aucune solution}, $S = \varnothing$ ;
   \item si $k > 0$ : on continue.
   \end{itemize}
\item \textbf{Appliquer $\ln$ aux deux membres} (licite car les deux membres sont strictement positifs) :
   \[ e^{ax+b} = k \iff \ln\left(e^{ax+b}\right) = \ln k \iff ax + b = \ln k. \]
\item \textbf{Résoudre l'équation du premier degré} obtenue : $x = \dfrac{\ln k - b}{a}$.
\end{enumerate}
Pour une \textbf{inéquation} $e^{ax+b} \leq k$ (avec $k > 0$), la démarche est identique, car $\ln$ est \emph{strictement croissante} : elle conserve le sens des inégalités. On obtient $ax + b \leq \ln k$, puis on divise par $a$ \textbf{en pensant à changer le sens de l'inégalité si $a < 0$}.
\end{methode}

\begin{attention}[Le piège du second membre négatif ou nul]
Face à $e^{2x-1} = -5$ ou $e^{x} = 0$, beaucoup d'élèves appliquent machinalement $\ln$ et écrivent « $2x-1 = \ln(-5)$ ». C'est une double erreur : $\ln(-5)$ n'existe pas, et surtout une exponentielle ne peut jamais valoir $-5$ ni $0$. La bonne réponse, immédiate et rédigée, est : « pour tout réel $x$, $e^{2x-1} > 0$, donc l'équation $e^{2x-1} = -5$ n'a pas de solution : $S = \varnothing$ ». De même, l'inéquation $e^{ax+b} \leq 0$ n'a aucune solution, tandis que $e^{ax+b} > 0$ est vraie pour tout $x \in \mathbb{R}$.
\end{attention}

\courssub{Exemple complet : étude de $f(x) = e^{2x-1}$}

\begin{exoresolu}[Étude complète et inéquation]
On considère la fonction $f$ définie par $f(x) = e^{2x-1}$.
\begin{enumerate}
\item Déterminer l'ensemble de définition de $f$ et ses limites aux bornes. Interpréter graphiquement.
\item Calculer $f'(x)$, en déduire le sens de variation et dresser le tableau de variations.
\item Donner l'équation de la tangente à la courbe au point d'abscisse $0$.
\item Résoudre dans $\mathbb{R}$ l'équation $e^{2x-1} = 3$, puis l'inéquation $e^{2x-1} \leq 3$.
\end{enumerate}
\tcblower
\textbf{1. Ensemble de définition et limites.} L'exponentielle est définie sur $\mathbb{R}$ et $2x-1$ est défini partout : $\mathcal{D}_f = \mathbb{R}$.

En $+\infty$ : $\lim\limits_{x\to+\infty}(2x-1) = +\infty$ et $\lim\limits_{X\to+\infty} e^X = +\infty$, donc par composition $\lim\limits_{x\to+\infty} f(x) = +\infty$.

En $-\infty$ : $\lim\limits_{x\to-\infty}(2x-1) = -\infty$ et $\lim\limits_{X\to-\infty} e^X = 0$, donc $\lim\limits_{x\to-\infty} f(x) = 0$. La droite $y = 0$ est asymptote horizontale à la courbe en $-\infty$.

\textbf{2. Dérivée et variations.} $f$ est de la forme $e^u$ avec $u(x) = 2x-1$, donc $u'(x) = 2$ :
\[ f'(x) = 2\,e^{2x-1}. \]
Pour tout réel $x$, $e^{2x-1} > 0$ et $2 > 0$, donc $f'(x) > 0$ : $f$ est \textbf{strictement croissante} sur $\mathbb{R}$.

\begin{center}
\begin{tabular}{|c|ccc|}
\hline
$x$ & $-\infty$ & & $+\infty$ \\
\hline
$f'(x)$ & & $+$ & \\
\hline
$f(x)$ & $0$ & $\nearrow$ & $+\infty$ \\
\hline
\end{tabular}
\end{center}

\textbf{3. Tangente en $0$.} $f(0) = e^{-1} = \dfrac{1}{e}$ et $f'(0) = 2e^{-1} = \dfrac{2}{e}$. L'équation de la tangente est :
\[ y = f'(0)\,x + f(0) = \dfrac{2}{e}\,x + \dfrac{1}{e}. \]

\textbf{4. Équation puis inéquation.} Comme $3 > 0$, on peut appliquer $\ln$ :
\begin{align*}
e^{2x-1} = 3 &\iff 2x - 1 = \ln 3 \iff x = \dfrac{1 + \ln 3}{2}.
\end{align*}
Donc $S = \left\{\dfrac{1+\ln 3}{2}\right\}$, soit numériquement $x \approx \num{1,05}$.

Pour l'inéquation, $\ln$ étant strictement croissante, elle conserve le sens :
\begin{align*}
e^{2x-1} \leq 3 &\iff 2x - 1 \leq \ln 3 \iff 2x \leq 1 + \ln 3 \iff x \leq \dfrac{1+\ln 3}{2}.
\end{align*}
(On divise par $2 > 0$ : le sens ne change pas.) Donc :
\[ S = \left]-\infty\,;\ \dfrac{1+\ln 3}{2}\right]. \]
\end{exoresolu}

\begin{attention}[Rédaction et erreurs fréquentes]
\begin{itemize}
\item Ne jamais écrire $\left(e^{ax+b}\right)' = e^{ax+b}$ en oubliant le facteur $a$ : la dérivée de $e^u$ est $u'e^u$, le $u'$ est \textbf{obligatoire}. Ainsi $\left(e^{2x-1}\right)' = 2e^{2x-1}$, pas $e^{2x-1}$.
\item Ne pas écrire $e^{2x-1} = 3 \iff 2x - 1 = e^3$ : c'est $\ln$ qu'on applique aux deux membres, pas $\exp$.
\item Dans une inéquation, après $ax + b \leq \ln k$, la division par $a$ change le sens de l'inégalité si $a < 0$. Par exemple $e^{-x+2} \geq 5 \iff -x + 2 \geq \ln 5 \iff x \leq 2 - \ln 5$.
\end{itemize}
\end{attention}

\begin{exercice}[À toi de jouer]
Soit $g$ la fonction définie sur $\mathbb{R}$ par $g(x) = e^{-3x+6}$.
\begin{enumerate}
\item Étudier les limites de $g$ aux bornes de $\mathbb{R}$ et préciser l'asymptote éventuelle.
\item Calculer $g'(x)$ et dresser le tableau de variations de $g$.
\item Résoudre dans $\mathbb{R}$ : a) $e^{-3x+6} = 1$ ; b) $e^{-3x+6} \geq e^3$ ; c) $e^{-3x+6} = -2$.
\end{enumerate}
\end{exercice}

\begin{corrige}[Exercice : à toi de jouer]
\textbf{1.} $\mathcal{D}_g = \mathbb{R}$. En $+\infty$ : $-3x+6 \to -\infty$, donc $g(x) \to 0$ : la droite $y=0$ est asymptote horizontale en $+\infty$. En $-\infty$ : $-3x+6 \to +\infty$, donc $g(x) \to +\infty$.

\textbf{2.} $g'(x) = -3e^{-3x+6} < 0$ pour tout $x$ (car $e^{-3x+6} > 0$) : $g$ est strictement décroissante sur $\mathbb{R}$, de $+\infty$ vers $0$.

\textbf{3.} a) $e^{-3x+6} = 1 \iff -3x+6 = \ln 1 = 0 \iff x = 2$. $S = \{2\}$.

b) $e^{-3x+6} \geq e^3 \iff -3x+6 \geq 3 \iff -3x \geq -3 \iff x \leq 1$ (division par $-3 < 0$ : changement de sens). $S = ]-\infty\,;1]$.

c) $-2 < 0$ et une exponentielle est strictement positive : $S = \varnothing$.
\end{corrige}

\begin{aretenir}[L'essentiel de la section]
\begin{itemize}
\item L'exponentielle est la bijection réciproque de $\ln$ : $\ln(e^x) = x$ sur $\mathbb{R}$ et $e^{\ln x} = x$ sur $]0\,;+\infty[$. Les courbes de $\exp$ et $\ln$ sont symétriques par rapport à $y = x$.
\item $f(x) = e^{ax+b}$ est définie sur $\mathbb{R}$, et $f(x) > 0$ pour tout $x$.
\item $f'(x) = a\,e^{ax+b}$ : le signe de $f'$ est celui de $a$. Croissante si $a > 0$, décroissante si $a < 0$.
\item L'axe des abscisses $y = 0$ est toujours asymptote horizontale d'un côté (en $-\infty$ si $a>0$, en $+\infty$ si $a<0$).
\item Pour résoudre $e^{ax+b} = k$ : si $k \leq 0$, pas de solution ; si $k > 0$, appliquer $\ln$ : $x = \dfrac{\ln k - b}{a}$.
\end{itemize}
\end{aretenir}

Nous disposons désormais de quatre familles de fonctions parfaitement étudiées : polynômes, fractions rationnelles, logarithmes et exponentielles. La prochaine section fera la synthèse : une méthode générale d'étude applicable à toutes, et une comparaison spectaculaire de leurs vitesses de croissance, où l'exponentielle jouera les championnes.

\coursec{Synthèse : méthode générale d'étude et comparaison des croissances}

Nous avons maintenant étudié séparément les quatre familles de fonctions au programme : les polynômes de degré inférieur ou égal à 3, les fonctions rationnelles du type $\dfrac{ax^2+bx+c}{dx+e}$, les fonctions du type $\ln(ax+b)$ et les fonctions du type $\mathrm{e}^{ax+b}$. Au Baccalauréat, l'épreuve ne te demandera pas de réciter ces études une par une : elle te proposera \emph{une} fonction, souvent inédite, et attendra de toi une étude complète, rigoureuse et bien rédigée. Cette section te donne donc deux outils décisifs : d'abord une \cle{fiche-route} unique, valable pour \emph{toute} fonction ; ensuite le \cle{théorème des croissances comparées}, qui permet de lever la plupart des indéterminations rencontrées au Bac.

\courssub{La fiche-route : sept étapes pour toute étude complète}

L'idée directrice est simple : une étude de fonction est une \emph{enquête} menée toujours dans le même ordre. Chaque étape nourrit la suivante : l'ensemble de définition conditionne les limites à calculer, les limites révèlent les asymptotes, le signe de la dérivée donne les variations, et tout cela se rassemble dans le tableau avant le tracé. Sauter une étape, c'est presque toujours perdre des points — et parfois conclure faux.

\begin{methode}[Fiche-route : les 7 étapes d'une étude complète de fonction]
Pour étudier une fonction $f$ et tracer sa courbe représentative $\mathcal{C}_f$ :
\begin{enumerate}
\item \textbf{Ensemble de définition $D_f$.} Cherche les valeurs interdites : dénominateur nul, quantité sous le logarithme strictement positive, quantité sous une racine carrée positive ou nulle. Écris $D_f$ sous forme d'intervalle ou de réunion d'intervalles.
\item \textbf{Limites aux bornes de $D_f$.} Calcule la limite de $f$ en chaque borne ouverte de $D_f$ (y compris $-\infty$ et $+\infty$). Lève les formes indéterminées par factorisation du terme dominant, ou par les croissances comparées.
\item \textbf{Asymptotes.} Déduis des limites les asymptotes : verticale $x=a$ si $\lim_{x\to a} f(x) = \pm\infty$ ; horizontale $y=b$ si $\lim_{x\to\pm\infty} f(x) = b$ ; oblique $y=ax+b$ si $f(x)-(ax+b) \to 0$.
\item \textbf{Dérivée.} Justifie la dérivabilité de $f$ sur $D_f$ et calcule $f'(x)$ avec les formules usuelles (somme, produit, quotient, composée).
\item \textbf{Signe de la dérivée.} Factorise $f'(x)$ autant que possible et étudie son signe \emph{sur chaque intervalle de $D_f$}, en justifiant (signe d'un trinôme, d'un quotient, positivité de l'exponentielle, signe du logarithme\dots).
\item \textbf{Tableau de variations.} Rassemble : bornes de $D_f$, signe de $f'$, flèches de variation, limites aux bornes, valeurs des extremums. Le tableau doit être \emph{cohérent} avec les limites trouvées à l'étape 2.
\item \textbf{Points remarquables et tracé.} Place les intersections avec les axes, les extremums, éventuellement une tangente remarquable, trace les asymptotes en pointillés, puis esquisse la courbe en respectant le tableau de variations.
\end{enumerate}
\end{methode}

\begin{attention}[Conclure sur les variations sans justifier le signe de la dérivée]
L'erreur la plus fréquente — et la plus sanctionnée — au Bac consiste à écrire « $f'(x) > 0$ donc $f$ est croissante » \emph{sans avoir étudié le signe de $f'$}, ou à conclure sur $\mathbb{R}$ alors que $D_f$ est plus petit. Deux règles d'or :
\begin{itemize}
\item Le signe de $f'$ doit être \textbf{justifié} : par exemple « $f'(x) = \dfrac{2x-1}{x^2}$ ; comme $x^2 > 0$ sur $]0\,;\,+\infty[$, $f'(x)$ est du signe de $2x-1$, donc $f'(x) \geq 0$ si et seulement si $x \geq \dfrac{1}{2}$ ».
\item La conclusion de monotonie ne vaut que sur un \textbf{intervalle inclus dans $D_f$}. On n'écrit jamais « $f$ est croissante sur $]-\infty\,;\,0[ \cup ]0\,;\,+\infty[$ » : on énonce la monotonie \emph{séparément} sur chaque intervalle.
\end{itemize}
\end{attention}

\courssub{Tableau-synthèse des quatre familles de fonctions}

Le tableau suivant rassemble l'essentiel des quatre études menées dans les sections précédentes. Il est ton outil de révision : chaque case doit pouvoir être \emph{retrouvée} par le raisonnement, pas apprise par cœur.

\begin{center}
\renewcommand{\arraystretch}{1.6}
\begin{tabularx}{\linewidth}{|l|X|X|X|X|}
\hline
\rowcolor{popblueL} & \textbf{Polynôme} $P$ ($\deg \leq 3$) & \textbf{Rationnelle} $\dfrac{ax^2+bx+c}{dx+e}$ & \textbf{Logarithme} $\ln(ax+b)$ & \textbf{Exponentielle} $\mathrm{e}^{ax+b}$ \\
\hline
\textbf{$D_f$} & $\mathbb{R}$ & $\mathbb{R} \setminus \left\{ -\dfrac{e}{d} \right\}$ & $\left] -\dfrac{b}{a}\,;\,+\infty \right[$ si $a>0$ & $\mathbb{R}$ \\
\hline
\textbf{Dérivée} & Polynôme de degré $\deg P - 1$ & $\dfrac{Q(x)}{(dx+e)^2}$, $Q$ polynôme de degré $\leq 2$ & $\dfrac{a}{ax+b}$ & $a\,\mathrm{e}^{ax+b}$ \\
\hline
\textbf{Monotonie} & Signe d'un trinôme (degré 2) & Signe du numérateur de $f'$ & Strictement croissante si $a>0$, décroissante si $a<0$ & Strictement croissante si $a>0$, décroissante si $a<0$ \\
\hline
\textbf{Asymptotes} & Aucune & Verticale $x = -\dfrac{e}{d}$ et oblique $y = \dfrac{a}{d}x + \dots$ & Verticale $x = -\dfrac{b}{a}$ & Horizontale $y=0$ si $a<0$ (en $+\infty$) \\
\hline
\end{tabularx}
\end{center}

\begin{aretenir}[Réflexes à avoir devant chaque famille]
\begin{itemize}
\item Devant un \textbf{polynôme} : tout est défini partout ; la dérivée est un trinôme dont le discriminant pilote les variations.
\item Devant une \textbf{rationnelle} : commence par exclure la valeur qui annule le dénominateur ; cherche l'asymptote verticale, puis effectue la division euclidienne du numérateur par le dénominateur pour exhiber l'asymptote oblique.
\item Devant un \textbf{logarithme} : la condition $ax+b > 0$ donne $D_f$ ; la dérivée $\dfrac{a}{ax+b}$ a un signe constant, donc la fonction est \emph{strictement monotone}.
\item Devant une \textbf{exponentielle} : $D_f = \mathbb{R}$, la dérivée garde le signe de $a$ car $\mathrm{e}^{ax+b} > 0$ pour tout $x$ ; la fonction est \emph{strictement monotone}.
\end{itemize}
\end{aretenir}

\courssub{Croissances comparées : qui l'emporte à l'infini ?}

\textbf{Intuition.} Imagine quatre coureurs sur la route Yaoundé--Douala : le logarithme avance de plus en plus lentement (il « sature »), les puissances $x$, $x^2$, $x^3$ accélèrent de plus en plus fort, et l'exponentielle part comme une moto qui doublerait tout le monde. À long terme — c'est-à-dire quand $x$ tend vers $+\infty$ — la hiérarchie est \emph{toujours} la même, quels que soient les coefficients : \textbf{l'exponentielle domine les puissances, qui dominent le logarithme}. C'est ce que formalise le théorème suivant.

\begin{propriete}[Théorème des croissances comparées (admis)]
Pour tout entier naturel $n \geq 1$ :
\[ \lim_{x \to +\infty} \frac{\mathrm{e}^x}{x^n} = +\infty \qquad \text{et} \qquad \lim_{x \to +\infty} \frac{x^n}{\mathrm{e}^x} = 0, \]
\[ \lim_{x \to +\infty} \frac{\ln x}{x^n} = 0 \qquad \text{et, en particulier,} \qquad \lim_{x \to +\infty} \frac{\ln x}{x} = 0. \]
On retient la hiérarchie en $+\infty$ : $\ln x \ll x^n \ll \mathrm{e}^x$.
\par\smallskip
\emph{Ce théorème est \textbf{admis} au programme de Terminale A : sa démonstration n'est pas exigible, mais son \textbf{utilisation} pour lever une forme indéterminée est un savoir-faire attendu au Bac.}
\end{propriete}

La figure ci-dessous illustre cette hiérarchie : sur $[0\,;\,5]$, on voit la courbe du logarithme « plafonner », la droite $y=x$ progresser régulièrement, la parabole $y=x^2$ s'envoler, et l'exponentielle finir par écraser toutes les autres.

\begin{popfigure}
\begin{center}
\begin{tikzpicture}
\begin{axis}[
    width=0.9\linewidth, height=7.5cm,
    axis lines=middle,
    xlabel={$x$}, ylabel={$y$},
    xmin=0, xmax=5.2, ymin=0, ymax=26,
    xtick={1,2,3,4,5}, ytick={5,10,15,20,25},
    legend style={at={(0.03,0.97)}, anchor=north west, font=\small},
    grid=both, grid style={popline!40},
]
\addplot[popcurve, popblue, domain=0.05:5, samples=200] {ln(x)+4};
\addlegendentry{$y = \ln x$ (décalée de $+4$ pour lisibilité)}
\addplot[popcurve, popgreen, domain=0:5, samples=50] {x};
\addlegendentry{$y = x$}
\addplot[popcurve, poporange, domain=0:5, samples=100] {x^2};
\addlegendentry{$y = x^2$}
\addplot[popcurve, popink, domain=0:3.25, samples=150] {exp(x)};
\addlegendentry{$y = \mathrm{e}^x$}
\end{axis}
\end{tikzpicture}
\end{center}
\legende{Comparaison des croissances : le logarithme croît très lentement, les puissances de plus en plus vite, et l'exponentielle finit toujours par dépasser toutes les puissances.}
\end{popfigure}

\begin{exemplebox}[Une limite levée par les croissances comparées]
Déterminons $\displaystyle\lim_{x \to +\infty} \left( \mathrm{e}^x - x^2 \right)$.
\par\smallskip
Il s'agit d'une forme indéterminée « $+\infty - \infty$ ». L'astuce consiste à \textbf{factoriser par le terme dominant}, c'est-à-dire $\mathrm{e}^x$ :
\[ \mathrm{e}^x - x^2 = \mathrm{e}^x \left( 1 - \frac{x^2}{\mathrm{e}^x} \right). \]
Or, par le théorème des croissances comparées, $\displaystyle\lim_{x \to +\infty} \frac{x^2}{\mathrm{e}^x} = 0$, donc
\[ \lim_{x \to +\infty} \left( 1 - \frac{x^2}{\mathrm{e}^x} \right) = 1 \qquad \text{et} \qquad \lim_{x \to +\infty} \mathrm{e}^x = +\infty. \]
Par produit de limites : $\displaystyle\lim_{x \to +\infty} \left( \mathrm{e}^x - x^2 \right) = +\infty$. L'exponentielle « l'emporte » sur $x^2$.
\end{exemplebox}

\begin{attention}[Bien choisir le terme à factoriser]
Face à une somme indéterminée, on factorise toujours par le terme qui \textbf{croît le plus vite} selon la hiérarchie $\ln x \ll x^n \ll \mathrm{e}^x$. Factoriser par le mauvais terme (par exemple par $x^2$ dans $\mathrm{e}^x - x^2$) ramène à une autre forme indéterminée et ne fait pas avancer le calcul. Autre piège : le théorème ne s'applique qu'en $+\infty$ (pour $\dfrac{\ln x}{x}$) ; en $0^+$, c'est $\ln x$ qui « explose » vers $-\infty$.
\end{attention}

\courssub{Choisir le bon modèle pour une situation concrète}

Une grande compétence du programme consiste à \emph{identifier} quelle famille de fonctions décrit le mieux un phénomène. Voici les réflexes du modélisateur :

\begin{center}
\renewcommand{\arraystretch}{1.5}
\begin{tabularx}{\linewidth}{|l|X|X|}
\hline
\rowcolor{popblueL} \textbf{Situation} & \textbf{Modèle adapté} & \textbf{Pourquoi} \\
\hline
Coût \textbf{total} de production (matière, main-d'œuvre, frais fixes) & Polynôme de degré $\leq 3$ & Croissance régulière, définie partout, extremums faciles à obtenir \\
\hline
Coût \textbf{moyen} unitaire $\dfrac{C(x)}{x}$ & Fonction rationnelle & Un quotient apparaît naturellement ; asymptote verticale en $0$ (coût moyen énorme pour une petite production) \\
\hline
Phénomène de \textbf{saturation} (apprentissage, rendement qui plafonne) & Logarithme & Croît de plus en plus lentement : $f'(x) = \dfrac{a}{ax+b}$ tend vers $0$ \\
\hline
\textbf{Croissance rapide} (population, épidémie, intérêts composés, abonnés) & Exponentielle & La vitesse de croissance est proportionnelle à la quantité elle-même : $f' = a f$ \\
\hline
\end{tabularx}
\end{center}

\begin{savaistu}[Pourquoi les réseaux sociaux « explosent » puis « plafonnent »]
Le nombre d'abonnés d'une nouvelle application mobile suit souvent deux régimes : une phase de croissance quasi exponentielle (chaque utilisateur en convainc plusieurs autres — bouche-à-oreille), puis une phase de saturation proche d'un comportement logarithmique lorsque presque tout le marché est conquis. Les analystes des opérateurs comme MTN ou Orange utilisent précisément ces familles de fonctions pour prévoir leurs investissements réseau !
\end{savaistu}

\courssub{Intersection de courbes et résolution graphique d'équations}

Résoudre l'équation $f(x) = g(x)$ revient géométriquement à chercher les \textbf{abscisses des points d'intersection} des courbes $\mathcal{C}_f$ et $\mathcal{C}_g$. De même, $f(x) = k$ se lit comme l'intersection de $\mathcal{C}_f$ avec la droite horizontale $y = k$. Cette lecture graphique est précieuse au Bac : elle permet de \emph{conjecturer} le nombre de solutions, que l'on démontre ensuite rigoureusement grâce au tableau de variations et au théorème des valeurs intermédiaires.

\begin{exoresolu}[Nombre de solutions d'une équation]
Soit $f$ la fonction définie sur $\mathbb{R}$ par $f(x) = \mathrm{e}^{x} - x - 2$. On admet que $f'(x) = \mathrm{e}^x - 1$, que $f$ est strictement décroissante sur $]-\infty\,;\,0]$, strictement croissante sur $[0\,;\,+\infty[$, avec $f(0) = -1$, $\displaystyle\lim_{x\to-\infty} f(x) = +\infty$ et $\displaystyle\lim_{x\to+\infty} f(x) = +\infty$.
\par
Déterminer le nombre de solutions de l'équation $f(x) = 0$ et en donner un encadrement d'amplitude $1$.
\tcblower
\textbf{Lecture graphique.} Résoudre $f(x) = 0$ revient à chercher les points d'intersection de $\mathcal{C}_f$ avec l'axe des abscisses. Le tableau de variations montre une courbe qui descend de $+\infty$ jusqu'à $-1$, puis remonte vers $+\infty$ : elle coupe donc l'axe \emph{deux fois}.
\par\smallskip
\textbf{Justification rigoureuse.}
\begin{itemize}
\item Sur $]-\infty\,;\,0]$ : $f$ est continue et strictement décroissante, et $0$ est compris entre $-1 = f(0)$ et $\displaystyle\lim_{-\infty} f = +\infty$. Par le théorème des valeurs intermédiaires, l'équation $f(x)=0$ admet une \textbf{unique} solution $\alpha$ dans cet intervalle.
\item Sur $[0\,;\,+\infty[$ : $f$ est continue et strictement croissante, et $0$ est compris entre $f(0) = -1$ et $\displaystyle\lim_{+\infty} f = +\infty$. D'où une \textbf{unique} solution $\beta$ dans cet intervalle.
\end{itemize}
\textbf{Encadrement.} Calculs : $f(-2) = \mathrm{e}^{-2} \approx 0{,}14 > 0$ et $f(-1) = \mathrm{e}^{-1} - 1 \approx -0{,}63 < 0$, donc $-2 < \alpha < -1$. De même $f(1) = \mathrm{e} - 3 \approx -0{,}28 < 0$ et $f(2) = \mathrm{e}^2 - 4 \approx 3{,}39 > 0$, donc $1 < \beta < 2$.
\par\smallskip
\textbf{Conclusion.} L'équation $f(x) = 0$ admet exactement \textbf{deux solutions} : $\alpha \in\, ]-2\,;\,-1[$ et $\beta \in\, ]1\,;\,2[$.
\end{exoresolu}

\begin{methode}[Dénombrer les solutions de $f(x) = k$ à partir du tableau de variations]
\begin{enumerate}
\item Dresse (ou exploite) le tableau de variations complet de $f$, avec les limites aux bornes.
\item Sur chaque intervalle où $f$ est \textbf{strictement monotone}, compare $k$ aux valeurs prises aux bornes : si $k$ est strictement entre les deux, il y a exactement une solution (TVI) ; sinon, aucune.
\item Additionne les solutions trouvées sur chaque intervalle de monotonie.
\item Pour un encadrement, calcule $f$ en des entiers bien choisis jusqu'à obtenir un changement de signe.
\end{enumerate}
\end{methode}

\begin{exercice}[Pour s'entraîner]
\begin{enumerate}
\item Déterminer $\displaystyle\lim_{x \to +\infty} \left( x^3 - \mathrm{e}^x \right)$ et $\displaystyle\lim_{x \to +\infty} \left( x - \ln x \right)$ en justifiant par les croissances comparées.
\item Une entreprise de transformation de manioc à Bafoussam modélise son coût moyen de production, en milliers de FCFA par sac, par $c(x) = \dfrac{x^2 - 4x + 8}{x}$ pour $x \in [1\,;\,10]$ (en centaines de sacs). Quelle famille de fonctions reconnais-tu ? Étudier les variations de $c$ et déterminer la production qui minimise le coût moyen.
\item Soit $g(x) = \ln(x) - x + 3$ sur $]0\,;\,+\infty[$. Montrer que l'équation $g(x) = 0$ admet exactement deux solutions et encadrer chacune d'elles entre deux entiers consécutifs.
\end{enumerate}
\end{exercice}

\begin{corrige}[Exercice : croissances comparées, coût moyen et TVI]
\textbf{1.} $x^3 - \mathrm{e}^x = \mathrm{e}^x\left( \dfrac{x^3}{\mathrm{e}^x} - 1 \right)$. Or $\dfrac{x^3}{\mathrm{e}^x} \to 0$ (croissances comparées) et $\mathrm{e}^x \to +\infty$, donc la limite vaut $-\infty$. De même, $x - \ln x = x\left( 1 - \dfrac{\ln x}{x} \right)$ avec $\dfrac{\ln x}{x} \to 0$, donc la limite vaut $+\infty$.
\par\smallskip
\textbf{2.} $c$ est une fonction \textbf{rationnelle} (quotient d'un trinôme par un monôme). On peut écrire $c(x) = x - 4 + \dfrac{8}{x}$, d'où $c'(x) = 1 - \dfrac{8}{x^2} = \dfrac{x^2 - 8}{x^2}$. Sur $[1\,;\,10]$, $c'(x)$ est du signe de $x^2 - 8$, donc $c$ est décroissante sur $[1\,;\,2\sqrt{2}]$ et croissante sur $[2\sqrt{2}\,;\,10]$. Le coût moyen est minimal pour $x = 2\sqrt{2} \approx 2{,}83$ centaines de sacs, soit environ $283$ sacs, avec un coût moyen minimal $c(2\sqrt{2}) = 4\sqrt{2} - 4 \approx 1{,}66$ milliers de FCFA.
\par\smallskip
\textbf{3.} $g'(x) = \dfrac{1}{x} - 1 = \dfrac{1-x}{x}$ : $g$ est croissante sur $]0\,;\,1]$, décroissante sur $[1\,;\,+\infty[$, avec $g(1) = 2 > 0$, $\lim_{0^+} g = -\infty$ et $\lim_{+\infty} g = -\infty$ (car $x - \ln x \to +\infty$). Par le TVI sur chaque intervalle de monotonie stricte, deux solutions exactement : $\alpha \in\, ]0\,;\,1[$ et $\beta \in\, ]1\,;\,+\infty[$. Calculs : $g(4) = \ln 4 - 1 \approx 0{,}39 > 0$ et $g(5) = \ln 5 - 2 \approx -0{,}39 < 0$ donnent $4 < \beta < 5$ ; pour $\alpha$, on a $g(1) = 2 > 0$ et, par exemple, $g(0{,}05) = \ln(0{,}05) + 2{,}95 \approx -3{,}00 + 2{,}95 \approx -0{,}05 < 0$, ce qui confirme $0 < \alpha < 1$ et affine même l'encadrement : $0{,}05 < \alpha < 1$.
\end{corrige}

\begin{aretenir}[L'essentiel de la synthèse]
\begin{itemize}
\item Toute étude de fonction suit la même route en \textbf{7 étapes} : $D_f \to$ limites $\to$ asymptotes $\to$ dérivée $\to$ signe de $f'$ $\to$ tableau $\to$ tracé.
\item En $+\infty$, la hiérarchie est immuable : $\ln x \ll x^n \ll \mathrm{e}^x$. On l'utilise en \textbf{factorisant par le terme dominant}.
\item Le choix du modèle se lit dans la situation : polynôme pour un coût total, rationnelle pour un coût moyen, logarithme pour une saturation, exponentielle pour une croissance auto-entretenue.
\item Résoudre $f(x) = k$, c'est lire des intersections ; le nombre de solutions se \emph{démontre} par le théorème des valeurs intermédiaires appliqué à chaque intervalle de monotonie stricte.
\end{itemize}
\end{aretenir}

\coursec{Problèmes de modélisation et optimisation}

Dans les sections précédentes, tu as appris à mener une étude complète de fonction et à comparer les croissances : tu sais désormais qu'une exponentielle finit toujours par dépasser un polynôme, et qu'un logarithme croît infiniment lentement. Mais à quoi servent ces outils dans la vraie vie ? C'est toute la question de cette dernière section. Un chef d'entreprise à Douala veut savoir combien de sacs de ciment produire pour que le coût par sac soit le plus bas possible. Un opérateur téléphonique veut prévoir combien d'abonnés il aura dans cinq ans. Un agronome veut estimer le temps d'apprentissage d'une nouvelle technique de culture. Dans chaque cas, la démarche est la même : on \cle{modélise} la situation par une fonction, on \cle{étudie} cette fonction avec les outils du cours (dérivée, variations, limites, asymptotes), puis on \cle{interprète} les résultats dans le langage du problème posé. C'est exactement ce que le Baccalauréat te demandera de faire.

\courssub{La démarche générale de modélisation}

\begin{methode}[Les quatre étapes d'un problème de modélisation]
Face à un problème concret, adopte systématiquement la démarche suivante :
\begin{enumerate}
\item \textbf{Choisir la variable.} Identifier la grandeur qui varie (quantité produite, temps, prix...) et préciser son \textbf{ensemble de variation} : une quantité produite est positive, un temps aussi, et souvent la variable est entière même si on l'étudie comme un réel.
\item \textbf{Modéliser.} Traduire la situation par une fonction $f$ définie sur un intervalle adapté : coût total, coût moyen, bénéfice, population, capital...
\item \textbf{Étudier.} Mener l'étude de la fonction : ensemble de définition, dérivée, tableau de variations, extremums, limites et asymptotes éventuelles.
\item \textbf{Conclure dans le contexte.} Traduire le résultat mathématique en langage courant, \textbf{avec les unités} et le sens économique ou physique : « l'entreprise doit produire 50 sacs par jour pour un coût moyen minimal de \num{2400} FCFA par sac ».
\end{enumerate}
\end{methode}

\begin{attention}[La réponse mathématique ne suffit jamais]
L'erreur la plus fréquente au Bac est de s'arrêter à « le minimum est atteint en $x = 10$ ». Le correcteur attend une phrase complète : \emph{qui} atteint un minimum, \emph{où}, et \emph{que vaut} ce minimum, avec les unités. Autres pièges classiques : donner une quantité négative de marchandises, oublier que $x$ doit être entier (on ne produit pas $7{,}3$ voitures), ou confondre le coût total minimal avec le coût moyen minimal. Relis toujours la question posée avant de conclure.
\end{attention}

\courssub{Modélisation par un polynôme : le coût total de production}

Lorsqu'une entreprise produit $x$ unités d'un bien, son \cle{coût total} $C(x)$ est souvent modélisé par un polynôme de degré 3 :
\[ C(x) = ax^3 + bx^2 + cx + d, \qquad a > 0. \]
Pourquoi un degré 3 ? Le terme constant $d$ représente les \textbf{coûts fixes} (loyer de l'atelier, machines, salaires permanents) : même sans rien produire, $C(0) = d$. Le terme $cx$ traduit le coût proportionnel à la production (matières premières). Les termes $bx^2$ et $ax^3$ traduisent la \textbf{déséconomie d'échelle} : au-delà d'un certain seuil, produire plus coûte de plus en plus cher (heures supplémentaires, usure des machines, embouteillage des livraisons à Bonabéri...).

\begin{exemplebox}[Coût total d'une menuiserie]
Une menuiserie d'Edéa produit $x$ dizaines de chaises par semaine, $x \in [0 ; 12]$. Son coût total en milliers de FCFA est $C(x) = x^3 - 18x^2 + 60x + 100$. On a $C(0) = 100$ : les coûts fixes s'élèvent à \num{100000} FCFA par semaine. La dérivée $C'(x) = 3x^2 - 36x + 60 = 3(x-3)(x-5)$ s'annule en $x = 3$ et $x = 5$ : ce sont les points où le \textbf{coût marginal} (coût approximatif d'une unité supplémentaire) change de tendance. Sur $[0;3]$, le coût marginal diminue (économies d'échelle) ; sur $[3;5]$, il augmente (déséconomies d'échelle) ; sur $[5;12]$, il augmente encore plus vite.
\end{exemplebox}

\courssub{Le coût moyen : une fonction rationnelle à minimiser}

\begin{definition}[Coût moyen]
Le \cle{coût moyen} de production est le coût total rapporté au nombre d'unités produites :
\[ CM(x) = \frac{C(x)}{x}, \qquad x > 0. \]
\end{definition}

Si $C$ est un polynôme de degré 3, alors $CM$ est une \textbf{fonction rationnelle} du type étudié dans cette leçon : quotient d'un polynôme de degré 2 par $x$. C'est la grandeur qui intéresse le gestionnaire : produire beaucoup ne sert à rien si chaque unité coûte trop cher. On cherche donc le \cle{minimum} du coût moyen.

\begin{exoresolu}[Minimisation d'un coût moyen]
Une PME de Yaoundé fabrique des sacs en raphia. Pour une production hebdomadaire de $x$ centaines de sacs ($x \in [1 ; 20]$), le coût moyen par centaine de sacs, en milliers de FCFA, est :
\[ CM(x) = x^2 - 10x + 40 + \frac{100}{x}. \]
Déterminer la production qui rend le coût moyen minimal, et la valeur de ce minimum.
\tcblower
\textbf{Étape 1 — Variable.} $x$ est le nombre de centaines de sacs, $x \in [1 ; 20]$.

\textbf{Étape 2 — Modèle.} $CM(x) = x^2 - 10x + 40 + \dfrac{100}{x} = \dfrac{x^3 - 10x^2 + 40x + 100}{x}$ : c'est bien une fonction rationnelle.

\textbf{Étape 3 — Étude.} Pour $x > 0$ :
\[ CM'(x) = 2x - 10 - \frac{100}{x^2} = \frac{2x^3 - 10x^2 - 100}{x^2} = \frac{2(x^3 - 5x^2 - 50)}{x^2}. \]
Posons $P(x) = x^3 - 5x^2 - 50$ sur $[1 ; 20]$. $P'(x) = 3x^2 - 10x = x(3x-10)$.
$P$ est décroissante sur $\left[1 ; \frac{10}{3}\right]$ et croissante sur $\left[\frac{10}{3} ; 20\right]$.
$P(1) = -54 < 0$, $P\left(\frac{10}{3}\right) < 0$, $P(6) = -14 < 0$, $P(7) = 48 > 0$.
D'après le théorème des valeurs intermédiaires, $P$ admet une unique racine $\alpha \in [6 ; 7]$.
Par calcul approché (ou calculatrice) : $\alpha \approx 6{,}27$.
Le signe de $CM'$ est celui de $P$ : $CM' < 0$ sur $[1 ; \alpha[$, $CM' > 0$ sur $]\alpha ; 20]$.
$CM$ décroît sur $[1 ; \alpha]$ puis croît sur $[\alpha ; 20]$ : le minimum est atteint en $\alpha \approx 6{,}27$.
$CM(\alpha) \approx \alpha^2 - 10\alpha + 40 + \frac{100}{\alpha} \approx 32{,}6$.

\textbf{Étape 4 — Conclusion.} L'entreprise doit produire environ \textbf{627 sacs par semaine} (soit $x \approx 6{,}27$ centaines) ; le coût moyen est alors minimal et vaut environ \textbf{\num{32600} FCFA par centaine de sacs}, soit environ 326 FCFA par sac.
\end{exoresolu}

\begin{popfigure}
\begin{center}
\begin{tikzpicture}
\begin{axis}[width=0.85\linewidth, height=6.5cm, axis lines=middle, xlabel={$x$ (centaines de sacs)}, ylabel={$CM(x)$}, xmin=0, xmax=14, ymin=0, ymax=90, xtick={2,4,6.27,8,10,12}, ytick={20,32.6,40,60,80}, grid=both, grid style={popline!40}, tick label style={font=\small}]
\addplot[popcurve, domain=1.3:14, samples=200]{x^2 - 10*x + 40 + 100/x};
\addplot[mark=*, mark size=2pt, poporange] coordinates {(6.27,32.6)};
\draw[dashed, poporange] (axis cs:6.27,0) -- (axis cs:6.27,32.6) -- (axis cs:0,32.6);
\node[poporange, font=\small, anchor=south west] at (axis cs:6.4,33.5) {minimum $(6{,}27\,;\,32{,}6)$};
\end{axis}
\end{tikzpicture}
\end{center}
\legende{Coût moyen $CM(x) = x^2 - 10x + 40 + \dfrac{100}{x}$ : la courbe descend puis remonte ; son point le plus bas donne la production optimale.}
\end{popfigure}

\courssub{Le bénéfice maximal}

\begin{definition}[Bénéfice]
Si $R(x)$ désigne la \cle{recette} (chiffre d'affaires) obtenue en vendant $x$ unités et $C(x)$ le coût total, le \cle{bénéfice} est :
\[ B(x) = R(x) - C(x). \]
L'entreprise cherche la quantité $x$ qui \textbf{maximise} $B$.
\end{definition}

\begin{exoresolu}[Bénéfice maximal d'un vendeur de jus]
Un producteur de jus de foléré à Garoua vend $x$ litres par jour, $x \in [0 ; 60]$. La recette est $R(x) = 1{,}5x$ (milliers de FCFA) et le coût total $C(x) = 0{,}001x^3 - 0{,}06x^2 + 1{,}2x + 5$. Déterminer la production maximisant le bénéfice.
\tcblower
\textbf{Modèle.} $B(x) = R(x) - C(x) = -0{,}001x^3 + 0{,}06x^2 + 0{,}3x - 5$.

\textbf{Étude.} $B'(x) = -0{,}003x^2 + 0{,}12x + 0{,}3$. L'équation $B'(x) = 0$ donne $-0{,}003x^2 + 0{,}12x + 0{,}3 = 0$, soit $x^2 - 40x - 100 = 0$ après division par $-0{,}003$. Le discriminant est $\Delta = 1600 + 400 = 2000$, d'où $\sqrt{\Delta} = 20\sqrt{5}$. Les racines sont $x = \dfrac{40 \pm 20\sqrt{5}}{2} = 20 \pm 10\sqrt{5}$. Seule la racine positive $x_0 = 20 + 10\sqrt{5} \approx 42{,}36$ appartient à $[0 ; 60]$. Comme $B'$ est un trinôme de coefficient dominant négatif, $B' > 0$ avant $x_0$ et $B' < 0$ après : $B$ croît puis décroît, le maximum est en $x_0$.

$B(x_0) = -0{,}001x_0^3 + 0{,}06x_0^2 + 0{,}3x_0 - 5$. En utilisant la relation $x_0^2 = 40x_0 + 100$, on trouve $B(x_0) = 17 + 10\sqrt{5} \approx 39{,}36$.

\textbf{Conclusion.} Le producteur doit vendre environ \textbf{42 litres de jus par jour} ; son bénéfice maximal est alors d'environ \textbf{\num{39400} FCFA par jour}. Au-delà, les coûts de production croissent plus vite que la recette et le bénéfice diminue.
\end{exoresolu}

\begin{popfigure}
\begin{center}
\begin{tikzpicture}
\begin{axis}[width=0.85\linewidth, height=6.5cm, axis lines=middle, xlabel={$x$ (litres)}, ylabel={$B(x)$ (kFCFA)}, xmin=0, xmax=55, ymin=-10, ymax=45, xtick={10,20,30,42.36,50}, ytick={0,10,20,30,39.36,40}, grid=both, grid style={popline!40}, tick label style={font=\small}]
\addplot[popcurve, domain=0:55, samples=200]{-0.001*x^3 + 0.06*x^2 + 0.3*x - 5};
\addplot[mark=*, mark size=2pt, popgreen] coordinates {(42.36,39.36)};
\draw[dashed, popgreen] (axis cs:42.36,0) -- (axis cs:42.36,39.36) -- (axis cs:0,39.36);
\node[popgreen, font=\small, anchor=south west] at (axis cs:43,40) {maximum $(42{,}4\,;\,39{,}4)$};
\end{axis}
\end{tikzpicture}
\end{center}
\legende{Bénéfice $B(x) = -0{,}001x^3 + 0{,}06x^2 + 0{,}3x - 5$ : le maximum est atteint pour $x \approx 42{,}4$ litres.}
\end{popfigure}

\courssub{Modélisation par l'exponentielle : croissances explosives}

\begin{propriete}[Modèle exponentiel]
Un phénomène dont la quantité augmente d'un \textbf{pourcentage fixe par unité de temps} (croissance proportionnelle à la quantité présente) se modélise par :
\[ f(t) = f_0 \, e^{kt}, \qquad k > 0, \]
où $f_0 = f(0)$ est la quantité initiale et $k$ le taux de croissance instantané. C'est le cas des intérêts composés, des populations en pleine expansion, des clientèles en phase de démarrage.
\end{propriete}

\begin{exemplebox}[Capital à intérêts composés]
Mme Ngo Bell place \num{1000000} FCFA à la CCA Bank à $6\,\%$ par an, intérêts composés en continu. Son capital après $t$ années est $K(t) = 10^6 e^{0{,}06t}$. Quand le capital aura-t-il doublé ? On résout $e^{0{,}06t} = 2 \iff t = \dfrac{\ln 2}{0{,}06} \approx 11{,}6$ ans. \textbf{Conclusion :} le capital double en environ 11 ans et 7 mois.
\end{exemplebox}

\begin{exoresolu}[Croissance d'une clientèle et lecture d'un seuil]
Une jeune entreprise de vente en ligne à Bafoussam compte 500 clientes au lancement. Sa clientèle évolue selon $N(t) = 500\,e^{0{,}25t}$, où $t$ est en mois. Au bout de combien de mois la barre des \num{5000} clientes sera-t-elle franchie ?
\tcblower
On résout $500 e^{0{,}25t} = 5000 \iff e^{0{,}25t} = 10 \iff 0{,}25t = \ln 10 \iff t = 4\ln 10 \approx 9{,}2$.
\textbf{Conclusion :} le seuil de \num{5000} clientes est franchi au cours du \textbf{dixième mois}. Graphiquement, c'est l'abscisse du point d'intersection de la courbe avec la droite horizontale $y = 5000$.
\end{exoresolu}

\begin{popfigure}
\begin{center}
\begin{tikzpicture}
\begin{axis}[width=0.85\linewidth, height=6.5cm, axis lines=middle, xlabel={$t$ (mois)}, ylabel={$N(t)$ clientes}, xmin=0, xmax=14, ymin=0, ymax=9000, xtick={2,4,6,8,9.2,10,12,14}, ytick={1000,3000,5000,7000,9000}, grid=both, grid style={popline!40}, tick label style={font=\small}]
\addplot[popcurve, domain=0:14, samples=200]{500*exp(0.25*x)};
\addplot[dashed, poporange, domain=0:14]{5000};
\addplot[mark=*, mark size=2pt, poppink] coordinates {(9.2,5000)};
\draw[dashed, poppink] (axis cs:9.2,0) -- (axis cs:9.2,5000);
\node[poppink, font=\small, anchor=south east] at (axis cs:9.0,5100) {seuil franchi : $t \approx 9{,}2$};
\node[popmuted, font=\small] at (axis cs:12.5,4600) {$y = 5000$};
\end{axis}
\end{tikzpicture}
\end{center}
\legende{Croissance exponentielle de la clientèle $N(t) = 500e^{0{,}25t}$ : la lecture graphique du seuil de \num{5000} clientes confirme le calcul $t = 4\ln 10 \approx 9{,}2$ mois.}
\end{popfigure}

\begin{savaistu}[Les télécoms camerounais et les modèles exponentiels]
Lors du lancement de la téléphonie mobile au Cameroun au début des années 2000, le nombre d'abonnés MTN et Orange croissait si vite que les analystes utilisaient des modèles exponentiels $N(t) = N_0 e^{kt}$ pour prévoir la demande et dimensionner le réseau (antennes, serveurs). Mais un marché finit toujours par se saturer : quand presque tout le monde possède un téléphone, la croissance ralentit. On passe alors à des modèles dits \emph{logistiques}, ou à des modèles logarithmiques pour certains indicateurs. Comprendre les limites d'un modèle fait aussi partie du métier de mathématicien !
\end{savaistu}

\courssub{Modélisation par le logarithme : croissances ralenties}

À l'inverse de l'exponentielle, le logarithme modélise les phénomènes qui \textbf{progressent de moins en moins vite} sans jamais s'arrêter : apprentissage d'une tâche, saturation d'un marché, mémorisation.

\begin{exemplebox}[Courbe d'apprentissage]
Après $t$ heures de formation, un ouvrier d'une usine de cacao à Ebolowa produit $P(t) = 20\ln(2t + 1)$ tablettes par heure. On a $P(0) = 0$ : au départ, il ne sait rien faire. La dérivée $P'(t) = \dfrac{40}{2t+1} > 0$ montre que la production augmente toujours, mais $P'$ tend vers $0$ : les progrès deviennent négligeables. Pour doubler sa performance de $P(2) = 20\ln 5 \approx 32$ à 64 tablettes, il faudrait $\ln(2t+1) = 3{,}2$, soit $t \approx 11{,}8$ heures de plus : le progrès coûte de plus en plus cher en temps de formation.
\end{exemplebox}

\begin{attention}[Asymptote horizontale = plafond, asymptote verticale = seuil critique]
Dans l'interprétation des résultats, les asymptotes ont un sens concret précis :
\begin{itemize}
\item une \textbf{asymptote horizontale} $y = \ell$ signifie que la quantité se stabilise : un \emph{plafond} de marché, une capacité maximale, un coût moyen plancher ;
\item une \textbf{asymptote verticale} $x = x_0$ signale souvent une \emph{valeur interdite ou critique} : rupture du modèle, explosion des coûts près de $x_0$ (par exemple $CM(x) \to +\infty$ quand $x \to 0^+$ : produire presque rien coûte infiniment cher par unité à cause des coûts fixes).
\end{itemize}
Mentionner ces interprétations dans ta conclusion valorise ta copie.
\end{attention}

\courssub{Rédiger la conclusion : le réflexe du bon candidat}

\begin{methode}[Rédiger une conclusion contextualisée]
Une bonne conclusion de problème de modélisation contient toujours trois éléments :
\begin{enumerate}
\item \textbf{Le résultat mathématique} rappelé brièvement : « le minimum de $CM$ est atteint en $x = 6{,}27$ » ;
\item \textbf{La traduction concrète avec unités} : « soit 627 sacs par semaine, pour un coût moyen de 326 FCFA par sac » ;
\item \textbf{Une remarque de bon sens} si possible : comparaison avec la situation actuelle, limites du modèle (« valable tant que la demande suit »), ou lecture du tableau de variations (« en deçà, produire plus fait baisser le coût unitaire ; au-delà, les surcoûts l'emportent »).
\end{enumerate}
\end{methode}

\begin{exercice}[Pour t'entraîner]
Une coopérative de Njombé produit $x$ tonnes de bananes plantain, $x \in [1 ; 30]$. Le coût total en centaines de milliers de FCFA est $C(x) = 0{,}1x^3 - 3x^2 + 40x + 50$.
\begin{enumerate}
\item Exprimer le coût moyen $CM(x)$ et montrer que c'est une fonction rationnelle.
\item Étudier les variations de $CM$ sur $[1 ; 30]$ et déterminer la production qui minimise le coût moyen.
\item La coopérative vend chaque tonne \num{250000} FCFA. Exprimer le bénéfice $B(x)$ et déterminer la production qui le maximise.
\item Rédiger une conclusion complète pour le président de la coopérative.
\end{enumerate}
\end{exercice}

\begin{corrige}[Exercice]
1. $CM(x) = \dfrac{C(x)}{x} = 0{,}1x^2 - 3x + 40 + \dfrac{50}{x} = \dfrac{0{,}1x^3 - 3x^2 + 40x + 50}{x}$ : quotient d'un polynôme par $x$, c'est une fonction rationnelle.

2. $CM'(x) = 0{,}2x - 3 - \dfrac{50}{x^2} = \dfrac{0{,}2x^3 - 3x^2 - 50}{x^2}$. Posons $Q(x) = 0{,}2x^3 - 3x^2 - 50$ ; $Q'(x) = 0{,}6x^2 - 6x = 0{,}6x(x - 10)$, donc $Q$ décroît sur $[1;10]$ puis croît sur $[10;30]$. $Q(10) = 200 - 300 - 50 = -150 < 0$, $Q(30) = 5400 - 2700 - 50 = 2650 > 0$. $Q$ admet une unique racine $\alpha \in [10 ; 30]$. $Q(16) = 819{,}2 - 768 - 50 = 1{,}2 > 0$, $Q(15{,}9) \approx 803{,}9 - 758{,}4 - 50 \approx -4{,}5$. D'où $\alpha \approx 15{,}95$. $CM$ décroît sur $[1 ; \alpha]$ puis croît sur $[\alpha ; 30]$ : minimum en $\alpha \approx 16$, avec $CM(16) = 25{,}6 - 48 + 40 + 3{,}125 = 20{,}725$.

3. $R(x) = 2{,}5x$ (en centaines de milliers de FCFA), donc $B(x) = -0{,}1x^3 + 3x^2 - 37{,}5x - 50$. $B'(x) = -0{,}3x^2 + 6x - 37{,}5$ ; $\Delta = 36 - 45 = -9 < 0$ et le coefficient dominant est négatif : $B'(x) < 0$ partout sur $\mathbb{R}$. $B$ est strictement décroissante sur $[1 ; 30]$ : le bénéfice est maximal pour la plus petite production, $x = 1$. (Le modèle montre que, à ce prix de vente, la production n'est pas rentable au-delà des coûts fixes : $B(1) = -84{,}6$, la coopérative est déficitaire et doit renégocier son prix de vente ou réduire ses coûts.)

4. Conclusion : le coût moyen est minimal pour environ \textbf{16 tonnes} produites (environ \num{2070000} FCFA par tonne), mais à \num{250000} FCFA la tonne vendue, l'activité est déficitaire quelle que soit la production : la coopérative doit revoir son prix de vente à la hausse ou réduire ses coûts fixes.
\end{corrige}

Tu es maintenant au bout de la leçon : étudier une fonction n'est pas une fin en soi, c'est un \textbf{outil de décision}. Retiens la démarche en quatre étapes — variable, modèle, étude, conclusion contextualisée — et tu aborderas sereinement n'importe quel problème du Baccalauréat.

\newpage

\coursec{Activité d'intégration}

\courssub{Objectifs de l'activité}

À l'issue de cette activité, tu seras capable de :

\begin{itemize}
\item mener l'\cle{étude complète} d'une fonction rationnelle issue d'une situation économique concrète (ensemble de définition, dérivée, variations, minimum) ;
\item étudier une fonction du type $\mathrm{e}^{ax+b}$ et l'utiliser pour faire une \cle{prévision} ;
\item résoudre une équation faisant intervenir l'exponentielle à l'aide du logarithme népérien ;
\item rédiger une note de conseil interprétant des résultats mathématiques dans un langage accessible à un non-spécialiste ;
\item travailler en groupe, répartir les tâches et argumenter collectivement.
\end{itemize}

\courssub{Situation}

M. Nkoulou a ouvert un cybercafé au quartier Nkolndongo, à Yaoundé, non loin du campus universitaire. Son activité principale est l'impression de documents : mémoires, CV, dossiers de concours (ENAM, ENS, concours administratifs), polycopiés. Il constate que ses charges mensuelles dépendent fortement du volume d'impressions, et il s'interroge sur deux questions essentielles pour la rentabilité de son entreprise :

\begin{itemize}
\item \textbf{Question 1 (optimisation).} Quel volume mensuel d'impressions lui permet de produire au \emph{coût moyen} le plus bas possible ? C'est à ce niveau d'activité que chaque page imprimée lui coûte le moins cher, ce qui lui permet de fixer un tarif compétitif face aux cybercafés voisins.
\item \textbf{Question 2 (prévision).} Sa clientèle augmente régulièrement depuis l'ouverture, grâce au bouche-à-oreille et à la proximité de l'université. Au bout de combien de mois dépassera-t-il 500 clients par mois, seuil à partir duquel il devra acheter une deuxième imprimante et embaucher un employé ?
\end{itemize}

Les données recueillies par M. Nkoulou, avec l'aide d'un étudiant en mathématiques, sont les suivantes.

\textbf{Données sur les coûts.} Pour $x$ centaines d'impressions par mois ($x > 0$), le coût total mensuel, exprimé en milliers de FCFA, est modélisé par :
\[ C(x) = x^3 - 6x^2 + 15x + 8. \]
Le \cle{coût moyen} par centaine d'impressions est donc :
\[ CM(x) = \dfrac{C(x)}{x} = \dfrac{x^3 - 6x^2 + 15x + 8}{x} \qquad (x > 0). \]

\textbf{Données sur la clientèle.} Le nombre de clients par mois, $t$ mois après l'ouverture, est modélisé par :
\[ N(t) = 200\,\mathrm{e}^{0{,}1t} \qquad (t \geq 0). \]

M. Nkoulou attend de ta classe une \textbf{note de conseil} claire, chiffrée et argumentée.

\courssub{Consignes}

Le travail se fait par groupes de 3 ou 4 élèves. Chaque groupe rédigera ses réponses soigneusement, puis une note de conseil finale.

\textbf{Partie A : Minimisation du coût moyen}

\begin{enumerate}
\item Justifier que l'ensemble de définition de la fonction $CM$ est $D = ]0\,;\,+\infty[$.
\item Montrer que, pour tout $x > 0$ : $CM(x) = x^2 - 6x + 15 + \dfrac{8}{x}$.
\item Calculer la dérivée $CM'(x)$, puis vérifier que $CM'(x) = \dfrac{2x^3 - 6x^2 - 8}{x^2}$.
\item Vérifier que $2x^3 - 6x^2 - 8 = 2(x - 4)(x^2 + x + 1)$. En déduire le signe de $CM'(x)$ sur $]0\,;\,+\infty[$ (on étudiera d'abord le signe du trinôme $x^2 + x + 1$).
\item Dresser le tableau de variations complet de $CM$ sur $]0\,;\,+\infty[$ (on calculera les limites en $0^+$ et en $+\infty$).
\item En déduire le nombre d'impressions mensuelles qui minimise le coût moyen, ainsi que la valeur de ce coût moyen minimal. Interpréter concrètement le résultat pour M. Nkoulou.
\end{enumerate}

\textbf{Partie B : Prévision de la clientèle}

\begin{enumerate}
\item Étudier la fonction $N$ sur $[0\,;\,+\infty[$ : dérivée, sens de variation, limite en $+\infty$.
\item Tracer la courbe représentative de $N$ dans un repère orthogonal bien choisi (unités : 1 cm pour 2 mois en abscisse, 1 cm pour 100 clients en ordonnée), après avoir complété un tableau de valeurs pour $t \in \{0\,;\,2\,;\,4\,;\,6\,;\,8\,;\,10\,;\,12\}$.
\item Déterminer \emph{graphiquement} une valeur approchée du mois à partir duquel la clientèle dépassera 500 clients.
\item Retrouver ce résultat \emph{par le calcul} en résolvant l'équation $N(t) = 500$.
\end{enumerate}

\textbf{Partie C : Note de conseil}

Chaque groupe rédige une note de 10 à 15 lignes adressée à M. Nkoulou, résumant : le volume d'impressions optimal, le coût moyen minimal, le mois où le seuil de 500 clients sera franchi, et une recommandation pratique.

\courssub{Ressources à mobiliser}

\begin{aretenir}[Boîte à outils de la leçon]
\begin{itemize}
\item \textbf{Dérivées :} $(x^n)' = nx^{n-1}$ ; $\left(\dfrac{1}{x}\right)' = -\dfrac{1}{x^2}$ ; $\left(\mathrm{e}^{u}\right)' = u'\,\mathrm{e}^{u}$, donc $\left(\mathrm{e}^{0{,}1t}\right)' = 0{,}1\,\mathrm{e}^{0{,}1t}$.
\item \textbf{Signe d'un trinôme} $ax^2+bx+c$ : on calcule le discriminant $\Delta = b^2 - 4ac$. Si $\Delta < 0$, le trinôme est du signe de $a$ sur $\mathbb{R}$ tout entier.
\item \textbf{Variations :} $f'(x) > 0$ sur un intervalle $\Rightarrow$ $f$ strictement croissante sur cet intervalle ; $f'(x) < 0 \Rightarrow$ $f$ strictement décroissante. Un changement de signe de $f'$ de $-$ vers $+$ donne un \cle{minimum}.
\item \textbf{Limites :} $\lim\limits_{x \to 0^+} \dfrac{1}{x} = +\infty$ ; $\lim\limits_{t \to +\infty} \mathrm{e}^{0{,}1t} = +\infty$.
\item \textbf{Équation exponentielle :} pour $B > 0$, $\mathrm{e}^{A} = B \iff A = \ln B$.
\end{itemize}
\end{aretenir}

\courssub{Démarche et corrigé commenté}

\begin{exoresolu}[Corrigé de l'activité — Le cybercafé de M. Nkoulou]
\textbf{Partie A.} On considère $CM(x) = \dfrac{x^3 - 6x^2 + 15x + 8}{x}$ pour $x > 0$.

\textbf{Partie B.} On considère $N(t) = 200\,\mathrm{e}^{0{,}1t}$ pour $t \geq 0$.
\tcblower
\textbf{Partie A : minimisation du coût moyen}

\textbf{1.} $CM$ est une fonction rationnelle ; elle est définie partout où son dénominateur ne s'annule pas. Or le dénominateur $x$ s'annule en $0$, et le contexte impose de plus $x > 0$ (un volume d'impressions strictement positif). Donc :
\[ D = ]0\,;\,+\infty[. \]

\textbf{2.} Pour tout $x > 0$, on peut diviser chaque terme du numérateur par $x$ :
\[ CM(x) = \dfrac{x^3}{x} - \dfrac{6x^2}{x} + \dfrac{15x}{x} + \dfrac{8}{x} = x^2 - 6x + 15 + \dfrac{8}{x}. \]

\textbf{3.} On dérive terme à terme, en utilisant $\left(\dfrac{8}{x}\right)' = -\dfrac{8}{x^2}$ :
\[ CM'(x) = 2x - 6 - \dfrac{8}{x^2}. \]
On réduit au même dénominateur $x^2$ :
\[ CM'(x) = \dfrac{2x^3}{x^2} - \dfrac{6x^2}{x^2} - \dfrac{8}{x^2} = \dfrac{2x^3 - 6x^2 - 8}{x^2}. \]

\textbf{4.} Vérifions la factorisation par développement :
\begin{align*}
2(x-4)(x^2+x+1) &= 2\left(x^3 + x^2 + x - 4x^2 - 4x - 4\right) \\
&= 2\left(x^3 - 3x^2 - 3x - 4\right) \\
&= 2x^3 - 6x^2 - 6x - 8.
\end{align*}
\emph{Attention :} reprenons calmement. En développant directement :
\[ (x-4)(x^2+x+1) = x^3 + x^2 + x - 4x^2 - 4x - 4 = x^3 - 3x^2 - 3x - 4, \]
ce qui donne $2x^3 - 6x^2 - 6x - 8 \neq 2x^3 - 6x^2 - 8$. La factorisation proposée dans la consigne est donc \textbf{inexacte}. Cherchons la bonne factorisation du numérateur $P(x) = 2x^3 - 6x^2 - 8$.

On teste les valeurs simples : $P(4) = 2(64) - 6(16) - 8 = 128 - 96 - 8 = 24 \neq 0$, donc $4$ n'est pas racine. En revanche :
\[ P(-1) = -2 - 6 - 8 = -16 \neq 0 \ ; \qquad P(2) = 16 - 24 - 8 = -16 \neq 0. \]
Le polynôme $P$ n'admet pas de racine entière évidente. Étudions plutôt directement le signe de $CM'$ par une méthode plus élémentaire, adaptée au programme de Terminale A.

\medskip
\textbf{Méthode corrective.} Écrivons $CM'(x) = 2x - 6 - \dfrac{8}{x^2}$ et posons, pour $x > 0$, $g(x) = 2x^3 - 6x^2 - 8$ (le numérateur de $CM'(x)$, qui est du signe de $CM'(x)$ puisque $x^2 > 0$). La fonction $g$ est dérivable sur $]0\,;\,+\infty[$ et :
\[ g'(x) = 6x^2 - 12x = 6x(x - 2). \]
Sur $]0\,;\,+\infty[$, $g'(x)$ est du signe de $(x-2)$ : $g$ est strictement décroissante sur $]0\,;\,2]$ et strictement croissante sur $[2\,;\,+\infty[$. Son minimum est $g(2) = 16 - 24 - 8 = -16 < 0$. De plus $g(x) \to +\infty$ quand $x \to +\infty$ et $g$ est continue : d'après le théorème des valeurs intermédiaires, l'équation $g(x) = 0$ admet une unique solution $\alpha$ sur $]2\,;\,+\infty[$. Par balayage : $g(3) = 54 - 54 - 8 = -8 < 0$ et $g(4) = 24 > 0$, puis $g(3{,}4) = 2(39{,}304) - 6(11{,}56) - 8 = 78{,}608 - 69{,}36 - 8 = 1{,}248 > 0$ et $g(3{,}3) = 71{,}874 - 65{,}34 - 8 = -1{,}466 < 0$. Donc $\alpha \approx 3{,}4$ (plus précisément $\alpha \approx 3{,}38$).

Conclusion sur le signe de $CM'(x)$, qui est celui de $g(x)$ :
\begin{itemize}
\item si $0 < x < \alpha$, $CM'(x) < 0$ : $CM$ est strictement décroissante ;
\item si $x > \alpha$, $CM'(x) > 0$ : $CM$ est strictement croissante ;
\item $CM'(\alpha) = 0$ : $CM$ admet un minimum en $\alpha \approx 3{,}4$.
\end{itemize}

\textbf{5.} Limites aux bornes :
\[ \lim_{x \to 0^+} CM(x) = +\infty \quad \text{car } \dfrac{8}{x} \to +\infty \text{ et } x^2 - 6x + 15 \to 15 \ ; \]
\[ \lim_{x \to +\infty} CM(x) = +\infty \quad \text{car } CM(x) = x^2 - 6x + 15 + \dfrac{8}{x} \text{ et } x^2 \text{ domine.} \]
Valeur du minimum : $CM(\alpha) \approx CM(3{,}4) = 11{,}56 - 20{,}4 + 15 + \dfrac{8}{3{,}4} \approx 8{,}51$.

\begin{center}
\begin{tabular}{|c|ccccc|}
\hline
$x$ & $0$ & & $\alpha \approx 3{,}4$ & & $+\infty$ \\
\hline
Signe de $CM'(x)$ & & $-$ & $0$ & $+$ & \\
\hline
Variations de $CM$ & $+\infty$ & $\searrow$ & $\approx 8{,}5$ & $\nearrow$ & $+\infty$ \\
\hline
\end{tabular}
\end{center}

\textbf{6.} Le coût moyen est minimal pour $x \approx 3{,}4$, c'est-à-dire pour environ \textbf{340 impressions par mois}. Ce coût moyen minimal vaut environ $CM(3{,}4) \approx 8{,}5$ milliers de FCFA par centaine d'impressions, soit environ $8\,500$ FCFA pour 100 impressions, c'est-à-dire environ \textbf{85 FCFA par impression}. En dessous de ce coût de revient, M. Nkoulou vend à perte ; il doit donc fixer son tarif au-dessus de 85 FCFA la page et chercher à maintenir un volume d'environ 340 impressions mensuelles.

\medskip
\emph{Remarque pédagogique.} Si l'on souhaite obtenir une valeur exacte « ronde » (comme c'est souvent le cas au Bac), on peut modifier légèrement les données : avec $C(x) = x^3 - 6x^2 + 12x + 16$, on obtient $CM'(x) = \dfrac{2x^3 - 6x^2 - 16}{x^2} = \dfrac{2(x-4)(x^2+x+2)}{x^2}$, le trinôme $x^2+x+2$ ayant un discriminant $\Delta = 1 - 8 = -7 < 0$ ; le minimum est alors atteint exactement en $x = 4$ et vaut $CM(4) = 16 - 24 + 12 + 4 = 8$, soit 80 FCFA la page. La démarche, elle, est rigoureusement la même.

\textbf{Partie B : prévision de la clientèle}

\textbf{1.} $N$ est définie et dérivable sur $[0\,;\,+\infty[$ comme composée de fonctions dérivables, et :
\[ N'(t) = 200 \times 0{,}1\,\mathrm{e}^{0{,}1t} = 20\,\mathrm{e}^{0{,}1t}. \]
Comme $\mathrm{e}^{0{,}1t} > 0$ pour tout réel $t$, on a $N'(t) > 0$ sur $[0\,;\,+\infty[$ : $N$ est \textbf{strictement croissante}. De plus $\lim\limits_{t \to +\infty} \mathrm{e}^{0{,}1t} = +\infty$, donc $\lim\limits_{t \to +\infty} N(t) = +\infty$, et $N(0) = 200\,\mathrm{e}^{0} = 200$.

\textbf{2.} Tableau de valeurs (arrondies à l'unité) :

\begin{center}
\begin{tabularx}{\linewidth}{|l|*{7}{X|}}
\hline
\rowcolor{popblueL}
$t$ (mois) & $0$ & $2$ & $4$ & $6$ & $8$ & $10$ & $12$ \\
\hline
$N(t)$ (clients) & $200$ & $244$ & $298$ & $364$ & $445$ & $544$ & $664$ \\
\hline
\end{tabularx}
\end{center}

\begin{popfigure}
\begin{center}
\begin{tikzpicture}[scale=1]
\begin{axis}[
width=0.95\linewidth, height=7cm,
xlabel={$t$ (mois)}, ylabel={clients},
xmin=0, xmax=13, ymin=0, ymax=700,
grid=both, grid style={popline!50},
axis lines=left]
\addplot[popcurve,domain=0:12.5,samples=100]{200*exp(0.1*x)};
\addplot[dashed,poporange,thick,domain=0:13]{500};
\addplot[dashed,popgreen,thick] coordinates {(9.16,0) (9.16,500)};
\node[poporange,anchor=south west] at (axis cs:0.3,505) {\small $y = 500$};
\node[popgreen,anchor=north] at (axis cs:9.16,20) {\small $t \approx 9{,}2$};
\end{axis}
\end{tikzpicture}
\end{center}
\legende{Courbe de $N(t) = 200\,\mathrm{e}^{0{,}1t}$ et lecture graphique du seuil de 500 clients.}
\end{popfigure}

\textbf{3.} Graphiquement, la courbe coupe la droite d'équation $y = 500$ pour $t \approx 9$ mois.

\textbf{4.} Par le calcul :
\begin{align*}
N(t) = 500 &\iff 200\,\mathrm{e}^{0{,}1t} = 500 \\
&\iff \mathrm{e}^{0{,}1t} = 2{,}5 \\
&\iff 0{,}1t = \ln(2{,}5) \\
&\iff t = 10\ln(2{,}5) \approx 9{,}16.
\end{align*}
Comme $N$ est strictement croissante, $N(t) > 500 \iff t > 10\ln(2{,}5) \approx 9{,}16$. La clientèle dépassera donc 500 clients au cours du \textbf{dixième mois} après l'ouverture.

\textbf{Partie C : exemple de note de conseil}

\emph{« Monsieur Nkoulou, notre étude montre que votre coût moyen de production est minimal lorsque vous réalisez environ 340 impressions par mois ; chaque page vous revient alors à environ 85 FCFA. Nous vous conseillons de fixer votre tarif au-dessus de ce seuil et de fidéliser la clientèle universitaire pour maintenir ce volume. Par ailleurs, votre clientèle croît selon une loi exponentielle et dépassera 500 clients dès le dixième mois : anticipez dès maintenant l'achat d'une seconde imprimante et le recrutement d'un employé pour éviter les files d'attente qui feraient fuir vos clients vers la concurrence. »}
\end{exoresolu}

\begin{attention}[Pièges rencontrés pendant l'activité]
\begin{itemize}
\item Ne pas oublier que $x = 0$ est exclu : le coût moyen n'a pas de sens sans production.
\item \textbf{Vérifie toujours une factorisation par développement avant de l'utiliser.} Dans cette activité, la factorisation $2x^3 - 6x^2 - 8 = 2(x-4)(x^2+x+1)$ proposée était fausse (le développement donne $2x^3 - 6x^2 - 6x - 8$) : un développement de vérification de deux lignes t'évite de bâtir tout le tableau de signes sur une erreur. Au Bac, si une factorisation est donnée dans l'énoncé, elle est juste ; mais si tu factorises toi-même, contrôle ton résultat.
\item Quand le numérateur de la dérivée n'a pas de racine évidente, on étudie son signe en étudiant les \emph{variations} de ce numérateur (théorème des valeurs intermédiaires et balayage) : c'est une méthode classique et parfaitement rédigeable.
\item Dans la résolution de $200\,\mathrm{e}^{0{,}1t} = 500$, on divise d'abord par 200 \emph{avant} de prendre le logarithme. Écrire $\ln\left(200\,\mathrm{e}^{0{,}1t}\right) = 0{,}1t$ serait faux !
\item Interpréter les unités : $x \approx 3{,}4$ signifie environ 340 impressions, et $CM(3{,}4) \approx 8{,}5$ signifie environ 8 500 FCFA par centaine, soit environ 85 FCFA la page. Une réponse sans unité perd des points en situation d'intégration.
\end{itemize}
\end{attention}

\courssub{Bilan de l'activité}

Cette activité a permis de mettre en évidence que :

\begin{itemize}
\item les fonctions rationnelles modélisent naturellement les \cle{coûts moyens} en économie, et leur étude complète (dérivée, signe, variations) permet de déterminer un \cle{optimum} concret et exploitable ;
\item l'étape technique décisive est l'étude du signe de la dérivée : soit par factorisation (quand une racine évidente existe), soit par l'étude des variations du numérateur et le théorème des valeurs intermédiaires ;
\item les fonctions du type $\mathrm{e}^{ax+b}$ décrivent les phénomènes de \cle{croissance rapide} (clientèle, population, épidémie, capital avec intérêts composés) et le passage par le logarithme népérien permet de résoudre les équations de seuil ;
\item la lecture graphique et le calcul exact se complètent : le graphique donne une estimation, le calcul la confirme et la précise ;
\item les mathématiques de Terminale A fournissent des outils de décision réels pour un commerçant camerounais : optimisation des coûts, tarification, anticipation des investissements.
\end{itemize}

\begin{savaistu}[Et après ?]
La démarche que tu viens d'utiliser — modéliser par une fonction, étudier ses variations, déterminer un optimum ou un seuil — est exactement celle des \emph{sciences économiques} et de la \emph{gestion} que tu rencontreras à l'université (microéconomie, recherche opérationnelle). Le coût moyen minimal que tu as calculé s'appelle en économie le \emph{seuil de rentabilité technique}, et la croissance exponentielle de la clientèle est le modèle de base du \emph{bouche-à-oreille} en marketing. Les mathématiques de la série A sont bel et bien des mathématiques pour décider.
\end{savaistu}

\newpage

\coursec{Méthodes et exercices résolus}

\coursec{Étude de fonctions}

\courssub{Méthode 1 : Étude complète d'une fonction polynôme de degré $\leq 3$}

\begin{methode}[Plan d'étude d'un polynôme de degré $\leq 3$]
Soit $f(x)=ax^3+bx^2+cx+d$ avec $a\neq 0$.
\begin{enumerate}
\item \textbf{Ensemble de définition} : un polynôme est défini sur $\mathbb{R}$ tout entier : $D_f=\mathbb{R}$.
\item \textbf{Limites aux bornes} : la limite d'un polynôme en $\pm\infty$ est celle de son \cle{terme de plus haut degré} :
\[
\lim_{x\to\pm\infty} f(x)=\lim_{x\to\pm\infty} ax^3.
\]
\item \textbf{Dérivée} : $f'(x)=3ax^2+2bx+c$ (trinôme du second degré).
\item \textbf{Signe de $f'$} : calculer le discriminant $\Delta=(2b)^2-4(3a)c$. Trois cas : $\Delta<0$ (signe constant), $\Delta=0$ (une racine double), $\Delta>0$ (deux racines $x_1,x_2$).
\item \textbf{Tableau de variations} : reporter le signe de $f'$, les limites et les valeurs des \cle{extremums} $f(x_1)$, $f(x_2)$.
\item \textbf{Courbe} : calculer quelques points bien choisis ($f(0)$, les extremums, une ou deux valeurs entières), puis tracer en respectant le tableau.
\end{enumerate}
\textbf{Réflexe} : une cubique admet toujours un point d'inflexion en $x=-\dfrac{b}{3a}$, centre de symétrie de la courbe (non exigible, mais utile pour vérifier le tracé).
\end{methode}

\begin{exoresolu}[Étude de $f(x)=x^3-3x+2$ (type Bac A)]
On considère la fonction $f$ définie sur $\mathbb{R}$ par $f(x)=x^3-3x+2$.
\begin{enumerate}
\item Déterminer les limites de $f$ en $-\infty$ et en $+\infty$.
\item Étudier les variations de $f$ et dresser son tableau de variations.
\item Déterminer les points d'intersection de la courbe $(\mathcal{C}_f)$ avec l'axe des abscisses.
\item Tracer $(\mathcal{C}_f)$ dans un repère orthonormé.
\end{enumerate}
\tcblower
\textbf{1. Limites.} $f$ est un polynôme : sa limite en $\pm\infty$ est celle de son terme de plus haut degré $x^3$.
\[
\lim_{x\to -\infty} x^3=-\infty \quad\text{donc}\quad \lim_{x\to -\infty} f(x)=-\infty ;
\qquad
\lim_{x\to +\infty} x^3=+\infty \quad\text{donc}\quad \lim_{x\to +\infty} f(x)=+\infty.
\]

\textbf{2. Variations.} $f$ est dérivable sur $\mathbb{R}$ et
\[
f'(x)=3x^2-3=3(x^2-1)=3(x-1)(x+1).
\]
$f'$ est un trinôme de racines $x=-1$ et $x=1$, du signe de $a=3>0$ à l'extérieur des racines :
\begin{itemize}
\item $f'(x)>0$ sur $]-\infty;-1[\ \cup\ ]1;+\infty[$ : $f$ est strictement croissante ;
\item $f'(x)<0$ sur $]-1;1[$ : $f$ est strictement décroissante ;
\item $f'(-1)=f'(1)=0$.
\end{itemize}
Valeurs remarquables :
\[
f(-1)=(-1)^3-3(-1)+2=-1+3+2=4 \ ;\qquad f(1)=1-3+2=0.
\]
Tableau de variations :
\begin{center}
\begin{tabular}{|c|ccccccc|}
\hline
$x$ & $-\infty$ & & $-1$ & & $1$ & & $+\infty$ \\
\hline
$f'(x)$ & & $+$ & $0$ & $-$ & $0$ & $+$ & \\
\hline
$f(x)$ & $-\infty$ & $\nearrow$ & $4$ & $\searrow$ & $0$ & $\nearrow$ & $+\infty$ \\
\hline
\end{tabular}
\end{center}
$f$ admet un \cle{maximum local} $4$ en $x=-1$ et un \cle{minimum local} $0$ en $x=1$.

\textbf{3. Intersection avec l'axe des abscisses.} On résout $f(x)=0$. On remarque que $f(1)=0$ : $1$ est racine, donc $x-1$ se met en facteur. Par identification :
\[
x^3-3x+2=(x-1)(x^2+x-2).
\]
Or $x^2+x-2=(x-1)(x+2)$ (car $1$ est racine évidente et le produit des racines vaut $-2$). D'où
\[
f(x)=(x-1)^2(x+2).
\]
$f(x)=0 \iff x=1$ (racine double) ou $x=-2$. La courbe coupe donc l'axe des abscisses en $A(-2\,;\,0)$ et le \emph{touche} en $B(1\,;\,0)$ (racine double : la courbe est tangente à l'axe en ce point, ce qui confirme le minimum local nul).

\textbf{4. Tracé.} Points utiles : $f(0)=2$, $f(-2)=0$, $f(-1)=4$, $f(1)=0$, $f(2)=4$.
\begin{popfigure}
\begin{center}
\begin{tikzpicture}[scale=0.85]
\draw[->,popmuted] (-3.2,0) -- (3.2,0) node[below left]{$x$};
\draw[->,popmuted] (0,-1.5) -- (0,5) node[below left]{$y$};
\draw[popblue,very thick,domain=-2.35:2.25,samples=200] plot (\x,{\x*\x*\x-3*\x+2});
\fill[poporange] (-2,0) circle (2pt) node[below left]{$A$};
\fill[poporange] (1,0) circle (2pt) node[below right]{$B$};
\fill[popink] (-1,4) circle (2pt) node[above]{max $(-1\,;\,4)$};
\draw (0,0) node[below left]{$O$};
\foreach \x in {-2,-1,1,2} \draw (\x,0.06) -- (\x,-0.06) node[below]{\footnotesize $\x$};
\foreach \y in {2,4} \draw (0.06,\y) -- (-0.06,\y) node[left]{\footnotesize $\y$};
\end{tikzpicture}
\end{center}
\legende{Courbe représentative de $f(x)=x^3-3x+2$}
\end{popfigure}
\textbf{Conclusion} : $f$ est croissante sur $]-\infty;-1]$ et $[1;+\infty[$, décroissante sur $[-1;1]$, et sa courbe rencontre l'axe des abscisses en $-2$ et en $1$ (point de tangence).
\end{exoresolu}

\courssub{Méthode 2 : Étude d'une fonction rationnelle homographique du type $\dfrac{ax^2+bx+c}{dx+e}$}

\begin{methode}[Plan d'étude de $f(x)=\dfrac{ax^2+bx+c}{dx+e}$]
\begin{enumerate}
\item \textbf{Ensemble de définition} : résoudre $dx+e=0$, soit $x=-\dfrac{e}{d}$. Alors $D_f=\mathbb{R}\setminus\left\{-\dfrac{e}{d}\right\}$. La droite $x=-\dfrac{e}{d}$ est \cle{asymptote verticale} (à justifier par une limite infinie).
\item \textbf{Limites aux bornes} : en $\pm\infty$, la limite est celle du quotient des termes de plus haut degré : $\lim\limits_{x\to\pm\infty}f(x)=\lim\limits_{x\to\pm\infty}\dfrac{ax^2}{dx}=\lim\limits_{x\to\pm\infty}\dfrac{a}{d}x=\pm\infty$ : \textbf{pas d'asymptote horizontale}.
\item \textbf{Asymptote oblique} : effectuer la \cle{division euclidienne} du numérateur par le dénominateur (ou identifier) pour écrire
\[
f(x)=\alpha x+\beta+\dfrac{\gamma}{dx+e}.
\]
La droite $y=\alpha x+\beta$ est asymptote oblique car $\lim\limits_{x\to\pm\infty}\dfrac{\gamma}{dx+e}=0$.
\item \textbf{Dérivée} : appliquer $\left(\dfrac{u}{v}\right)'=\dfrac{u'v-uv'}{v^2}$. Le numérateur obtenu est un trinôme dont on étudie le signe ; le dénominateur est un carré, donc strictement positif sur $D_f$.
\item \textbf{Tableau de variations} : en n'oubliant pas la \cle{double barre} en $x=-\dfrac{e}{d}$.
\item \textbf{Tracé} : asymptotes en pointillés d'abord, extremums, quelques points.
\end{enumerate}
\textbf{Réflexe} : pour la position de la courbe par rapport à son asymptote oblique, étudier le signe de $\dfrac{\gamma}{dx+e}$.
\end{methode}

\begin{exoresolu}[Étude de $f(x)=\dfrac{x^2-3x+3}{x-1}$ (type Bac A)]
Soit $f$ la fonction définie par $f(x)=\dfrac{x^2-3x+3}{x-1}$ et $(\mathcal{C})$ sa courbe dans un repère orthonormé.
\begin{enumerate}
\item Déterminer l'ensemble de définition de $f$ et les limites aux bornes. Interpréter graphiquement.
\item Montrer que pour tout $x\neq 1$, $f(x)=x-2+\dfrac{1}{x-1}$. En déduire une asymptote oblique et la position relative de $(\mathcal{C})$ par rapport à elle.
\item Étudier les variations de $f$ et dresser le tableau de variations.
\item Tracer $(\mathcal{C})$ et ses asymptotes.
\end{enumerate}
\tcblower
\textbf{1. Ensemble de définition et limites.} $f$ est définie si et seulement si $x-1\neq 0$, donc
\[
D_f=\mathbb{R}\setminus\{1\}=]-\infty;1[\ \cup\ ]1;+\infty[.
\]
\emph{En $1$} : le numérateur en $1$ vaut $1-3+3=1\neq 0$. On a $\lim\limits_{x\to 1}(x^2-3x+3)=1$ et :
\begin{itemize}
\item pour $x>1$, $x-1>0$ et $x-1\to 0^+$ : $\lim\limits_{\substack{x\to 1\\x>1}} f(x)=+\infty$ ;
\item pour $x<1$, $x-1<0$ et $x-1\to 0^-$ : $\lim\limits_{\substack{x\to 1\\x<1}} f(x)=-\infty$.
\end{itemize}
Ces limites infinies prouvent que la droite d'équation $x=1$ est \cle{asymptote verticale} à $(\mathcal{C})$.

\emph{En $\pm\infty$} : $\lim\limits_{x\to\pm\infty} f(x)=\lim\limits_{x\to\pm\infty}\dfrac{x^2}{x}=\lim\limits_{x\to\pm\infty} x$, donc
\[
\lim_{x\to -\infty} f(x)=-\infty \quad\text{et}\quad \lim_{x\to +\infty} f(x)=+\infty.
\]

\textbf{2. Asymptote oblique.} Effectuons la division euclidienne de $x^2-3x+3$ par $x-1$ :
\[
x^2-3x+3=(x-1)(x-2)+1.
\]
Vérification : $(x-1)(x-2)+1=x^2-3x+2+1=x^2-3x+3$. Donc pour tout $x\neq 1$ :
\[
f(x)=x-2+\frac{1}{x-1}.
\]
Comme $\lim\limits_{x\to\pm\infty}\dfrac{1}{x-1}=0$, on a $\lim\limits_{x\to\pm\infty}\big[f(x)-(x-2)\big]=0$ : la droite $(\Delta)$ d'équation $y=x-2$ est \cle{asymptote oblique} à $(\mathcal{C})$ en $-\infty$ et en $+\infty$.

\emph{Position relative} : le signe de $f(x)-(x-2)=\dfrac{1}{x-1}$ est celui de $x-1$ :
\begin{itemize}
\item si $x>1$ : $(\mathcal{C})$ est \textbf{au-dessus} de $(\Delta)$ ;
\item si $x<1$ : $(\mathcal{C})$ est \textbf{en dessous} de $(\Delta)$.
\end{itemize}

\textbf{3. Variations.} $f$ est dérivable sur $D_f$ comme quotient de polynômes. Avec $u=x^2-3x+3$ et $v=x-1$ :
\[
f'(x)=\frac{(2x-3)(x-1)-(x^2-3x+3)\times 1}{(x-1)^2}
=\frac{2x^2-5x+3-x^2+3x-3}{(x-1)^2}
=\frac{x^2-2x}{(x-1)^2}
=\frac{x(x-2)}{(x-1)^2}.
\]
Le dénominateur $(x-1)^2$ est strictement positif sur $D_f$ ; le signe de $f'$ est donc celui du trinôme $x(x-2)$, de racines $0$ et $2$, positif à l'extérieur des racines :
\begin{itemize}
\item $f'(x)>0$ sur $]-\infty;0[\ \cup\ ]2;+\infty[$ ;
\item $f'(x)<0$ sur $]0;1[\ \cup\ ]1;2[$.
\end{itemize}
Valeurs : $f(0)=\dfrac{3}{-1}=-3$ ; $f(2)=\dfrac{4-6+3}{1}=1$.
\begin{center}
\begin{tabular}{|c|ccccccccc|}
\hline
$x$ & $-\infty$ & & $0$ & & $1$ & & $2$ & & $+\infty$ \\
\hline
$f'(x)$ & & $+$ & $0$ & $-$ & $\|$ & $-$ & $0$ & $+$ & \\
\hline
$f(x)$ & $-\infty$ & $\nearrow$ & $-3$ & $\searrow$ & $\|\ ^{-\infty}_{+\infty}$ & $\searrow$ & $1$ & $\nearrow$ & $+\infty$ \\
\hline
\end{tabular}
\end{center}
$f$ admet un maximum local $-3$ en $0$ et un minimum local $1$ en $2$.

\textbf{4. Tracé.}
\begin{popfigure}
\begin{center}
\begin{tikzpicture}[scale=0.7]
\draw[->,popmuted] (-4.5,0) -- (5.5,0) node[below left]{$x$};
\draw[->,popmuted] (0,-5.5) -- (0,5.5) node[below left]{$y$};
\draw[poporange,dashed,thick] (1,-5.5) -- (1,5.5);
\draw[popgreen,dashed,thick,domain=-4:5.2] plot (\x,{\x-2});
\draw[popblue,very thick,domain=-4.2:0.82,samples=200] plot (\x,{(\x*\x-3*\x+3)/(\x-1)});
\draw[popblue,very thick,domain=1.24:5.3,samples=200] plot (\x,{(\x*\x-3*\x+3)/(\x-1)});
\fill[popink] (0,-3) circle (2pt) node[below left]{\footnotesize max $(0;-3)$};
\fill[popink] (2,1) circle (2pt) node[above right]{\footnotesize min $(2;1)$};
\draw (0,0) node[below left]{$O$};
\node[poporange] at (1.35,4.6) {\footnotesize $x=1$};
\node[popgreen] at (4.6,2.1) {\footnotesize $y=x-2$};
\end{tikzpicture}
\end{center}
\legende{Courbe de $f$ et ses deux asymptotes}
\end{popfigure}
\textbf{Conclusion} : $(\mathcal{C})$ admet l'asymptote verticale $x=1$, l'asymptote oblique $y=x-2$, un maximum local en $(0;-3)$ et un minimum local en $(2;1)$.
\end{exoresolu}

\courssub{Méthode 3 : Étude d'une fonction du type $x\mapsto \ln(ax+b)$}

\begin{methode}[Plan d'étude de $f(x)=\ln(ax+b)$, $a\neq 0$]
\begin{enumerate}
\item \textbf{Ensemble de définition} : le logarithme n'existe que pour un argument \cle{strictement positif} : résoudre $ax+b>0$.
\begin{itemize}
\item si $a>0$ : $D_f=\left]-\dfrac{b}{a};+\infty\right[$ ;
\item si $a<0$ : $D_f=\left]-\infty;-\dfrac{b}{a}\right[$.
\end{itemize}
\item \textbf{Limites} : en $+\infty$ (si $a>0$), $\lim\limits_{x\to+\infty}\ln(ax+b)=+\infty$. À la borne finie $x\to -\dfrac{b}{a}$, l'argument tend vers $0^+$ : $\lim f(x)=-\infty$ : \cle{asymptote verticale} $x=-\dfrac{b}{a}$.
\item \textbf{Dérivée} : formule $(\ln u)'=\dfrac{u'}{u}$ :
\[
f'(x)=\frac{a}{ax+b}.
\]
Sur $D_f$, $ax+b>0$, donc $f'$ a le \textbf{signe de $a$} : si $a>0$, $f$ est strictement croissante ; si $a<0$, strictement décroissante.
\item \textbf{Tableau de variations} puis \textbf{tracé} : point d'intersection avec l'axe des abscisses obtenu en résolvant $\ln(ax+b)=0 \iff ax+b=1 \iff x=\dfrac{1-b}{a}$.
\end{enumerate}
\textbf{Formules à retenir} : $\ln 1=0$, $\ln e=1$, $\ln(ab)=\ln a+\ln b$, $\ln\dfrac{1}{x}=-\ln x$, et $\ln x<0 \iff 0<x<1$.
\end{methode}

\begin{exoresolu}[Étude de $f(x)=\ln(2x-4)$ (type Bac A)]
Soit $f$ la fonction définie par $f(x)=\ln(2x-4)$ et $(\mathcal{C})$ sa courbe représentative dans un repère orthonormé.
\begin{enumerate}
\item Déterminer l'ensemble de définition de $f$.
\item Calculer les limites de $f$ aux bornes de son ensemble de définition et interpréter graphiquement le résultat.
\item Étudier les variations de $f$ et dresser son tableau de variations.
\item Déterminer les coordonnées du point d'intersection de $(\mathcal{C})$ avec l'axe des abscisses, puis tracer $(\mathcal{C})$.
\item Résoudre dans $D_f$ l'inéquation $f(x)\geq 0$.
\end{enumerate}
\tcblower
\textbf{1. Ensemble de définition.} $f(x)$ existe si et seulement si $2x-4>0 \iff x>2$. Donc
\[
D_f=]2;+\infty[.
\]

\textbf{2. Limites.}
\emph{En $2^+$} : $\lim\limits_{\substack{x\to 2\\x>2}}(2x-4)=0^+$ et $\lim\limits_{X\to 0^+}\ln X=-\infty$, donc par composition :
\[
\lim_{\substack{x\to 2\\x>2}} f(x)=-\infty.
\]
La droite d'équation $x=2$ est donc \cle{asymptote verticale} à $(\mathcal{C})$.

\emph{En $+\infty$} : $\lim\limits_{x\to+\infty}(2x-4)=+\infty$ et $\lim\limits_{X\to+\infty}\ln X=+\infty$, donc
\[
\lim_{x\to+\infty} f(x)=+\infty.
\]

\textbf{3. Variations.} $f$ est dérivable sur $D_f$ (composée de fonctions dérivables, l'argument restant strictement positif) et, par la formule $(\ln u)'=\dfrac{u'}{u}$ :
\[
f'(x)=\frac{2}{2x-4}=\frac{1}{x-2}.
\]
Sur $]2;+\infty[$, $x-2>0$, donc $f'(x)>0$ : $f$ est \textbf{strictement croissante} sur $D_f$.
\begin{center}
\begin{tabular}{|c|ccc|}
\hline
$x$ & $2$ & & $+\infty$ \\
\hline
$f'(x)$ & $\|$ & $+$ & \\
\hline
$f(x)$ & $\|\ -\infty$ & $\nearrow$ & $+\infty$ \\
\hline
\end{tabular}
\end{center}

\textbf{4. Intersection avec l'axe des abscisses.}
\[
f(x)=0 \iff \ln(2x-4)=0 \iff 2x-4=e^0=1 \iff x=\frac{5}{2}.
\]
Comme $\dfrac{5}{2}=2,5\in D_f$, la courbe coupe l'axe des abscisses au point $A\left(\dfrac{5}{2}\,;\,0\right)$.
Autre point utile : $f(4)=\ln 4\approx 1,39$.
\begin{popfigure}
\begin{center}
\begin{tikzpicture}[scale=0.9]
\draw[->,popmuted] (-0.5,0) -- (6,0) node[below left]{$x$};
\draw[->,popmuted] (0,-3.5) -- (0,3) node[below left]{$y$};
\draw[poporange,dashed,thick] (2,-3.5) -- (2,3);
\draw[popblue,very thick,domain=2.06:5.8,samples=200] plot (\x,{ln(2*\x-4)});
\fill[popink] (2.5,0) circle (2pt) node[above left]{\footnotesize $A\left(\tfrac{5}{2};0\right)$};
\draw (0,0) node[below left]{$O$};
\foreach \x in {1,2,3,4,5} \draw (\x,0.06) -- (\x,-0.06) node[below]{\footnotesize $\x$};
\foreach \y in {1,2} \draw (0.06,\y) -- (-0.06,\y) node[left]{\footnotesize $\y$};
\node[poporange] at (2.4,2.4) {\footnotesize $x=2$};
\end{tikzpicture}
\end{center}
\legende{Courbe de $f(x)=\ln(2x-4)$ et son asymptote verticale}
\end{popfigure}

\textbf{5. Inéquation $f(x)\geq 0$.} La fonction $\ln$ est croissante et $\ln 1=0$ :
\[
\ln(2x-4)\geq 0 \iff 2x-4\geq 1 \iff x\geq \frac{5}{2}.
\]
Toutes les valeurs obtenues appartiennent à $D_f$, donc
\[
\mathcal{S}=\left[\frac{5}{2};+\infty\right[.
\]
\textbf{Conclusion} : $f$ est strictement croissante sur $]2;+\infty[$, de $-\infty$ à $+\infty$, s'annule en $x=\dfrac{5}{2}$ et admet l'asymptote verticale $x=2$.
\end{exoresolu}

\courssub{Méthode 4 : Étude d'une fonction du type $x\mapsto e^{ax+b}$}

\begin{methode}[Plan d'étude de $f(x)=e^{ax+b}$, $a\neq 0$]
\begin{enumerate}
\item \textbf{Ensemble de définition} : l'exponentielle est définie partout : $D_f=\mathbb{R}$.
\item \textbf{Signe} : pour tout réel $x$, $f(x)>0$ : la courbe est toujours \cle{au-dessus} de l'axe des abscisses.
\item \textbf{Limites} (par composition) :
\begin{itemize}
\item si $a>0$ : $\lim\limits_{x\to+\infty} e^{ax+b}=+\infty$ et $\lim\limits_{x\to-\infty} e^{ax+b}=0$ : l'axe des abscisses est \cle{asymptote horizontale} en $-\infty$ ;
\item si $a<0$ : c'est l'inverse (asymptote horizontale $y=0$ en $+\infty$).
\end{itemize}
\item \textbf{Dérivée} : formule $(e^u)'=u'\,e^u$ :
\[
f'(x)=a\,e^{ax+b}.
\]
Comme $e^{ax+b}>0$, $f'$ a le \textbf{signe de $a$} : croissante si $a>0$, décroissante si $a<0$.
\item \textbf{Tableau et tracé} : point remarquable $f(0)=e^{b}$ ; on résout $e^{ax+b}=k$ (avec $k>0$) par $ax+b=\ln k$.
\end{enumerate}
\textbf{Formules à retenir} : $e^0=1$, $e^{a+b}=e^a\,e^b$, $e^{-x}=\dfrac{1}{e^x}$, $\ln(e^x)=x$, $e^{\ln x}=x$ pour $x>0$.
\end{methode}

\begin{exoresolu}[Étude de $f(x)=e^{2x-1}$ (type Bac A)]
Soit $f$ la fonction définie sur $\mathbb{R}$ par $f(x)=e^{2x-1}$ et $(\mathcal{C})$ sa courbe dans un repère orthonormé.
\begin{enumerate}
\item Calculer les limites de $f$ en $-\infty$ et en $+\infty$. Interpréter graphiquement.
\item Étudier les variations de $f$ et dresser son tableau de variations.
\item Déterminer une équation de la tangente $(T)$ à $(\mathcal{C})$ au point d'abscisse $\dfrac{1}{2}$.
\item Tracer $(T)$ et $(\mathcal{C})$.
\item Résoudre dans $\mathbb{R}$ l'équation $f(x)=e^3$.
\end{enumerate}
\tcblower
\textbf{1. Limites.} Par composition, avec $X=2x-1$ :
\[
\lim_{x\to +\infty}(2x-1)=+\infty \ \text{et}\ \lim_{X\to+\infty}e^X=+\infty \quad\Longrightarrow\quad \lim_{x\to+\infty} f(x)=+\infty ;
\]
\[
\lim_{x\to -\infty}(2x-1)=-\infty \ \text{et}\ \lim_{X\to-\infty}e^X=0 \quad\Longrightarrow\quad \lim_{x\to-\infty} f(x)=0.
\]
La droite d'équation $y=0$ (l'axe des abscisses) est \cle{asymptote horizontale} à $(\mathcal{C})$ en $-\infty$.

\textbf{2. Variations.} $f$ est dérivable sur $\mathbb{R}$ et, par la formule $(e^u)'=u'e^u$ :
\[
f'(x)=2\,e^{2x-1}.
\]
Pour tout réel $x$, $e^{2x-1}>0$, donc $f'(x)>0$ : $f$ est \textbf{strictement croissante} sur $\mathbb{R}$.
\begin{center}
\begin{tabular}{|c|ccc|}
\hline
$x$ & $-\infty$ & & $+\infty$ \\
\hline
$f'(x)$ & & $+$ & \\
\hline
$f(x)$ & $0$ & $\nearrow$ & $+\infty$ \\
\hline
\end{tabular}
\end{center}

\textbf{3. Tangente en $x=\dfrac{1}{2}$.} Équation de la tangente : $y=f'\left(\tfrac12\right)\left(x-\tfrac12\right)+f\left(\tfrac12\right)$.
\[
f\left(\tfrac12\right)=e^{2\times\frac12-1}=e^0=1 \ ;\qquad f'\left(\tfrac12\right)=2e^0=2.
\]
Donc
\[
(T):\ y=2\left(x-\frac12\right)+1=2x.
\]
Remarque élégante : la tangente passe par l'origine.

\textbf{4. Tracé.} Points utiles : $f(0)=e^{-1}\approx 0,37$ ; $f\left(\tfrac12\right)=1$ ; $f(1)=e\approx 2,72$.
\begin{popfigure}
\begin{center}
\begin{tikzpicture}[scale=0.95]
\draw[->,popmuted] (-2.5,0) -- (2.5,0) node[below left]{$x$};
\draw[->,popmuted] (0,-1) -- (0,4) node[below left]{$y$};
\draw[popgreen,dashed,thick,domain=-1.6:1.7] plot (\x,{2*\x});
\draw[popblue,very thick,domain=-2.3:1.35,samples=200] plot (\x,{exp(2*\x-1)});
\fill[popink] (0.5,1) circle (2pt) node[above left]{\footnotesize $\left(\tfrac12;1\right)$};
\draw (0,0) node[below left]{$O$};
\foreach \x in {-2,-1,1,2} \draw (\x,0.05) -- (\x,-0.05) node[below]{\footnotesize $\x$};
\foreach \y in {1,2,3} \draw (0.05,\y) -- (-0.05,\y) node[left]{\footnotesize $\y$};
\node[popgreen] at (1.55,2.6) {\footnotesize $(T):y=2x$};
\end{tikzpicture}
\end{center}
\legende{Courbe de $f(x)=e^{2x-1}$ et sa tangente en $x=\tfrac12$}
\end{popfigure}

\textbf{5. Équation $f(x)=e^3$.} La fonction exponentielle est une bijection (strictement croissante), donc :
\[
e^{2x-1}=e^3 \iff 2x-1=3 \iff x=2.
\]
$\mathcal{S}=\{2\}$.
\textbf{Conclusion} : $f$ est strictement croissante de $0$ vers $+\infty$, admet l'asymptote horizontale $y=0$ en $-\infty$, et sa tangente au point $\left(\tfrac12;1\right)$ a pour équation $y=2x$.
\end{exoresolu}

\courssub{Méthode 5 : Synthèse — conduire seul une étude complète de fonction}

\begin{methode}[La grille d'étude universelle en 7 questions]
Face à \emph{n'importe quelle} fonction au Bac, réponds dans l'ordre :
\begin{enumerate}
\item \textbf{Où vit-elle ?} Ensemble de définition : dénominateur $\neq 0$, argument de $\ln$ $>0$, contenu d'une racine $\geq 0$.
\item \textbf{Que fait-elle aux bords ?} Limites aux bornes \emph{ouvertes} de $D_f$ (il y en a 2, 3 ou 4 selon les trous du domaine).
\item \textbf{Y a-t-il des murs ou des rails ?} Limite infinie en un point fini $\Rightarrow$ asymptote verticale ; limite finie $\ell$ en $\pm\infty$ $\Rightarrow$ asymptote horizontale $y=\ell$ ; sinon, pour les rationnelles, chercher l'asymptote oblique par division.
\item \textbf{Monte-t-elle ou descend-elle ?} Calcul de $f'$, factorisation maximale, étude de signe (tableau de signes si nécessaire).
\item \textbf{Tableau de variations} : signe de $f'$ en première ligne, flèches de $f$ en seconde, avec limites et extremums ; double barre aux valeurs interdites.
\item \textbf{Points stratégiques} : intersections avec les axes, extremums, tangentes demandées.
\item \textbf{Tracé soigné} : asymptotes en pointillés d'abord, points remarquables, puis la courbe en respectant le tableau.
\end{enumerate}
\textbf{Réflexe de vérification} : la courbe tracée doit être \emph{cohérente} avec le tableau de variations (croissance, limites, asymptotes). Toute contradiction signale une erreur.
\end{methode}

\begin{savaistu}[Comparaison des croissances et applications au Cameroun]
En Terminale, on utilise souvent la hiérarchie des croissances au voisinage de $+\infty$ :
\[
\ln x \ll x^a \ll e^x \ll x^x \quad (a>0).
\]
Cela permet de deviner les limites sans calcul formel.
\par\textbf{Exemples concrets :}
\begin{itemize}
\item \textbf{Démographie} : la population du Cameroun ($\approx 28$ millions en 2024) suit grossièrement un modèle exponentiel $P(t)=P_0 e^{kt}$ tant que les ressources ne limitent pas la croissance.
\item \textbf{Finance (MTN MoMo / Orange Money)} : un placement à taux composé suit une exponentielle ; le temps pour doubler un capital est $\frac{\ln 2}{k}$ (règle des 70 : $70/taux\%$).
\item \textbf{Électrification rurale (AER/SONATREL)} : le coût de raccordement par km de ligne suit souvent une fonction rationnelle (coût fixe + coût variable) ou une racine carrée selon le terrain.
\item \textbf{Télécoms} : la capacité d'un lien (théorème de Shannon) est logarithmique en le rapport signal/bruit : $C = B \log_2(1+S/N)$.
\end{itemize}
\end{savaistu}

\begin{exoresolu}[Étude guidée de $f(x)=\dfrac{x^2+1}{x}$ puis lecture graphique (type Bac A)]
Soit $f$ la fonction définie par $f(x)=\dfrac{x^2+1}{x}$.
\begin{enumerate}
\item Préciser $D_f$ et montrer que $f(x)=x+\dfrac{1}{x}$.
\item Calculer les limites aux quatre bornes de $D_f$ et en déduire les asymptotes de $(\mathcal{C})$.
\item Étudier les variations de $f$ et dresser le tableau de variations complet.
\item Montrer que $f$ est impaire. Qu'en déduit-on pour $(\mathcal{C})$ ?
\end{enumerate}
\tcblower
\textbf{1. Ensemble de définition et écriture simplifiée.} $f$ est définie pour $x\neq 0$ : $D_f=\mathbb{R}^*=]-\infty;0[\ \cup\ ]0;+\infty[$. Pour tout $x\neq 0$ :
\[
f(x)=\frac{x^2}{x}+\frac{1}{x}=x+\frac{1}{x}.
\]

\textbf{2. Limites et asymptotes.}
\emph{En $0$} : $\lim\limits_{x\to 0}x=0$ et $\lim\limits_{\substack{x\to 0\\x>0}}\dfrac{1}{x}=+\infty$, $\lim\limits_{\substack{x\to 0\\x<0}}\dfrac{1}{x}=-\infty$. Donc
\[
\lim_{\substack{x\to 0\\x>0}} f(x)=+\infty \ ;\qquad \lim_{\substack{x\to 0\\x<0}} f(x)=-\infty.
\]
La droite $x=0$ (l'axe des ordonnées) est \cle{asymptote verticale}.

\emph{En $\pm\infty$} : $\lim\limits_{x\to\pm\infty}\dfrac{1}{x}=0$ donc $\lim\limits_{x\to+\infty}f(x)=+\infty$ et $\lim\limits_{x\to-\infty}f(x)=-\infty$. De plus
\[
\lim_{x\to\pm\infty}\big[f(x)-x\big]=\lim_{x\to\pm\infty}\frac{1}{x}=0 :
\]
la droite $y=x$ est \cle{asymptote oblique} en $-\infty$ et en $+\infty$. La courbe est au-dessus de cette asymptote pour $x>0$ et en dessous pour $x<0$.

\textbf{3. Variations.} $f$ est dérivable sur $\mathbb{R}^*$ et
\[
f'(x)=1-\frac{1}{x^2}=\frac{x^2-1}{x^2}=\frac{(x-1)(x+1)}{x^2}.
\]
Le dénominateur $x^2$ est strictement positif sur $\mathbb{R}^*$ ; le signe de $f'$ est celui de $(x-1)(x+1)$ : positif sur $]-\infty;-1[\ \cup\ ]1;+\infty[$, négatif sur $]-1;0[\ \cup\ ]0;1[$.
Valeurs : $f(-1)=-2$ et $f(1)=2$.
\begin{center}
\begin{tabular}{|c|ccccccccc|}
\hline
$x$ & $-\infty$ & & $-1$ & & $0$ & & $1$ & & $+\infty$ \\
\hline
$f'(x)$ & & $+$ & $0$ & $-$ & $\|$ & $-$ & $0$ & $+$ & \\
\hline
$f(x)$ & $-\infty$ & $\nearrow$ & $-2$ & $\searrow$ & $\|\ ^{-\infty}_{+\infty}$ & $\searrow$ & $2$ & $\nearrow$ & $+\infty$ \\
\hline
\end{tabular}
\end{center}

\textbf{4. Parité.} Pour tout $x\neq 0$, $-x\in D_f$ et
\[
f(-x)=-x+\frac{1}{-x}=-\left(x+\frac{1}{x}\right)=-f(x).
\]
$f$ est \cle{impaire} : sa courbe est symétrique par rapport à l'origine $O$ du repère. Il suffisait donc d'étudier $f$ sur $]0;+\infty[$ et de compléter par symétrie centrale.

\textbf{5. Tracé.}
\begin{popfigure}
\begin{center}
\begin{tikzpicture}[scale=0.75]
\draw[->,popmuted] (-4.5,0) -- (4.5,0) node[below left]{$x$};
\draw[->,popmuted] (0,-4.5) -- (0,4.5) node[below left]{$y$};
\draw[poporange,dashed,thick] (0,-4.5) -- (0,4.5);
\draw[popgreen,dashed,thick,domain=-4.2:4.2] plot (\x,{\x});
\draw[popblue,very thick,domain=-4.2:-0.3,samples=200] plot (\x,{\x+1/\x});
\draw[popblue,very thick,domain=0.3:4.2,samples=200] plot (\x,{\x+1/\x});
\fill[popink] (-1,-2) circle (2pt) node[below left]{\footnotesize max $(-1;-2)$};
\fill[popink] (1,2) circle (2pt) node[above right]{\footnotesize min $(1;2)$};
\draw (0,0) node[below left]{$O$};
\node[poporange] at (-0.5,4) {\footnotesize $x=0$};
\node[popgreen] at (3.8,3.5) {\footnotesize $y=x$};
\end{tikzpicture}
\end{center}
\legende{Courbe de $f(x)=x+\frac{1}{x}$, asymptotes et extremums}
\end{popfigure}
\textbf{Conclusion} : $(\mathcal{C})$ admet les asymptotes $x=0$ et $y=x$, un maximum local $-2$ en $-1$, un minimum local $2$ en $1$, et est symétrique par rapport à $O$.
\end{exoresolu}

\courssub{Méthode 6 : Résoudre un problème concret de modélisation et d'optimisation}

\begin{methode}[Modéliser puis optimiser avec une fonction]
\begin{enumerate}
\item \textbf{Identifier la variable} $x$ (préciser son unité et son domaine réel : souvent $x\geq 0$, parfois un intervalle borné imposé par le contexte).
\item \textbf{Traduire l'énoncé} en une fonction $f(x)$ : coût, recette, bénéfice, distance, durée, population...
\item \textbf{Restreindre l'étude} au domaine réel du problème, même si la fonction est définie plus largement.
\item \textbf{Optimiser} : calculer $f'$, étudier son signe, conclure sur le maximum ou le minimum \emph{sur le domaine du problème} (attention aux bornes !).
\item \textbf{Conclure en français} avec une phrase complète et les unités : « Le bénéfice est maximal pour $x=\ldots$ unités et vaut $\ldots$ FCFA ».
\end{enumerate}
\textbf{Réflexes} : bénéfice $=$ recette $-$ coût ; recette $=$ prix $\times$ quantité ; un extremum local n'est la solution que s'il appartient au domaine réel et s'il y est bien global (vérifier les bornes).
\end{methode}

\begin{exoresolu}[Bénéfice maximal d'un atelier de couture à Douala (type Bac A)]
Un atelier de confection de Douala produit chaque semaine $x$ dizaines de pagnes, avec $x\in[0;8]$. Le bénéfice hebdomadaire, exprimé en \textbf{dizaines de milliers de FCFA}, est modélisé par
\[
B(x)=-x^3+9x^2-15x \qquad (x\in[0;8]).
\]
\begin{enumerate}
\item Étudier les variations de $B$ sur $[0;8]$ et dresser son tableau de variations.
\item En déduire la production hebdomadaire qui maximise le bénéfice, et calculer ce bénéfice maximal en FCFA.
\item Déterminer les productions pour lesquelles l'atelier est bénéficiaire ($B(x)>0$).
\end{enumerate}
\tcblower
\textbf{1. Variations de $B$ sur $[0;8]$.} $B$ est un polynôme, dérivable sur $\mathbb{R}$, donc sur $[0;8]$ :
\[
B'(x)=-3x^2+18x-15=-3(x^2-6x+5)=-3(x-1)(x-5).
\]
Les racines de $B'$ sont $x=1$ et $x=5$. Le trinôme $(x-1)(x-5)$ est positif à l'extérieur des racines ; avec le coefficient $-3<0$ :
\begin{itemize}
\item $B'(x)<0$ sur $[0;1[\ \cup\ ]5;8]$ : $B$ décroît ;
\item $B'(x)>0$ sur $]1;5[$ : $B$ croît.
\end{itemize}
Valeurs aux points clés :
\begin{align*}
B(0)&=0 ;\\
B(1)&=-1+9-15=-7 ;\\
B(5)&=-125+225-75=25 ;\\
B(8)&=-512+576-120=-56.
\end{align*}
\begin{center}
\begin{tabular}{|c|ccccccc|}
\hline
$x$ & $0$ & & $1$ & & $5$ & & $8$ \\
\hline
$B'(x)$ & & $-$ & $0$ & $+$ & $0$ & $-$ & \\
\hline
$B(x)$ & $0$ & $\searrow$ & $-7$ & $\nearrow$ & $25$ & $\searrow$ & $-56$ \\
\hline
\end{tabular}
\end{center}

\textbf{2. Bénéfice maximal.} D'après le tableau, sur $[0;8]$, le maximum de $B$ est atteint en $x=5$ et vaut $B(5)=25$.
Or $x$ est en dizaines de pagnes et $B$ en dizaines de milliers de FCFA :
\[
x=5 \iff 50\ \text{pagnes} \ ;\qquad B(5)=25 \iff 25\times \num{10000}=\num{250000}\ \text{FCFA}.
\]
\textbf{Conclusion} : l'atelier maximise son bénéfice hebdomadaire en produisant \textbf{50 pagnes par semaine} ; le bénéfice maximal est de \textbf{250 000 FCFA}.

\textbf{3. Productions bénéficiaires.} On résout $B(x)>0$ sur $[0;8]$. Factorisons :
\[
B(x)=x(-x^2+9x-15).
\]
Discriminant de $-x^2+9x-15$ : $\Delta=81-60=21$, racines $x_1=\dfrac{9-\sqrt{21}}{2}\approx 2,21$ et $x_2=\dfrac{9+\sqrt{21}}{2}\approx 6,79$. Sur $[0;8]$, $B(x)=x(-x^2+9x-15)$ : pour $x>0$, le signe de $B$ est celui du trinôme $-x^2+9x-15$, négatif à l'extérieur de $[x_1;x_2]$ et positif à l'intérieur. Donc
\[
B(x)>0 \iff x\in\left]\frac{9-\sqrt{21}}{2}\,;\,\frac{9+\sqrt{21}}{2}\right[ \approx \ ]2,21\,;\,6,79[.
\]
\textbf{Conclusion} : l'atelier est bénéficiaire pour une production comprise entre environ $23$ et $67$ pagnes par semaine (bornes exclues), c'est-à-dire en pratique entre $23$ et $67$ pagnes.

\textbf{4. Illustration graphique.}
\begin{popfigure}
\begin{center}
\begin{tikzpicture}[scale=0.6]
\draw[->,popmuted] (-0.5,0) -- (8.5,0) node[below left]{$x$ (dizaines de pagnes)};
\draw[->,popmuted] (0,-6) -- (0,28) node[below left]{$B(x)$ (dizaines de milliers FCFA)};
\draw[popblue,very thick,domain=0:8,samples=200] plot (\x,{-\x*\x*\x+9*\x*\x-15*\x});
\draw[popmuted,dashed] (5,0) -- (5,25) node[above]{\footnotesize $x=5$};
\draw[popmuted,dashed] (0,25) -- (5,25) node[left]{\footnotesize $25$};
\fill[popink] (0,0) circle (2pt) node[below left]{\footnotesize $O$};
\fill[popink] (1,-7) circle (2pt) node[below]{\footnotesize min};
\fill[popink] (5,25) circle (2pt) node[above right]{\footnotesize max};
\fill[popink] (8,-56) circle (2pt) node[below right]{\footnotesize fin};
\foreach \x in {1,2,3,4,5,6,7,8} \draw (\x,0.2) -- (\x,-0.2) node[below]{\footnotesize $\x$};
\foreach \y in {5,10,15,20,25} \draw (0.1,\y) -- (-0.1,\y) node[left]{\footnotesize $\y$};
\node[popblue] at (4,15) {\footnotesize $B(x)$};
\end{tikzpicture}
\end{center}
\legende{Variations du bénéfice $B(x)$ sur $[0;8]$}
\end{popfigure}
\end{exoresolu}

\begin{attention}[Les 5 erreurs qui coûtent des points au Bac]
\begin{enumerate}
\item \textbf{Oublier l'ensemble de définition ou le fausser.} Pour $\ln(ax+b)$, la condition est $ax+b>0$ (\emph{strictement}) ; pour $\dfrac{P(x)}{dx+e}$, exclure $x=-\dfrac{e}{d}$. Résoudre une équation ou une inéquation \emph{hors de $D_f$} et donner des solutions interdites est une faute éliminatoire : toujours vérifier que les solutions appartiennent au domaine.
\item \textbf{Annoncer une asymptote sans limite justificative.} On n'écrit jamais « $x=1$ est asymptote » sans avoir calculé $\lim\limits_{x\to 1} f(x)=\pm\infty$. De même, une asymptote oblique $y=\alpha x+\beta$ exige la justification $\lim\limits_{x\to\pm\infty}[f(x)-(\alpha x+\beta)]=0$.
\item \textbf{Se tromper dans les dérivées composées.} Les formules sont $\big(\ln(ax+b)\big)'=\dfrac{a}{ax+b}$ et $\big(e^{ax+b}\big)'=a\,e^{ax+b}$ : le facteur $a$ est obligatoire. Pour $\dfrac{u}{v}$, c'est $\dfrac{u'v-uv'}{v^2}$ : attention au signe moins et à ne pas « simplifier » le carré du dénominateur avec un seul facteur du numérateur.
\item \textbf{Étudier le signe de $f'$ au lieu de celui de son numérateur.} Pour $f'(x)=\dfrac{x(x-2)}{(x-1)^2}$, le dénominateur est un carré strictement positif : inutile de le faire figurer dans le tableau de signes, mais il faut la \textbf{double barre} en $x=1$ dans le tableau de variations. Oublier cette double barre, ou faire « sauter » la flèche par-dessus la valeur interdite, est sanctionné.
\item \textbf{Répondre hors contexte dans les problèmes.} Dans un problème d'optimisation, chercher l'extremum \emph{sur le domaine réel} (ex. $[0;8]$), comparer avec les valeurs aux bornes, puis conclure par une phrase avec les bonnes unités et la bonne échelle ($x=5$ dizaines $=50$ pagnes ; $B=25$ dizaines de milliers $= \num{250000}$ FCFA). Une réponse « $x=5$ » sans conversion ni phrase de conclusion perd des points.
\end{enumerate}
\end{attention}

\newpage

\coursec{Exercices}

\courssub{Je vérifie mes connaissances}

\begin{exercice}[QCM : dérivées et limites]
Pour chaque question, une seule réponse est exacte. Recopie la lettre correspondante.
\begin{enumerate}
\item La dérivée de $f(x) = x^3 - 3x^2 - 9x + 5$ est :
\begin{enumerate}
\item[a)] $f'(x) = 3x^2 - 6x - 9$ \quad b) $f'(x) = x^2 - 6x - 9$ \quad c) $f'(x) = 3x^2 - 3x - 9$
\end{enumerate}
\item L'ensemble de définition de $g(x) = \ln(2x + 4)$ est :
\begin{enumerate}
\item[a)] $]0\ ; +\infty[$ \quad b) $]-2\ ; +\infty[$ \quad c) $[-2\ ; +\infty[$
\end{enumerate}
\item $\displaystyle\lim_{x \to -\infty} \mathrm{e}^{2x}$ est égale à :
\begin{enumerate}
\item[a)] $+\infty$ \quad b) $0$ \quad c) $-\infty$
\end{enumerate}
\item La fonction $h(x) = \dfrac{x^2 - 3x + 3}{x - 2}$ admet une asymptote verticale d'équation :
\begin{enumerate}
\item[a)] $x = 2$ \quad b) $y = 2$ \quad c) $x = -2$
\end{enumerate}
\end{enumerate}
\end{exercice}

\begin{exercice}[Vrai ou faux, en justifiant]
On considère la fonction $f$ définie sur $\mathbb{R}$ par $f(x) = \mathrm{e}^{2x} - 2\mathrm{e}^x + 1$. Pour chacune des affirmations suivantes, dire si elle est vraie ou fausse, en justifiant la réponse.
\begin{enumerate}
\item Pour tout réel $x$, $f(x) = \left(\mathrm{e}^x - 1\right)^2$.
\item $f'(x) = 2\mathrm{e}^{2x} - 2$.
\item $\displaystyle\lim_{x \to +\infty} f(x) = +\infty$.
\item L'équation $f(x) = 0$ admet une unique solution dans $\mathbb{R}$.
\end{enumerate}
\end{exercice}

\begin{exercice}[Définitions et cours]
\begin{enumerate}
\item Donner la définition d'une \cle{asymptote verticale} à la courbe d'une fonction en un point d'abscisse $a$.
\item Donner la définition d'une \cle{asymptote oblique} d'équation $y = ax + b$ en $+\infty$.
\item Rappeler la formule donnant l'équation de la tangente à la courbe de $f$ au point d'abscisse $a$.
\item Citer les limites de référence : $\displaystyle\lim_{x \to 0^+} \ln x$ et $\displaystyle\lim_{x \to +\infty} \dfrac{\ln x}{x}$.
\end{enumerate}
\end{exercice}

\begin{exercice}[Calculs directs]
\begin{enumerate}
\item Déterminer l'ensemble de définition de chacune des fonctions suivantes :
\[ f(x) = \dfrac{x^2 + 1}{x - 3} \;; \qquad g(x) = \ln(3x - 6) \;; \qquad h(x) = \mathrm{e}^{-2x + 1}. \]
\item Calculer la dérivée de chacune de ces trois fonctions.
\item Calculer $\displaystyle\lim_{x \to +\infty} \left(x^3 - 3x^2 - 9x + 5\right)$ et $\displaystyle\lim_{x \to -\infty} \left(x^3 - 3x^2 - 9x + 5\right)$.
\end{enumerate}
\end{exercice}

\courssub{Je m'entraîne}

\begin{exercice}[Variations d'un polynôme de degré 3]
Soit $f$ la fonction définie sur $\mathbb{R}$ par $f(x) = x^3 - 3x^2 - 9x + 5$.
\begin{enumerate}
\item Calculer $f'(x)$ et dresser le tableau de variations de $f$ sur $\mathbb{R}$.
\item Déterminer les extremums locaux de $f$ et préciser leur nature.
\item Déterminer une équation de la tangente à la courbe de $f$ au point d'abscisse $0$.
\item Tracer la courbe de $f$ dans un repère orthonormé (unité : \SI{1}{cm}), puis résoudre graphiquement l'équation $f(x) = 0$ (donner un encadrement d'amplitude $0,5$ de chaque solution).
\end{enumerate}
\end{exercice}

\begin{exercice}[Fonction rationnelle et coût moyen]
Une menuiserie de Douala produit $x$ meubles par semaine ($x > 0$). Le coût moyen de production d'un meuble, en dizaines de milliers de FCFA, est donné par :
\[ CM(x) = \dfrac{x^2 + 4x + 16}{x}. \]
\begin{enumerate}
\item Montrer que pour tout $x > 0$ : $CM(x) = x + 4 + \dfrac{16}{x}$.
\item Calculer $CM'(x)$ et étudier son signe sur $]0\ ; +\infty[$.
\item Dresser le tableau de variations de $CM$.
\item En déduire le nombre de meubles à produire pour que le coût moyen soit minimal. Quel est alors ce coût moyen minimal, en FCFA ?
\end{enumerate}
\end{exercice}

\begin{exercice}[Étude d'une fonction logarithme]
Soit $f$ la fonction définie par $f(x) = \ln(2x + 4) - x$.
\begin{enumerate}
\item Déterminer l'ensemble de définition $D_f$ de $f$.
\item Calculer les limites de $f$ aux bornes de $D_f$.
\item Calculer $f'(x)$, étudier son signe et dresser le tableau de variations de $f$.
\item En déduire le signe de $f(x)$ sur $D_f$.
\item Résoudre alors l'inéquation $\ln(2x + 4) < x$.
\end{enumerate}
\end{exercice}

\begin{exercice}[Croissance démographique]
La population d'une ville du Cameroun, en millions d'habitants, est modélisée par $P(t) = 5\,\mathrm{e}^{0{,}03t}$, où $t$ est le temps écoulé en années depuis 2020.
\begin{enumerate}
\item Étudier le sens de variation de $P$ sur $[0\ ; +\infty[$.
\item Quelle était la population de la ville en 2020 ?
\item On appelle \emph{temps de doublement} la durée $T$ telle que $P(T) = 2\,P(0)$. Montrer que $T = \dfrac{\ln 2}{0{,}03}$ et donner une valeur approchée de $T$ à une année près.
\item En quelle année la population dépassera-t-elle 10 millions d'habitants ?
\end{enumerate}
\end{exercice}

\courssub{Je me prépare au Bac}

\begin{exercice}[Bac type A : étude complète d'une fonction rationnelle]
On considère la fonction $f$ définie par $f(x) = \dfrac{x^2 - 3x + 3}{x - 2}$ et on note $(\mathcal{C})$ sa courbe représentative dans un repère orthonormé $(O\,; \vec{\imath}, \vec{\jmath})$ (unité : \SI{1}{cm}).
\begin{enumerate}
\item Déterminer l'ensemble de définition $D_f$ de $f$.
\item Calculer les limites de $f$ aux bornes de $D_f$. En déduire que $(\mathcal{C})$ admet une asymptote verticale dont on donnera une équation.
\item Montrer qu'il existe trois réels $a$, $b$ et $c$ tels que, pour tout $x \in D_f$ :
\[ f(x) = ax + b + \dfrac{c}{x - 2}. \]
\item En déduire que la droite $(\Delta)$ d'équation $y = x - 1$ est asymptote oblique à $(\mathcal{C})$ en $+\infty$ et en $-\infty$.
\item Étudier la position relative de $(\mathcal{C})$ par rapport à $(\Delta)$ suivant les valeurs de $x$.
\item Calculer $f'(x)$, étudier son signe et dresser le tableau de variations complet de $f$.
\item Déterminer une équation de la tangente $(T)$ à $(\mathcal{C})$ au point d'abscisse $3$.
\item Tracer soigneusement $(\Delta)$, $(T)$, l'asymptote verticale et $(\mathcal{C})$.
\item Discuter graphiquement, suivant les valeurs du réel $m$, le nombre de solutions de l'équation $f(x) = m$.
\end{enumerate}
\end{exercice}

\begin{exercice}[Bac type A : fonction exponentielle et changement de variable]
Soit $f$ la fonction définie sur $\mathbb{R}$ par $f(x) = \mathrm{e}^{2x} - 2\mathrm{e}^x + 1$, et $(\mathcal{C})$ sa courbe dans un repère orthogonal (unités : \SI{2}{cm} sur l'axe des abscisses, \SI{1}{cm} sur l'axe des ordonnées).
\begin{enumerate}
\item Calculer $\displaystyle\lim_{x \to -\infty} f(x)$. Interpréter graphiquement le résultat.
\item Vérifier que $f(x) = \mathrm{e}^x\left(\mathrm{e}^x - 2\right) + 1$, puis calculer $\displaystyle\lim_{x \to +\infty} f(x)$.
\item Montrer que pour tout réel $x$ : $f'(x) = 2\mathrm{e}^x\left(\mathrm{e}^x - 1\right)$.
\item Étudier le signe de $f'(x)$ et dresser le tableau de variations de $f$.
\item En posant $X = \mathrm{e}^x$, résoudre dans $\mathbb{R}$ l'équation $f(x) = 0$.
\item En déduire le signe de $f(x)$ sur $\mathbb{R}$ et retrouver le résultat de la question 1 de l'exercice « Vrai ou faux » : $f(x) \geq 0$ pour tout $x$.
\item Tracer $(\mathcal{C})$ en plaçant au moins quatre points d'abscisses $-2$, $-1$, $0$ et $1$.
\item Résoudre graphiquement puis algébriquement l'inéquation $f(x) < 1$.
\end{enumerate}
\end{exercice}

\begin{exercice}[Bac type A : problème d'optimisation en entreprise]
Une entreprise agroalimentaire de Bafoussam fabrique et vend chaque mois $x$ tonnes de farine de maïs, avec $x \in [0\ ; 10]$. Le bénéfice mensuel, exprimé en millions de FCFA, est modélisé par :
\[ B(x) = -x^3 + 12x^2 - 36x. \]
\begin{enumerate}
\item Calculer $B(0)$ et $B(10)$. Interpréter ces résultats pour l'entreprise.
\item Calculer $B'(x)$ et vérifier que $B'(x) = -3(x - 2)(x - 6)$.
\item Étudier le signe de $B'(x)$ sur $[0\ ; 10]$ et dresser le tableau de variations de $B$.
\item Déterminer la production mensuelle qui maximise le bénéfice. Quel est alors le bénéfice maximal, en FCFA ?
\item L'entreprise souhaite que son bénéfice soit strictement positif. Montrer que $B(x) = -x(x - 6)^2$ et en déduire l'ensemble des productions pour lesquelles l'entreprise est bénéficiaire.
\item Le directeur financier affirme : « Produire 6 tonnes par mois est la pire décision possible. » A-t-il raison ? Justifier à l'aide du tableau de variations.
\end{enumerate}
\end{exercice}

\newpage

\coursec{Corrigés des exercices}

\begin{corrige}[Exercice 1 — QCM : dérivées et limites]
\begin{enumerate}
\item \textbf{Réponse a).} $f'(x) = 3x^2 - 6x - 9$. En effet, $(x^3)' = 3x^2$, $(-3x^2)' = -6x$, $(-9x)' = -9$ et la dérivée d'une constante est nulle.
\item \textbf{Réponse b).} $g(x) = \ln(2x+4)$ est définie si et seulement si $2x + 4 > 0$, c'est-à-dire $x > -2$. Donc $D_g = ]-2\ ; +\infty[$. L'inégalité est \emph{stricte} car le logarithme n'est pas défini en $0$.
\item \textbf{Réponse c).} Quand $x \to -\infty$, $2x \to -\infty$, et par composition $\displaystyle\lim_{X \to -\infty} \mathrm{e}^X = 0$, donc $\displaystyle\lim_{x \to -\infty} \mathrm{e}^{2x} = 0$.
\item \textbf{Réponse d).} Le dénominateur s'annule en $x = 2$ et le numérateur vaut $4 - 6 + 3 = 1 \neq 0$ en $x = 2$ : la limite est infinie, donc la droite d'équation $x = 2$ est asymptote verticale. (Une asymptote verticale a une équation de la forme $x = a$, jamais $y = a$.)
\end{enumerate}
\end{corrige}

\begin{corrige}[Exercice 2 — Vrai ou faux, en justifiant]
\begin{enumerate}
\item \textbf{Vrai.} En posant $X = \mathrm{e}^x$, on a $\mathrm{e}^{2x} = \left(\mathrm{e}^x\right)^2 = X^2$, donc
\[ f(x) = X^2 - 2X + 1 = (X - 1)^2 = \left(\mathrm{e}^x - 1\right)^2. \]
\item \textbf{Faux.} $f'(x) = 2\mathrm{e}^{2x} - 2\mathrm{e}^x$ (la dérivée de $\mathrm{e}^{2x}$ est $2\mathrm{e}^{2x}$, celle de $-2\mathrm{e}^x$ est $-2\mathrm{e}^x$). L'affirmation oublie le facteur $\mathrm{e}^x$ dans le second terme.
\item \textbf{Vrai.} On factorise : $f(x) = \mathrm{e}^{2x}\left(1 - 2\mathrm{e}^{-x} + \mathrm{e}^{-2x}\right)$. Quand $x \to +\infty$, $\mathrm{e}^{-x} \to 0$ et $\mathrm{e}^{-2x} \to 0$, donc la parenthèse tend vers $1$ et $\mathrm{e}^{2x} \to +\infty$. Par produit, $\displaystyle\lim_{x \to +\infty} f(x) = +\infty$.
\item \textbf{Vrai.} D'après la question 1, $f(x) = 0 \iff \left(\mathrm{e}^x - 1\right)^2 = 0 \iff \mathrm{e}^x = 1 \iff x = 0$. L'équation admet une unique solution : $x = 0$.
\end{enumerate}
\end{corrige}

\begin{corrige}[Exercice 3 — Définitions et cours]
\begin{enumerate}
\item La droite d'équation $x = a$ est \cle{asymptote verticale} à la courbe de $f$ si $\displaystyle\lim_{x \to a^+} f(x) = \pm\infty$ ou $\displaystyle\lim_{x \to a^-} f(x) = \pm\infty$ (au moins l'une des deux limites est infinie).
\item La droite d'équation $y = ax + b$ est \cle{asymptote oblique} à la courbe de $f$ en $+\infty$ si $\displaystyle\lim_{x \to +\infty} \big[f(x) - (ax + b)\big] = 0$.
\item La tangente à la courbe de $f$ au point d'abscisse $a$ a pour équation :
\[ y = f'(a)(x - a) + f(a). \]
\item Limites de référence : $\displaystyle\lim_{x \to 0^+} \ln x = -\infty$ (asymptote verticale $x = 0$ pour le logarithme) et $\displaystyle\lim_{x \to +\infty} \dfrac{\ln x}{x} = 0$ (croissance comparée : toute puissance de $x$ l'emporte sur $\ln x$).
\end{enumerate}
\end{corrige}

\begin{corrige}[Exercice 4 — Calculs directs]
\begin{enumerate}
\item \textbf{Ensembles de définition.}
\begin{itemize}
\item $f(x) = \dfrac{x^2+1}{x-3}$ est définie si et seulement si $x - 3 \neq 0$ : $D_f = \mathbb{R} \setminus \{3\}$.
\item $g(x) = \ln(3x - 6)$ est définie si et seulement si $3x - 6 > 0 \iff x > 2$ : $D_g = ]2\ ; +\infty[$.
\item $h(x) = \mathrm{e}^{-2x+1}$ : l'exponentielle est définie pour tout réel : $D_h = \mathbb{R}$.
\end{itemize}
\item \textbf{Dérivées.}
\begin{itemize}
\item $f = \dfrac{u}{v}$ avec $u = x^2 + 1$, $u' = 2x$, $v = x - 3$, $v' = 1$ :
\[ f'(x) = \dfrac{2x(x-3) - (x^2+1)}{(x-3)^2} = \dfrac{2x^2 - 6x - x^2 - 1}{(x-3)^2} = \dfrac{x^2 - 6x - 1}{(x-3)^2}. \]
\item $g$ est du type $\ln(ax+b)$, de dérivée $\dfrac{a}{ax+b}$ : $g'(x) = \dfrac{3}{3x - 6} = \dfrac{1}{x - 2}$.
\item $h$ est du type $\mathrm{e}^{ax+b}$, de dérivée $a\,\mathrm{e}^{ax+b}$ : $h'(x) = -2\,\mathrm{e}^{-2x+1}$.
\end{itemize}
\item \textbf{Limites du polynôme $P(x) = x^3 - 3x^2 - 9x + 5$.} Un polynôme se comporte à l'infini comme son terme de plus haut degré :
\[ \lim_{x \to +\infty} P(x) = \lim_{x \to +\infty} x^3 = +\infty \;; \qquad \lim_{x \to -\infty} P(x) = \lim_{x \to -\infty} x^3 = -\infty. \]
\end{enumerate}
\end{corrige}

\begin{corrige}[Exercice 5 — Variations d'un polynôme de degré 3]
\begin{enumerate}
\item $f'(x) = 3x^2 - 6x - 9 = 3(x^2 - 2x - 3)$. Le trinôme $x^2 - 2x - 3$ a pour discriminant $\Delta = 4 + 12 = 16$, de racines $x = \dfrac{2 - 4}{2} = -1$ et $x = \dfrac{2 + 4}{2} = 3$. Donc $f'(x) = 3(x+1)(x-3)$, du signe de $(x+1)(x-3)$ : positif à l'extérieur des racines, négatif entre les racines.
\[ f(-1) = -1 - 3 + 9 + 5 = 10 \;; \qquad f(3) = 27 - 27 - 27 + 5 = -22. \]
\begin{center}
\begin{tabular}{|c|ccccccc|}\hline
$x$ & $-\infty$ & & $-1$ & & $3$ & & $+\infty$ \\ \hline
$f'(x)$ & & $+$ & $0$ & $-$ & $0$ & $+$ & \\ \hline
$f$ & $-\infty$ & $\nearrow$ & $10$ & $\searrow$ & $-22$ & $\nearrow$ & $+\infty$ \\ \hline
\end{tabular}
\end{center}
\item $f$ admet un \cle{maximum local} en $x = -1$, égal à $f(-1) = 10$, et un \cle{minimum local} en $x = 3$, égal à $f(3) = -22$. (Ce ne sont pas des extremums globaux puisque $f$ tend vers $\pm\infty$.)
\item Tangente en $0$ : $f(0) = 5$ et $f'(0) = -9$, donc $y = -9(x - 0) + 5$, c'est-à-dire $\boxed{y = -9x + 5}$.
\item \textbf{Lecture graphique et encadrement des racines (TVI).} La courbe coupe l'axe des abscisses en trois points. Justifions et encadrons chaque solution par le théorème des valeurs intermédiaires (TVI) :
\begin{itemize}
\item Sur $[-3\ ; -2]$ : $f(-3) = -27 - 27 + 27 + 5 = -22 < 0$ et $f(-2) = -8 - 12 + 18 + 5 = 3 > 0$. De plus $f(-2{,}5) = -15{,}625 - 18{,}75 + 22{,}5 + 5 = -6{,}875 < 0$. Donc $x_1 \in ]-2{,}5\ ; -2[$.
\item Sur $[0\ ; 1]$ : $f(0) = 5 > 0$ et $f(1) = 1 - 3 - 9 + 5 = -6 < 0$. De plus $f(0{,}5) = 0{,}125 - 0{,}75 - 4{,}5 + 5 = -0{,}125 < 0$. Donc $x_2 \in ]0\ ; 0{,}5[$.
\item Sur $[4\ ; 5]$ : $f(4) = 64 - 48 - 36 + 5 = -15 < 0$ et $f(5) = 125 - 75 - 45 + 5 = 10 > 0$. De plus $f(4{,}5) = 91{,}125 - 60{,}75 - 40{,}5 + 5 = -5{,}125 < 0$. Donc $x_3 \in ]4{,}5\ ; 5[$.
\end{itemize}
$f$ étant continue et strictement monotone sur chacun des trois intervalles, le TVI garantit l'existence et l'unicité de chaque solution.
\end{enumerate}
\end{corrige}

\begin{corrige}[Exercice 6 — Fonction rationnelle et coût moyen]
\begin{enumerate}
\item Pour tout $x > 0$ : $CM(x) = \dfrac{x^2 + 4x + 16}{x} = \dfrac{x^2}{x} + \dfrac{4x}{x} + \dfrac{16}{x} = x + 4 + \dfrac{16}{x}$.
\item $CM'(x) = 1 - \dfrac{16}{x^2} = \dfrac{x^2 - 16}{x^2} = \dfrac{(x-4)(x+4)}{x^2}$. Sur $]0\ ; +\infty[$, $x^2 > 0$ et $x + 4 > 0$, donc $CM'(x)$ est du signe de $x - 4$ : négatif sur $]0\ ; 4[$, nul en $4$, positif sur $]4\ ; +\infty[$.
\item Limites aux bornes : $\displaystyle\lim_{x \to 0^+} CM(x) = +\infty$ (car $\dfrac{16}{x} \to +\infty$) et $\displaystyle\lim_{x \to +\infty} CM(x) = +\infty$. De plus $CM(4) = 4 + 4 + 4 = 12$.
\begin{center}
\begin{tabular}{|c|ccccc|}\hline
$x$ & $0$ & & $4$ & & $+\infty$ \\ \hline
$CM'(x)$ & & $-$ & $0$ & $+$ & \\ \hline
$CM$ & $+\infty$ & $\searrow$ & $12$ & $\nearrow$ & $+\infty$ \\ \hline
\end{tabular}
\end{center}
\item Le coût moyen est minimal pour $x = 4$ meubles par semaine. Le coût moyen minimal est $CM(4) = 12$ dizaines de milliers de FCFA, soit $\boxed{120\,000\ \text{FCFA}}$ par meuble.
\end{enumerate}
\end{corrige}

\begin{corrige}[Exercice 7 — Étude d'une fonction logarithme]
\begin{enumerate}
\item $f(x) = \ln(2x+4) - x$ est définie si et seulement si $2x + 4 > 0 \iff x > -2$ : $D_f = ]-2\ ; +\infty[$.
\item \textbf{En $-2^+$ :} $2x + 4 \to 0^+$, donc $\ln(2x+4) \to -\infty$ et $-x \to 2$ : $\displaystyle\lim_{x \to -2^+} f(x) = -\infty$ (asymptote verticale $x = -2$).

\textbf{En $+\infty$ :} forme indéterminée « $+\infty - \infty$ ». On factorise par $x$ :
\[ f(x) = x\left(\dfrac{\ln(2x+4)}{x} - 1\right). \]
Or $\dfrac{\ln(2x+4)}{x} = \dfrac{\ln(2x+4)}{2x+4} \times \dfrac{2x+4}{x} \to 0 \times 2 = 0$ par croissance comparée. Donc la parenthèse tend vers $-1$ et $\displaystyle\lim_{x \to +\infty} f(x) = -\infty$.
\item $f'(x) = \dfrac{2}{2x+4} - 1 = \dfrac{2 - (2x+4)}{2x+4} = \dfrac{-2x - 2}{2x+4} = \dfrac{-(x+1)}{x+2}$.

Sur $D_f$, $x + 2 > 0$, donc $f'(x)$ est du signe de $-(x+1)$ : positif sur $]-2\ ; -1[$, nul en $-1$, négatif sur $]-1\ ; +\infty[$. Maximum : $f(-1) = \ln 2 + 1 \approx 1{,}69$.
\begin{center}
\begin{tabular}{|c|ccccc|}\hline
$x$ & $-2$ & & $-1$ & & $+\infty$ \\ \hline
$f'(x)$ & & $+$ & $0$ & $-$ & \\ \hline
$f$ & $-\infty$ & $\nearrow$ & $1 + \ln 2$ & $\searrow$ & $-\infty$ \\ \hline
\end{tabular}
\end{center}
\item Le maximum $1 + \ln 2$ est strictement positif et les limites aux bornes sont $-\infty$. $f$ est continue et strictement monotone sur $]-2\ ; -1]$ et sur $[-1\ ; +\infty[$ : par le corollaire du TVI, l'équation $f(x) = 0$ admet exactement deux solutions $\alpha \in ]-2\ ; -1[$ et $\beta \in ]-1\ ; +\infty[$. Par balayage : $f(-1{,}92) \approx 0{,}09 > 0$ et $f(-1{,}93) \approx -0{,}04 < 0$, donc $\alpha \approx -1{,}92$ ; $f(2{,}10) \approx 0{,}004 > 0$ et $f(2{,}11) \approx -0{,}003 < 0$, donc $\beta \approx 2{,}10$.

Conclusion : $f(x) < 0$ sur $]-2\ ; \alpha[ \cup ]\beta\ ; +\infty[$, $f(x) > 0$ sur $]\alpha\ ; \beta[$.
\item $\ln(2x+4) < x \iff f(x) < 0 \iff x \in ]-2\ ; \alpha[ \cup ]\beta\ ; +\infty[$, avec $\alpha \approx -1{,}92$ et $\beta \approx 2{,}10$.
\end{enumerate}
\end{corrige}

\begin{corrige}[Exercice 8 — Croissance démographique]
\begin{enumerate}
\item $P'(t) = 5 \times 0{,}03\,\mathrm{e}^{0{,}03t} = 0{,}15\,\mathrm{e}^{0{,}03t} > 0$ pour tout $t \geq 0$ : $P$ est \cle{strictement croissante} sur $[0\ ; +\infty[$.
\item $P(0) = 5\,\mathrm{e}^0 = 5$ : la ville comptait \textbf{5 millions d'habitants en 2020}.
\item $P(T) = 2P(0) \iff 5\,\mathrm{e}^{0{,}03T} = 10 \iff \mathrm{e}^{0{,}03T} = 2 \iff 0{,}03\,T = \ln 2 \iff T = \dfrac{\ln 2}{0{,}03}$.

Numériquement : $T \approx \dfrac{0{,}693}{0{,}03} \approx 23{,}1$, soit environ \textbf{23 ans}.
\item $P(t) > 10 \iff \mathrm{e}^{0{,}03t} > 2 \iff t > T \approx 23{,}1$. La population dépassera 10 millions d'habitants à partir de $2020 + 23{,}1 \approx 2043{,}1$, c'est-à-dire \textbf{au cours de l'année 2043}.
\end{enumerate}
\end{corrige}

\begin{corrige}[Exercice 9 — Bac type A : étude complète d'une fonction rationnelle]
\emph{Barème indicatif (sur 10 pts) : question 1 : 0,5 ; question 2 : 1,5 ; question 3 : 1 ; question 4 : 1 ; question 5 : 1 ; question 6 : 2 ; question 7 : 1 ; question 8 : 1,5 ; question 9 : 0,5.}

\begin{enumerate}
\item $f$ est définie si et seulement si $x - 2 \neq 0$ : $D_f = \mathbb{R} \setminus \{2\} = ]-\infty\ ; 2[ \cup ]2\ ; +\infty[$.
\item \textbf{En $2$ :} le numérateur vaut $2^2 - 3 \times 2 + 3 = 1 > 0$ en $x = 2$.
\begin{itemize}
\item Si $x \to 2^+$, $x - 2 \to 0^+$, donc $\displaystyle\lim_{x \to 2^+} f(x) = +\infty$ ;
\item si $x \to 2^-$, $x - 2 \to 0^-$, donc $\displaystyle\lim_{x \to 2^-} f(x) = -\infty$.
\end{itemize}
La limite étant infinie, la droite d'équation $\boxed{x = 2}$ est \cle{asymptote verticale} à $(\mathcal{C})$.

\textbf{À l'infini :} $f$ se comporte comme $\dfrac{x^2}{x} = x$ : $\displaystyle\lim_{x \to +\infty} f(x) = +\infty$ et $\displaystyle\lim_{x \to -\infty} f(x) = -\infty$.
\item Effectuons la division euclidienne de $x^2 - 3x + 3$ par $x - 2$ : $x^2 - 3x + 3 = (x - 2)(x - 1) + 1$. Vérification : $(x-2)(x-1) + 1 = x^2 - 3x + 2 + 1 = x^2 - 3x + 3$. Donc pour tout $x \in D_f$ :
\[ f(x) = x - 1 + \dfrac{1}{x - 2} \qquad (a = 1,\ b = -1,\ c = 1). \]
\item $f(x) - (x - 1) = \dfrac{1}{x-2}$, et $\displaystyle\lim_{x \to \pm\infty} \dfrac{1}{x-2} = 0$. Donc la droite $(\Delta)$ d'équation $y = x - 1$ est \cle{asymptote oblique} à $(\mathcal{C})$ en $+\infty$ et en $-\infty$.
\item Le signe de $f(x) - (x-1) = \dfrac{1}{x-2}$ est celui de $x - 2$ :
\begin{itemize}
\item si $x > 2$, $f(x) - (x-1) > 0$ : $(\mathcal{C})$ est \textbf{au-dessus} de $(\Delta)$ ;
\item si $x < 2$, $f(x) - (x-1) < 0$ : $(\mathcal{C})$ est \textbf{en dessous} de $(\Delta)$.
\end{itemize}
\item Avec $f = \dfrac{u}{v}$, $u = x^2 - 3x + 3$, $u' = 2x - 3$, $v = x - 2$, $v' = 1$ :
\[ f'(x) = \dfrac{(2x-3)(x-2) - (x^2 - 3x + 3)}{(x-2)^2} = \dfrac{2x^2 - 7x + 6 - x^2 + 3x - 3}{(x-2)^2} = \dfrac{x^2 - 4x + 3}{(x-2)^2} = \dfrac{(x-1)(x-3)}{(x-2)^2}. \]
Le dénominateur est strictement positif sur $D_f$ ; $f'(x)$ est du signe de $(x-1)(x-3)$ : positif sur $]-\infty\ ; 1[$ et $]3\ ; +\infty[$, négatif sur $]1\ ; 2[$ et $]2\ ; 3[$.

$f(1) = \dfrac{1 - 3 + 3}{-1} = -1$ (maximum local) ; $f(3) = \dfrac{9 - 9 + 3}{1} = 3$ (minimum local).
\begin{center}
\begin{tabular}{|c|c|c|c|c||c|c|c|c|cc|}\hline
$x$ & $-\infty$ & & $1$ & & $2^-$ & $2^+$ & & $3$ & & $+\infty$ \\ \hline
$f'(x)$ & & $+$ & $0$ & $-$ & & & $-$ & $0$ & $+$ & \\ \hline
$f$ & $-\infty$ & $\nearrow$ & $-1$ & $\searrow$ & $-\infty$ & $+\infty$ & $\searrow$ & $3$ & $\nearrow$ & $+\infty$ \\ \hline
\end{tabular}
\end{center}
\item $f(3) = 3$ et $f'(3) = 0$ : la tangente $(T)$ en $3$ est horizontale, d'équation $\boxed{y = 3}$.
\item Tracé : on place l'asymptote verticale $x = 2$, l'asymptote oblique $y = x - 1$ (passant par $(0\,; -1)$ et $(2\,; 1)$), la tangente horizontale $y = 3$, le sommet $(1\,; -1)$ de la branche de gauche et le point $(3\,; 3)$ de la branche de droite. Quelques points : $f(0) = -\dfrac{3}{2}$, $f(4) = \dfrac{16 - 12 + 3}{2} = \dfrac{7}{2}$, $f(5) = \dfrac{25 - 15 + 3}{3} = \dfrac{13}{3} \approx 4{,}3$.
\item L'équation $f(x) = m$ revient à chercher les intersections de $(\mathcal{C})$ avec la droite horizontale $y = m$. D'après le tableau de variations :
\begin{itemize}
\item si $m < -1$ : \textbf{2 solutions} (la droite coupe la branche de gauche deux fois) ;
\item si $m = -1$ : \textbf{1 solution} ($x = 1$) ;
\item si $-1 < m < 3$ : \textbf{aucune solution} ;
\item si $m = 3$ : \textbf{1 solution} ($x = 3$) ;
\item si $m > 3$ : \textbf{2 solutions} (la droite coupe la branche de droite deux fois).
\end{itemize}
\end{enumerate}

\textbf{Points de vigilance.} Citer la valeur du numérateur en $2$ (non nulle) avant de conclure à l'asymptote verticale ; distinguer limite à droite et limite à gauche ; pour l'asymptote oblique, écrire explicitement $\lim [f(x) - (x-1)] = 0$ ; ne jamais faire « traverser » la courbe à la valeur interdite $x = 2$ dans le tableau (utiliser la double barre $||$) ; pour la discussion graphique, justifier par le tableau de variations et non par une simple affirmation.
\end{corrige}

\begin{corrige}[Exercice 10 — Bac type A : fonction exponentielle et changement de variable]
\emph{Barème indicatif (sur 10 pts) : question 1 : 1 ; question 2 : 1,5 ; question 3 : 1 ; question 4 : 1,5 ; question 5 : 1,5 ; question 6 : 1 ; question 7 : 1,5 ; question 8 : 1.}

\begin{enumerate}
\item Quand $x \to -\infty$ : $\mathrm{e}^{2x} \to 0$ et $\mathrm{e}^x \to 0$, donc $\displaystyle\lim_{x \to -\infty} f(x) = 0 - 0 + 1 = 1$.

\emph{Interprétation graphique :} la droite d'équation $\boxed{y = 1}$ est \cle{asymptote horizontale} à $(\mathcal{C})$ en $-\infty$.
\item $\mathrm{e}^x\left(\mathrm{e}^x - 2\right) + 1 = \mathrm{e}^{2x} - 2\mathrm{e}^x + 1 = f(x)$. Quand $x \to +\infty$ : $\mathrm{e}^x \to +\infty$ et $\mathrm{e}^x - 2 \to +\infty$, donc par produit $\mathrm{e}^x(\mathrm{e}^x - 2) \to +\infty$, puis $\displaystyle\lim_{x \to +\infty} f(x) = +\infty$.
\item $f'(x) = 2\mathrm{e}^{2x} - 2\mathrm{e}^x = 2\mathrm{e}^x \cdot \mathrm{e}^x - 2\mathrm{e}^x = 2\mathrm{e}^x\left(\mathrm{e}^x - 1\right)$.
\item Pour tout réel $x$, $2\mathrm{e}^x > 0$, donc $f'(x)$ est du signe de $\mathrm{e}^x - 1$ :
\[ \mathrm{e}^x - 1 > 0 \iff \mathrm{e}^x > \mathrm{e}^0 \iff x > 0. \]
$f$ est donc strictement décroissante sur $]-\infty\ ; 0]$, strictement croissante sur $[0\ ; +\infty[$, avec un minimum $f(0) = 1 - 2 + 1 = 0$.
\begin{center}
\begin{tabular}{|c|ccccc|}\hline
$x$ & $-\infty$ & & $0$ & & $+\infty$ \\ \hline
$f'(x)$ & & $-$ & $0$ & $+$ & \\ \hline
$f$ & $1$ & $\searrow$ & $0$ & $\nearrow$ & $+\infty$ \\ \hline
\end{tabular}
\end{center}
\item Posons $X = \mathrm{e}^x$ ($X > 0$). Alors $f(x) = 0 \iff X^2 - 2X + 1 = 0 \iff (X - 1)^2 = 0 \iff X = 1$, soit $\mathrm{e}^x = 1 \iff x = 0$. L'équation admet pour unique solution $\boxed{x = 0}$.
\item $f(x) = \left(\mathrm{e}^x - 1\right)^2$ est un carré, donc $f(x) \geq 0$ pour tout réel $x$, et $f(x) = 0$ uniquement en $x = 0$. On retrouve ainsi le résultat de l'exercice « Vrai ou faux » : $f$ est positive sur $\mathbb{R}$, ce qui est cohérent avec le minimum $0$ atteint en $0$.
\item Valeurs approchées :
\begin{center}
\begin{tabularx}{\linewidth}{|l|X|X|X|X|}\hline
\rowcolor{popblueL} $x$ & $-2$ & $-1$ & $0$ & $1$ \\ \hline
$f(x)$ & $0{,}75$ & $0{,}40$ & $0$ & $2{,}95$ \\ \hline
\end{tabularx}
\end{center}
En effet : $f(-2) = \mathrm{e}^{-4} - 2\mathrm{e}^{-2} + 1 \approx 0{,}018 - 0{,}271 + 1 \approx 0{,}75$ ; $f(-1) = \mathrm{e}^{-2} - 2\mathrm{e}^{-1} + 1 \approx 0{,}135 - 0{,}736 + 1 \approx 0{,}40$ ; $f(1) = \mathrm{e}^2 - 2\mathrm{e} + 1 \approx 7{,}389 - 5{,}437 + 1 \approx 2{,}95$. On trace la courbe en la faisant approcher l'asymptote $y = 1$ vers la gauche et en la faisant passer par l'origine avec tangente horizontale.
\item \textbf{Graphiquement :} $f(x) < 1$ là où $(\mathcal{C})$ est strictement sous la droite $y = 1$ : on lit $x < \ln 2 \approx 0{,}69$.

\textbf{Algébriquement :}
\[ f(x) < 1 \iff \left(\mathrm{e}^x - 1\right)^2 < 1 \iff -1 < \mathrm{e}^x - 1 < 1 \iff 0 < \mathrm{e}^x < 2. \]
La condition $0 < \mathrm{e}^x$ est toujours vraie ; $\mathrm{e}^x < 2 \iff x < \ln 2$. Solution : $\boxed{\mathcal{S} = ]-\infty\ ; \ln 2[}$.
\end{enumerate}

\textbf{Points de vigilance.} Pour la limite en $+\infty$, factoriser par $\mathrm{e}^x$ (ou $\mathrm{e}^{2x}$) pour lever l'indétermination — ne jamais écrire « $+\infty - \infty = \ldots$ » ; la dérivée de $\mathrm{e}^{2x}$ est $2\mathrm{e}^{2x}$ (ne pas oublier la dérivée de l'exposant) ; au changement de variable, préciser $X = \mathrm{e}^x > 0$ ; pour l'inéquation avec un carré, revenir à la double inégalité $-1 < \mathrm{e}^x - 1 < 1$ plutôt que d'extraire une racine carrée sans précaution.
\end{corrige}

\begin{corrige}[Exercice 11 — Bac type A : problème d'optimisation en entreprise]
\emph{Barème indicatif (sur 10 pts) : question 1 : 1 ; question 2 : 1,5 ; question 3 : 2 ; question 4 : 1,5 ; question 5 : 2 ; question 6 : 2.}

\begin{enumerate}
\item $B(0) = 0$ : sans production, l'entreprise ne réalise aucun bénéfice.

$B(10) = -1000 + 12 \times 100 - 360 = -1000 + 1200 - 360 = -160$ : en produisant 10 tonnes, l'entreprise perd 160 millions de FCFA par mois (les coûts fixes et variables dépassent largement les recettes à ce niveau de production).
\item $B'(x) = -3x^2 + 24x - 36 = -3(x^2 - 8x + 12)$. Or $(x-2)(x-6) = x^2 - 8x + 12$, donc $B'(x) = -3(x-2)(x-6)$.
\item Sur $[0\ ; 10]$, étudions le signe de $-3(x-2)(x-6)$. Le produit $(x-2)(x-6)$ est positif à l'extérieur de $[2\ ; 6]$ et négatif entre $2$ et $6$ ; avec le facteur $-3 < 0$ :
\begin{itemize}
\item $B'(x) < 0$ sur $[0\ ; 2[$ et sur $]6\ ; 10]$ ;
\item $B'(x) > 0$ sur $]2\ ; 6[$ ;
\item $B'(2) = B'(6) = 0$.
\end{itemize}
$B(2) = -8 + 48 - 72 = -32$ ; $B(6) = -216 + 432 - 216 = 0$.
\begin{center}
\begin{tabular}{|c|ccccccc|}\hline
$x$ & $0$ & & $2$ & & $6$ & & $10$ \\ \hline
$B'(x)$ & & $-$ & $0$ & $+$ & $0$ & $-$ & \\ \hline
$B$ & $0$ & $\searrow$ & $-32$ & $\nearrow$ & $0$ & $\searrow$ & $-160$ \\ \hline
\end{tabular}
\end{center}
\item D'après le tableau, le maximum de $B$ sur $[0\ ; 10]$ est $0$, atteint en $x = 6$ (et en $x = 0$, qui correspond à l'absence de production). La meilleure décision de production est donc $\boxed{x = 6\ \text{tonnes par mois}}$, pour un bénéfice maximal de $\boxed{0\ \text{FCFA}}$ : l'entreprise couvre exactement ses coûts, sans dégager de gain.
\item Développons : $-x(x-6)^2 = -x(x^2 - 12x + 36) = -x^3 + 12x^2 - 36x = B(x)$.

Pour $x \in [0\ ; 10]$ : $-x \leq 0$ et $(x-6)^2 \geq 0$, donc $B(x) = -x(x-6)^2 \leq 0$ pour \emph{toute} production. Le bénéfice n'est \textbf{jamais strictement positif} : l'ensemble des productions pour lesquelles l'entreprise est bénéficiaire est \textbf{vide}. Le mieux qu'elle puisse faire est d'équilibrer ses comptes ($B = 0$ en $x = 6$).
\item \textbf{Non, le directeur a tort.} En $x = 6$, le bénéfice vaut $0$, ce qui est le \emph{maximum} de $B$ sur $[0\ ; 10]$ : c'est au contraire la meilleure production possible. La pire décision est de produire $10$ tonnes, avec une perte de $160$ millions de FCFA. Le directeur a sans doute confondu le point où $B'$ s'annule en « montant » ($x = 6$, maximum) avec un minimum ; or le minimum local est en $x = 2$ ($B(2) = -32$) et le minimum global en $x = 10$.
\end{enumerate}

\textbf{Points de vigilance.} Dans un problème concret, toujours respecter l'ensemble d'étude imposé ($[0\ ; 10]$ ici) et comparer les valeurs aux bornes avec les extremums locaux avant de conclure sur un maximum \emph{global} ; conclure avec une phrase en langage économique (quantité, unité monétaire) ; pour la question 6, la réponse doit s'appuyer explicitement sur le tableau de variations.
\end{corrige}

\begin{methode}[Méthode générale d'étude d'une fonction]
\emph{Applicable à tous les types de fonctions du programme (polynômes, rationnelles, logarithmes, exponentielles).}
\begin{enumerate}
\item \textbf{Ensemble de définition $D_f$} : identifier les valeurs interdites (dénominateur nul, argument du $\ln \le 0$, etc.).
\item \textbf{Parité / Périodicité} (si pertinent) : simplifie l'étude à un demi-plan.
\item \textbf{Limites aux bornes de $D_f$} : 
\begin{itemize}
\item Bornes finies interdites $\to$ asymptotes verticales ?
\item $\pm\infty$ (si dans $D_f$) $\to$ asymptotes horizontales ou obliques ?
\end{itemize}
\item \textbf{Dérivée $f'$} : calculer soigneusement (règles de dérivation, dérivation composée). Simplifier l'expression (factorisation) pour en étudier le signe.
\item \textbf{Tableau de variations} : dresser le tableau complet (valeurs interdites séparées par $||$, signes de $f'$, flèches de $f$, valeurs remarquables $f(x)$ aux extrema et bornes).
\item \textbf{Asymptotes} : 
\begin{itemize}
\item Verticales : $x = a$ si limite infinie en $a$.
\item Obliques/Horizontales : $y = ax+b$ si $\lim_{x\to\pm\infty}[f(x)-(ax+b)]=0$. Préciser la position relative (au-dessus / en dessous).
\end{itemize}
\item \textbf{Points particuliers} : intersections avec les axes ($f(0)$, $f(x)=0$), tangentes remarquables.
\item \textbf{Tracé de la courbe} : repère orthonormé, asymptotes (pointillés), points caractéristiques, courbe lisse respectant les variations et la position par rapport aux asymptotes.
\item \textbf{Résolution d'équations / inéquations} : lecture graphique (nombre de solutions) ou algébrique (factorisation, changement de variable), encadrements par TVI si nécessaire.
\end{enumerate}
\end{methode}

\begin{aretenir}[Comparaison des croissances (références pour les limites)]
\begin{itemize}
\item \textbf{Puissances vs Logarithme :} $\forall \alpha > 0,\ \displaystyle\lim_{x\to+\infty} \frac{\ln x}{x^\alpha} = 0$. Le logarithme est négligeable devant toute puissance de $x$.
\item \textbf{Exponentielle vs Puissances :} $\forall n \in \mathbb{N},\ \displaystyle\lim_{x\to+\infty} \frac{x^n}{\mathrm{e}^x} = 0$. L'exponentielle l'emporte sur toute puissance polynomiale.
\item \textbf{Exponentielle vs Logarithme :} $\displaystyle\lim_{x\to+\infty} \frac{\ln x}{\mathrm{e}^x} = 0$.
\item \textbf{En $-\infty$ pour l'exponentielle :} $\displaystyle\lim_{x\to-\infty} \mathrm{e}^x = 0$ et $\displaystyle\lim_{x\to-\infty} x^n \mathrm{e}^x = 0$ (pour tout $n$).
\item \textbf{Formes indéterminées classiques à lever :} 
$+\infty - \infty$ (factorisation par le terme dominant), 
$\frac{\infty}{\infty}$ (factorisation ou croissance comparée), 
$0 \times \infty$ (réécriture en quotient).
\end{itemize}
\end{aretenir}

\newpage

\coursec{Fiche bilan}

\coursec{Étude de fonctions}

\courssub{Étude des fonctions polynômes de degré $\leq 3$}

\begin{definition}[Fonction polynôme de degré $\leq 3$]
Une fonction $f$ est un \cle{polynôme de degré $\leq 3$} si elle s'écrit sous la forme :
\[ f(x) = ax^3 + bx^2 + cx + d \quad (a,b,c,d \in \mathbb{R}) \]
\begin{itemize}
  \item Si $a \neq 0$, le degré est $3$.
  \item Si $a = 0$ et $b \neq 0$, le degré est $2$ (trinôme).
  \item Si $a = b = 0$ et $c \neq 0$, le degré est $1$ (affine).
  \item Si $a = b = c = 0$, c'est une fonction constante.
\end{itemize}
L'\cle{ensemble de définition} est $D_f = \mathbb{R}$.
\end{definition}

\begin{propriete}[Limites en $\pm\infty$]
Pour $f(x) = ax^3 + bx^2 + cx + d$ avec $a \neq 0$ :
\[ \lim_{x \to +\infty} f(x) = \operatorname{signe}(a) \cdot (+\infty), \qquad \lim_{x \to -\infty} f(x) = \operatorname{signe}(a) \cdot (-\infty) \]
Pour un trinôme $ax^2+bx+c$ ($a \neq 0$), les deux limites sont $+\infty$ si $a>0$ et $-\infty$ si $a<0$.
\end{propriete}

\begin{propriete}[Dérivée et étude du signe]
La dérivée sur $\mathbb{R}$ est :
\[ f'(x) = 3ax^2 + 2bx + c \]
C'est un trinôme. Son discriminant réduit est $\Delta' = b^2 - 3ac$.
\begin{itemize}
  \item Si $\Delta' < 0$ : $f'$ a le signe de $3a$ sur $\mathbb{R}$ ; $f$ est strictement monotone.
  \item Si $\Delta' = 0$ : $f'$ a le signe de $3a$ sur $\mathbb{R}$ et s'annule en $x_0 = -\frac{b}{3a}$ ; $f$ est strictement monotone (point d'inflexion à tangente horizontale).
  \item Si $\Delta' > 0$ : $f'$ a deux zéros $x_1 < x_2$ donnés par $\frac{-b \pm \sqrt{\Delta'}}{3a}$. Le signe de $f'$ est celui de $3a$ en dehors de $[x_1, x_2]$ et opposé à l'intérieur. $f$ admet un maximum local en $x_1$ et un minimum local en $x_2$ (si $a>0$), ou l'inverse (si $a<0$).
\end{itemize}
\end{propriete}

\begin{methode}[Étude complète d'un polynôme degré 3 : $f(x)=ax^3+bx^2+cx+d$]
\begin{enumerate}
  \item $D_f = \mathbb{R}$.
  \item Limites en $-\infty$ et $+\infty$ (signe du terme de plus haut degré).
  \item Calculer $f'(x) = 3ax^2+2bx+c$, $\Delta' = b^2-3ac$, racines éventuelles.
  \item Tableau de signes de $f'$ puis tableau de variations de $f$.
  \item Calculer les extremums locaux (images des racines de $f'$) et $f(0)=d$ (intersection avec l'axe des ordonnées). Résoudre $f(x)=0$ si possible (intersections avec l'axe des abscisses).
  \item Tracer la courbe $\mathcal{C}_f$ dans un repère orthonormé.
\end{enumerate}
\end{methode}

\begin{exemplebox}[Étude de $f(x) = x^3 - 3x + 2$]
$D_f = \mathbb{R}$.
$\lim_{x\to-\infty} f(x) = -\infty$, $\lim_{x\to+\infty} f(x) = +\infty$.
$f'(x) = 3x^2 - 3 = 3(x-1)(x+1)$.
Zéros de $f'$ : $-1$ et $1$.
Signe de $f'$ : $+$ sur $]-\infty, -1]$, $-$ sur $[-1, 1]$, $+$ sur $[1, +\infty[$.
\begin{center}
\begin{tabular}{|c|ccccccc|}\hline
$x$ & $-\infty$ & & $-1$ & & $1$ & & $+\infty$ \\ \hline
$f'(x)$ & & $+$ & $0$ & $-$ & $0$ & $+$ & \\ \hline
$f(x)$ & $-\infty$ & $\nearrow$ & $4$ & $\searrow$ & $0$ & $\nearrow$ & $+\infty$ \\ \hline
\end{tabular}
\end{center}
Maximum $4$ en $-1$, minimum $0$ en $1$.
$f(0)=2$. $f(x)=(x-1)^2(x+2)$ : racines $-2$ (simple) et $1$ (double, tangente à l'axe).
\tcblower
\begin{popfigure}
\begin{tikzpicture}[scale=0.7, domain=-2.5:2]
\draw[->] (-3,0) -- (3,0) node[right] {$x$};
\draw[->] (0,-1) -- (0,5) node[above] {$y$};
\foreach \x in {-2,-1,1,2} \draw (\x,0.1) -- (\x,-0.1) node[below] {\small $\x$};
\foreach \y in {1,2,3,4} \draw (0.1,\y) -- (-0.1,\y) node[left] {\small $\y$};
\draw[popcurve, thick, domain=-2.5:2, samples=200] plot (\x, {\x*\x*\x - 3*\x + 2});
\draw[dashed, popmuted] (-1,0) node[below] {\small $-1$} -- (-1,4) node[left] {\small $4$};
\draw[dashed, popmuted] (1,0) node[below] {\small $1$} -- (1,0);
\draw[dashed, popmuted] (0,2) node[left] {\small $2$} -- (0,0);
\draw[dashed, popmuted] (-2,0) node[below] {\small $-2$} -- (-2,0);
\end{tikzpicture}
\legende{Courbe de $f(x)=x^3-3x+2$.}
\end{popfigure}
\end{exemplebox}

\begin{attention}[Pièges classiques]
\begin{itemize}
  \item Oublier que $D_f = \mathbb{R}$ (pas de restriction).
  \item Confondre le discriminant $\Delta = b^2-4ac$ du trinôme $f(x)$ et le discriminant réduit $\Delta' = b^2-3ac$ de $f'(x)$.
  \item Ne pas justifier le sens de variation par le signe de $f'$ (et non par "la courbe monte").
  \item Pour un trinôme ($a=0$), la dérivée est affine : une seule racine, pas de $\Delta'$.
\end{itemize}
\end{attention}

\courssub{Étude des fonctions rationnelles du type $\dfrac{ax^2+bx+c}{dx+e}$}

\begin{definition}[Fonction rationnelle du second degré sur le premier]
Soient $a,b,c,d,e \in \mathbb{R}$ avec $d \neq 0$ (sinon c'est un polynôme). La fonction
\[ f(x) = \frac{ax^2+bx+c}{dx+e} \]
est définie sur $D_f = \mathbb{R} \setminus \left\{ -\frac{e}{d} \right\}$.
La valeur $x_0 = -\frac{e}{d}$ est une \cle{valeur interdite} (zéro du dénominateur).
\end{definition}

\begin{propriete}[Asymptote verticale]
La droite d'équation $x = -\frac{e}{d}$ est une \cle{asymptote verticale} (A.V.).
Les limites sont infinies :
\[ \lim_{x \to \left(-\frac{e}{d}\right)^\pm} f(x) = \pm\infty \quad \text{(signes à déterminer selon le numérateur en $x_0$)} \]
Plus précisément, si $N(x_0) = a x_0^2 + b x_0 + c \neq 0$ :
\begin{itemize}
  \item $x \to x_0^-$ : signe de $\frac{N(x_0)}{d(x-x_0)}$.
  \item $x \to x_0^+$ : signe de $\frac{N(x_0)}{d(x-x_0)}$.
\end{itemize}
Si $N(x_0)=0$, il y a simplification (trou) : ce n'est plus une asymptote verticale.
\end{propriete}

\begin{propriete}[Asymptote oblique]
On effectue la division euclidienne du numérateur par le dénominateur :
\[ ax^2+bx+c = \left( \frac{a}{d}x + \frac{bd-ae}{d^2} \right)(dx+e) + \frac{cd^2 - bde + ae^2}{d^2} \]
On pose $\alpha = \frac{a}{d}$, $\beta = \frac{bd-ae}{d^2}$, $\gamma = \frac{cd^2 - bde + ae^2}{d^2}$.
Alors $f(x) = \alpha x + \beta + \frac{\gamma}{dx+e}$.
La droite d'équation $y = \alpha x + \beta$ est l'\cle{asymptote oblique} (A.O.).
La position relative de la courbe par rapport à l'asymptote est donnée par le signe de $\frac{\gamma}{dx+e}$ :
\begin{itemize}
  \item Si $\frac{\gamma}{dx+e} > 0$, la courbe est \textbf{au-dessus} de l'asymptote.
  \item Si $\frac{\gamma}{dx+e} < 0$, la courbe est \textbf{en-dessous} de l'asymptote.
\end{itemize}
\end{propriete}

\begin{propriete}[Dérivée]
\[ f'(x) = \frac{(2ax+b)(dx+e) - d(ax^2+bx+c)}{(dx+e)^2} = \frac{adx^2 + 2aex + (be-cd)}{(dx+e)^2} \]
Le dénominateur est toujours $>0$ sur $D_f$. Le signe de $f'$ est celui du numérateur $N'(x) = adx^2 + 2aex + (be-cd)$.
Son discriminant est $\Delta_{f'} = 4a^2e^2 - 4ad(be-cd) = 4a(ae^2 - bde + cd^2) = 4ad^2 \cdot \frac{\gamma}{d}$ (lien avec le reste de la division).
\end{propriete}

\begin{methode}[Étude complète d'une fonction rationnelle $\frac{ax^2+bx+c}{dx+e}$]
\begin{enumerate}
  \item $D_f = \mathbb{R} \setminus \{ -e/d \}$.
  \item Limites aux bornes de $D_f$ : $\pm\infty$ (A.V.) et $\pm\infty$ (comportement en $\pm\infty$ pour A.O.).
  \item Division euclidienne pour trouver A.O. $y=\alpha x+\beta$ et position relative (signe de $\gamma/(dx+e)$).
  \item Calculer $f'$, étudier son signe (trinôme $N'(x)$), dresser tableau de variations sur chaque intervalle de $D_f$.
  \item Points remarquables : intersections axes ($f(0)$, $f(x)=0$), extremums.
  \item Tracer : asymptotes (pointillés), courbe respectant variations et position relative à l'A.O.
\end{enumerate}
\end{methode}

\begin{exemplebox}[Étude de $f(x) = \frac{x^2+1}{x-1}$]
$D_f = \mathbb{R} \setminus \{1\}$.
A.V. $x=1$. $\lim_{x\to 1^-} f(x) = -\infty$, $\lim_{x\to 1^+} f(x) = +\infty$.
Division : $x^2+1 = (x+1)(x-1) + 2$. Donc $f(x) = x+1 + \frac{2}{x-1}$.
A.O. $y = x+1$. Position : $\frac{2}{x-1} > 0 \iff x>1$ (au-dessus), $<0 \iff x<1$ (en-dessous).
$f'(x) = \frac{(2x)(x-1) - (x^2+1)}{(x-1)^2} = \frac{x^2 - 2x - 1}{(x-1)^2}$.
Zéros de $f'$ : $1-\sqrt{2} \approx -0,4$ et $1+\sqrt{2} \approx 2,4$.
Variations sur $]-\infty, 1[$ : $\nearrow$ jusqu'à $1-\sqrt{2}$ (max), puis $\searrow$ vers $-\infty$.
Variations sur $]1, +\infty[$ : $\searrow$ jusqu'à $1+\sqrt{2}$ (min), puis $\nearrow$ vers $+\infty$.
\tcblower
\begin{popfigure}
\begin{tikzpicture}[scale=0.6, domain=-4:5]
\draw[->] (-5,0) -- (6,0) node[right] {$x$};
\draw[->] (0,-5) -- (0,6) node[above] {$y$};
\foreach \x in {-4,-3,-2,-1,1,2,3,4,5} \draw (\x,0.1) -- (\x,-0.1) node[below] {\small $\x$};
\foreach \y in {-4,-3,-2,-1,1,2,3,4,5} \draw (0.1,\y) -- (-0.1,\y) node[left] {\small $\y$};
% Asymptotes
\draw[dashed, poporange] (1,-5) -- (1,6) node[above] {\small $x=1$};
\draw[dashed, poporange] (-5,-4) -- (6,7) node[right] {\small $y=x+1$};
% Curve branches
\draw[popcurve, thick, domain=-4:0.9, samples=200] plot (\x, {(\x*\x+1)/(\x-1)});
\draw[popcurve, thick, domain=1.1:5, samples=200] plot (\x, {(\x*\x+1)/(\x-1)});
% Points
\fill[popblue] (0,-1) circle (2pt) node[below left] {$f(0)=-1$};
\fill[popblue] (1-1.414, -0.828) circle (2pt) node[above left] {Max};
\fill[popblue] (1+1.414, 3.828) circle (2pt) node[below right] {Min};
\end{tikzpicture}
\legende{Courbe de $f(x)=\frac{x^2+1}{x-1}$ avec asymptotes.}
\end{popfigure}
\end{exemplebox}

\begin{attention}[Erreurs fréquentes]
\begin{itemize}
  \item Oublier d'exclure la valeur interdite de $D_f$ \textbf{avant} tout calcul.
  \item Confondre asymptote verticale (limite infinie en un point fini) et asymptote oblique (comportement à l'infini).
  \item Ne pas donner la position relative par rapport à l'asymptote oblique (signe du reste).
  \item Faire le tableau de variations sur $\mathbb{R}$ au lieu de le faire sur les \textbf{deux intervalles} de $D_f$ séparément.
  \item Oublier que le dénominateur de $f'$ est un carré (toujours positif), le signe ne dépend que du numérateur.
\end{itemize}
\end{attention}

\courssub{Étude des fonctions du type $\ln(ax+b)$}

\begin{definition}[Fonction logarithme népérien affine]
Soient $a,b \in \mathbb{R}$, $a \neq 0$. La fonction $f(x) = \ln(ax+b)$ est définie si et seulement si l'argument est strictement positif :
\[ D_f = \{ x \in \mathbb{R} \mid ax+b > 0 \} \]
\begin{itemize}
  \item Si $a > 0$ : $D_f = ]-\frac{b}{a}, +\infty[$.
  \item Si $a < 0$ : $D_f = ]-\infty, -\frac{b}{a}[$.
\end{itemize}
La borne finie de $D_f$ ($-\frac{b}{a}$) n'appartient \textbf{pas} à $D_f$.
\end{definition}

\begin{propriete}[Limites aux bornes de $D_f$]
Rappel : $\lim_{u \to 0^+} \ln u = -\infty$, $\lim_{u \to +\infty} \ln u = +\infty$.
\begin{itemize}
  \item Cas $a > 0$ ($D_f = ]-\frac{b}{a}, +\infty[$) :
    $\lim_{x \to (-\frac{b}{a})^+} f(x) = -\infty$ (A.V. $x = -\frac{b}{a}$),
    $\lim_{x \to +\infty} f(x) = +\infty$.
  \item Cas $a < 0$ ($D_f = ]-\infty, -\frac{b}{a}[$) :
    $\lim_{x \to -\infty} f(x) = +\infty$,
    $\lim_{x \to (-\frac{b}{a})^-} f(x) = -\infty$ (A.V. $x = -\frac{b}{a}$).
\end{itemize}
Il n'y a \textbf{jamais} d'asymptote horizontale ni oblique (croissance logarithmique).
\end{propriete}

\begin{propriete}[Dérivée et variations]
\[ f'(x) = \frac{a}{ax+b} \]
Sur $D_f$, le dénominateur $ax+b$ est strictement positif. Donc \textbf{le signe de $f'$ est constant et égal au signe de $a$}.
\begin{itemize}
  \item Si $a > 0$ : $f' > 0$ sur $D_f$, $f$ est \textbf{strictement croissante} sur $]-\frac{b}{a}, +\infty[$.
  \item Si $a < 0$ : $f' < 0$ sur $D_f$, $f$ est \textbf{strictement décroissante} sur $]-\infty, -\frac{b}{a}[$.
\end{itemize}
Il n'y a \textbf{aucun extremum} local.
\end{propriete}

\begin{methode}[Étude complète de $f(x) = \ln(ax+b)$]
\begin{enumerate}
  \item Déterminer $D_f$ en résolvant $ax+b > 0$ (inéquation stricte).
  \item Calculer les limites aux deux bornes de $D_f$ (une finie $\to -\infty$, une infinie $\to \pm\infty$). Identifier l'A.V.
  \item Calculer $f'(x) = \frac{a}{ax+b}$, conclure sur la monotonie stricte (signe de $a$).
  \item Tableau de variations (une seule flèche par intervalle).
  \item Point(s) remarquable(s) : intersection avec l'axe des abscisses $\ln(ax+b)=0 \iff ax+b=1 \iff x=\frac{1-b}{a}$ (si dans $D_f$). Intersection avec l'axe des ordonnées $f(0)=\ln b$ (si $b>0$).
  \item Tracer la courbe : asymptote verticale, passage par les points trouvés, courbure (concave/convexe selon $f''$ si demandé).
\end{enumerate}
\end{methode}

\begin{exemplebox}[Étude de $f(x) = \ln(2x-3)$]
$D_f : 2x-3 > 0 \iff x > \frac{3}{2}$. Donc $D_f = ]\frac{3}{2}, +\infty[$.
A.V. $x = \frac{3}{2}$. $\lim_{x\to \frac{3}{2}^+} f(x) = -\infty$. $\lim_{x\to +\infty} f(x) = +\infty$.
$f'(x) = \frac{2}{2x-3} > 0$ sur $D_f$. $f$ strictement croissante.
$f(x)=0 \iff 2x-3=1 \iff x=2 \in D_f$.
$f(0)$ non défini ($0 \notin D_f$).
\begin{center}
\begin{tabular}{|c|ccccc|}\hline
$x$ & $\frac{3}{2}$ & & $2$ & & $+\infty$ \\ \hline
$f'(x)$ & & $+$ & $+$ & $+$ & \\ \hline
$f(x)$ & $-\infty$ & $\nearrow$ & $0$ & $\nearrow$ & $+\infty$ \\ \hline
\end{tabular}
\end{center}
\tcblower
\begin{popfigure}
\begin{tikzpicture}[scale=0.7, domain=1.51:5]
\draw[->] (1,0) -- (6,0) node[right] {$x$};
\draw[->] (0,-3) -- (0,3) node[above] {$y$};
\foreach \x in {1.5,2,3,4,5} \draw (\x,0.1) -- (\x,-0.1) node[below] {\small $\x$};
\foreach \y in {-2,-1,1,2} \draw (0.1,\y) -- (-0.1,\y) node[left] {\small $\y$};
\draw[dashed, poporange] (1.5,-3) -- (1.5,3) node[above] {\small $x=1,5$};
\draw[popcurve, thick, domain=1.51:5, samples=200] plot (\x, {ln(2*\x-3)});
\fill[popblue] (2,0) circle (2pt) node[below right] {$2$};
\end{tikzpicture}
\legende{Courbe de $f(x)=\ln(2x-3)$.}
\end{popfigure}
\end{exemplebox}

\begin{attention}[Points de vigilance]
\begin{itemize}
  \item L'inéquation $ax+b > 0$ est \textbf{stricte} : la borne n'est pas incluse. Écrire $] \frac{3}{2}, +\infty[$ et non $[\frac{3}{2}, +\infty[$.
  \item La dérivée est $\frac{a}{ax+b}$, pas $\frac{1}{ax+b}$ (oublier le $a$ de la dérivée de $ax+b$).
  \item Le tableau de variations ne comporte qu'\textbf{une seule colonne de variation} (pas d'extremum).
  \item Vérifier que les racines trouvées ($f(x)=0$) appartiennent bien à $D_f$.
\end{itemize}
\end{attention}

\courssub{Étude des fonctions du type $e^{ax+b}$}

\begin{definition}[Fonction exponentielle affine]
Soient $a,b \in \mathbb{R}$, $a \neq 0$. La fonction $f(x) = e^{ax+b}$ (ou $\exp(ax+b)$) est définie sur :
\[ D_f = \mathbb{R} \]
Ses valeurs sont \textbf{strictement positives} : $f(x) > 0$ pour tout $x \in \mathbb{R}$.
\end{definition}

\begin{propriete}[Limites aux bornes de $D_f$]
Rappel : $\lim_{u \to -\infty} e^u = 0$, $\lim_{u \to +\infty} e^u = +\infty$.
\begin{itemize}
  \item Cas $a > 0$ :
    $\lim_{x \to -\infty} f(x) = 0$ (A.H. $y=0$ à gauche),
    $\lim_{x \to +\infty} f(x) = +\infty$.
  \item Cas $a < 0$ :
    $\lim_{x \to -\infty} f(x) = +\infty$,
    $\lim_{x \to +\infty} f(x) = 0$ (A.H. $y=0$ à droite).
\end{itemize}
L'asymptote horizontale est toujours l'axe des abscisses ($y=0$) du côté où la limite est $0$.
\end{propriete}

\begin{propriete}[Dérivée et variations]
\[ f'(x) = a e^{ax+b} \]
Comme $e^{ax+b} > 0$ toujours, \textbf{le signe de $f'$ est constant et égal au signe de $a$}.
\begin{itemize}
  \item Si $a > 0$ : $f' > 0$ sur $\mathbb{R}$, $f$ \textbf{strictement croissante}.
  \item Si $a < 0$ : $f' < 0$ sur $\mathbb{R}$, $f$ \textbf{strictement décroissante}.
\end{itemize}
Aucun extremum local.
\end{propriete}

\begin{methode}[Étude complète de $f(x) = e^{ax+b}$]
\begin{enumerate}
  \item $D_f = \mathbb{R}$.
  \item Limites en $-\infty$ et $+\infty$ selon le signe de $a$. Identifier l'A.H. $y=0$ d'un côté.
  \item Calculer $f'(x) = a e^{ax+b}$, conclure sur la monotonie stricte (signe de $a$).
  \item Tableau de variations sur $\mathbb{R}$ (une seule flèche).
  \item Points remarquables : $f(0) = e^b$ (ordonnée à l'origine). $f(x)=0$ n'a \textbf{jamais} de solution (courbe au-dessus de l'axe des abscisses).
  \item Tracer : asymptote horizontale $y=0$ (pointillés), courbe strictement monotone, convexe ($f''>0$), passant par $(0, e^b)$.
\end{enumerate}
\end{methode}

\begin{exemplebox}[Étude de $f(x) = e^{2-x} = e^{-x+2}$]
Ici $a=-1, b=2$.
$D_f = \mathbb{R}$.
$a<0$ : $\lim_{x\to -\infty} f(x) = +\infty$, $\lim_{x\to +\infty} f(x) = 0$. A.H. $y=0$ à droite.
$f'(x) = -e^{2-x} < 0$ sur $\mathbb{R}$. $f$ strictement décroissante.
$f(0) = e^2 \approx 7,39$.
\begin{center}
\begin{tabular}{|c|ccccc|}\hline
$x$ & $-\infty$ & & $0$ & & $+\infty$ \\ \hline
$f'(x)$ & & $-$ & $-$ & $-$ & \\ \hline
$f(x)$ & $+\infty$ & $\searrow$ & $e^2$ & $\searrow$ & $0$ \\ \hline
\end{tabular}
\end{center}
\tcblower
\begin{popfigure}
\begin{tikzpicture}[scale=0.7, domain=-1:5]
\draw[->] (-1.5,0) -- (6,0) node[right] {$x$};
\draw[->] (0,-0.5) -- (0,8) node[above] {$y$};
\foreach \x in {1,2,3,4,5} \draw (\x,0.1) -- (\x,-0.1) node[below] {\small $\x$};
\foreach \y in {1,2,4,6,8} \draw (0.1,\y) -- (-0.1,\y) node[left] {\small $\y$};
\draw[dashed, poporange] (-1.5,0) -- (6,0) node[right] {\small $y=0$};
\draw[popcurve, thick, domain=-1:5, samples=200] plot (\x, {exp(2-\x)});
\fill[popblue] (0,7.389) circle (2pt) node[left] {$e^2$};
\end{tikzpicture}
\legende{Courbe de $f(x)=e^{2-x}$.}
\end{popfigure}
\end{exemplebox}

\begin{savaistu}[Croissances comparées : Hiérarchie à l'infini]
C'est un outil décisif pour les limites de formes indéterminées ($\frac{\infty}{\infty}$, $\infty-\infty$) et les asymptotes.
\begin{center}
\begin{tabular}{|c|c|}\hline
\rowcolor{popblueL} \textbf{Comparaison} & \textbf{Résultat} \\ \hline
Puissance vs Logarithme & $\forall \alpha > 0,\ \lim_{x\to+\infty} \frac{\ln x}{x^\alpha} = 0$ \\ \hline
Exponentielle vs Puissance & $\forall \alpha > 0,\ \lim_{x\to+\infty} \frac{x^\alpha}{e^x} = 0$ \\ \hline
Exponentielle vs Logarithme & $\lim_{x\to+\infty} \frac{\ln x}{e^x} = 0$ \\ \hline
\end{tabular}
\end{center}
\textbf{Mnémotechnique} : « L'exponentielle gagne contre tout, le logarithme perd contre tout, les puissances sont au milieu. »
\emph{Application} : Pour $f(x) = \frac{\ln x}{x}$, limite $0$ en $+\infty$ $\Rightarrow$ A.H. $y=0$. Pour $f(x) = x e^{-x}$, limite $0$ $\Rightarrow$ A.H. $y=0$.
\end{savaistu}

\courssub{Synthèse : Méthode générale d'étude et comparaison des croissances}

\begin{popfigure}
\begin{tikzpicture}[mindmap, grow cyclic, every node/.style={concept}, concept color=popblue!60,
  level 1/.append style={level distance=3.6cm, sibling angle=60},
  level 2/.append style={level distance=2.2cm},
  text width=2.4cm, align=center]
\node[concept color=popdark!80, text=white, text width=2.8cm]{\textbf{Étude de fonctions}}
  child[concept color=popblue!70]{ node {Polynômes degré $\leq 3$}
    child[level 2/.append style={sibling angle=50}]{ node[concept color=popblueL, text width=1.9cm]{$f'$ polynôme degré $2$} }
    child[level 2/.append style={sibling angle=-50}]{ node[concept color=popblueL, text width=1.9cm]{$D_f=\mathbb{R}$, pas d'asymptote} }
  }
  child[concept color=poporange!70]{ node {Rationnelles $\frac{ax^2+bx+c}{dx+e}$}
    child[level 2/.append style={sibling angle=50}]{ node[concept color=poporangeL, text width=1.9cm]{A.V. $x=-e/d$} }
    child[level 2/.append style={sibling angle=-50}]{ node[concept color=poporangeL, text width=1.9cm]{A.O. par division euclidienne} }
  }
  child[concept color=popgreen!70]{ node {Logarithme $\ln(ax+b)$}
    child[level 2/.append style={sibling angle=50}]{ node[concept color=popgreenL, text width=1.9cm]{$D_f : ax+b>0$} }
    child[level 2/.append style={sibling angle=-50}]{ node[concept color=popgreenL, text width=1.9cm]{A.V. en borne finie} }
  }
  child[concept color=poppurple!70]{ node {Exponentielle $e^{ax+b}$}
    child[level 2/.append style={sibling angle=50}]{ node[concept color=poppurpleL, text width=1.9cm]{Toujours $>0$} }
    child[level 2/.append style={sibling angle=-50}]{ node[concept color=poppurpleL, text width=1.9cm]{A.H. $y=0$ d'un côté} }
  }
  child[concept color=popgold!80]{ node {Limites et asymptotes}
    child[level 2/.append style={sibling angle=50}]{ node[concept color=popgoldL, text width=1.9cm]{Croissances comparées} }
    child[level 2/.append style={sibling angle=-50}]{ node[concept color=popgoldL, text width=1.9cm]{Formes indéterminées} }
  }
  child[concept color=poppink!70]{ node {Modélisation et optimisation}
    child[level 2/.append style={sibling angle=50}]{ node[concept color=poppinkL, text width=1.9cm]{Extremum $\Leftrightarrow f'$ s'annule} }
    child[level 2/.append style={sibling angle=-50}]{ node[concept color=poppinkL, text width=1.9cm]{Contexte réel : coût, volume} }
  };
\end{tikzpicture}
\legende{Carte mentale de la leçon : les quatre familles de fonctions et les deux compétences transversales.}
\end{popfigure}

\begin{aretenir}[L'essentiel de la leçon : Formules de référence]
\begin{center}
\begin{tabularx}{\linewidth}{|l|X|X|X|}
\hline
\rowcolor{popblueL} \textbf{Fonction $f(x)$} & \textbf{$D_f$} & \textbf{Dérivée $f'(x)$} & \textbf{Limites / Asymptotes clés} \\
\hline
$ax^3+bx^2+cx+d$ ($a\neq 0$) & $\mathbb{R}$ & $3ax^2+2bx+c$ & $\pm\infty$ selon signe $a$ ; Pas d'asymptote \\
\hline
$\dfrac{ax^2+bx+c}{dx+e}$ ($d\neq 0$) & $\mathbb{R}\setminus\{-e/d\}$ & $\dfrac{adx^2+2aex+(be-cd)}{(dx+e)^2}$ & A.V. $x=-e/d$ ; A.O. $y=\frac{a}{d}x+\frac{bd-ae}{d^2}$ \\
\hline
$\ln(ax+b)$ ($a\neq 0$) & $ax+b>0$ & $\dfrac{a}{ax+b}$ & A.V. en borne finie ($-\infty$) ; $+\infty$ autre borne \\
\hline
$e^{ax+b}$ ($a\neq 0$) & $\mathbb{R}$ & $a\,e^{ax+b}$ & A.H. $y=0$ d'un côté ; $+\infty$ de l'autre \\
\hline
\end{tabularx}
\end{center}
\end{aretenir}

\begin{propriete}[Théorèmes clés pour le Bac]
\begin{itemize}
  \item \textbf{Théorème des variations} : Si $f$ est dérivable sur un intervalle $I$, alors $f' \geq 0 \iff f$ croissante ; $f' \leq 0 \iff f$ décroissante. (Stricte si $f'$ ne s'annule pas sur un intervalle non réduit à un point).
  \item \textbf{Théorème des valeurs intermédiaires (TVI)} : Si $f$ continue sur $[a,b]$, pour tout $k$ entre $f(a)$ et $f(b)$, $\exists c \in [a,b], f(c)=k$. Utile pour montrer existence d'une solution à $f(x)=k$.
  \item \textbf{Croissances comparées} (exigible) :
    $\lim_{x\to+\infty} \frac{\ln x}{x} = 0$,
    $\lim_{x\to+\infty} \frac{e^x}{x} = +\infty$,
    $\lim_{x\to+\infty} \frac{x^n}{e^x} = 0\ (n\in\mathbb{N})$.
  \item \textbf{Asymptotes} :
    \begin{itemize}
      \item Verticale $x=x_0$ : $\lim_{x\to x_0^\pm} f(x) = \pm\infty$.
      \item Horizontale $y=\ell$ : $\lim_{x\to\pm\infty} f(x) = \ell \in \mathbb{R}$.
      \item Oblique $y=mx+p$ : $\lim_{x\to\pm\infty} (f(x)-mx-p) = 0$ avec $m = \lim \frac{f(x)}{x}$, $p = \lim (f(x)-mx)$.
    \end{itemize}
\end{itemize}
\end{propriete}

\begin{methode}[Méthode express d'étude complète (Checklist universelle)]
\begin{enumerate}
  \item \textbf{Ensemble de définition $D_f$} : Dénominateur $\neq 0$ ; Argument $\ln > 0$ ; Racine carrée $\geq 0$.
  \item \textbf{Limites aux bornes de $D_f$} (pas seulement $\pm\infty$ !) et \textbf{Asymptotes} (V, H, O).
  \item \textbf{Dérivée $f'$} : Formules usuelles, règles de dérivation (somme, produit, quotient, composée). Simplifier $f'$.
  \item \textbf{Signe de $f'$} : Résoudre $f'(x)=0$, tableau de signes de $f'$ sur $D_f$.
  \item \textbf{Tableau de variations} : Une ligne $f'$, une ligne $f$. Colonnes = bornes de $D_f$ + zéros de $f'$. Flèches $\nearrow/\searrow$ uniquement sur ligne $f$.
  \item \textbf{Points remarquables} : Extremums (images des zéros de $f'$), intersections axes ($f(0)$, $f(x)=0$).
  \item \textbf{Tracé} : Repère, asymptotes (pointillés), points, courbe respectant variations et convexité si connue.
\end{enumerate}
\end{methode}

\begin{aretenir}[Checklist avant le Bac — « Zéro faute »]
Avant de rendre ta copie, vérifie chaque point :
\begin{itemize}
  \item[$\square$] J'ai déterminé l'\cle{ensemble de définition} avant tout calcul (et je l'ai écrit en notation d'intervalle).
  \item[$\square$] Pour $\ln(ax+b)$, j'ai résolu l'inéquation \textbf{stricte} $ax+b>0$ (attention au sens si $a<0$).
  \item[$\square$] Pour une fonction rationnelle, j'ai exclu la valeur qui annule le dénominateur de $D_f$.
  \item[$\square$] J'ai calculé les limites aux \textbf{bornes ouvertes} de $D_f$ (ex: $1^-$ et $1^+$ pour A.V. en $1$), pas seulement en $\pm\infty$.
  \item[$\square$] J'ai dérivé correctement : $(\ln u)'=\frac{u'}{u}$, $(e^u)'=u'e^u$, $(\frac{u}{v})'=\frac{u'v-uv'}{v^2}$.
  \item[$\square$] J'ai factorisé/simplifié $f'$ pour en étudier le signe facilement.
  \item[$\square$] Le signe de $f'$ sur le tableau est cohérent avec mes calculs (zéros, signe au-delà).
  \item[$\square$] Les limites écrites dans le tableau de variations correspondent exactement aux valeurs calculées à l'étape 2.
  \item[$\square$] J'ai identifié et \textbf{nommé} toutes les asymptotes (Verticale, Horizontale, Oblique) avec leurs équations.
  \item[$\square$] Pour l'A.O., j'ai précisé la position relative (courbe au-dessus / en-dessous) grâce au signe du reste.
  \item[$\square$] J'ai placé les points remarquables (extremums, intersections axes) avec leurs coordonnées exactes ou approchées.
  \item[$\square$] Ma courbe respecte le tableau de variations et \textbf{ne traverse jamais ses asymptotes verticales}.
  \item[$\square$] En optimisation, j'ai justifié qu'il s'agit bien d'un maximum ou minimum global (par variations ou TVI).
  \item[$\square$] J'ai répondu à la question posée avec une \textbf{phrase complète} et les \textbf{unités} du contexte.
\end{itemize}
\end{aretenir}

\courssub{Problèmes de modélisation et optimisation}

\begin{methode}[Démarche de modélisation / optimisation]
\begin{enumerate}
  \item \textbf{Variables} : Identifier la variable indépendante $x$ (quantité à choisir) et la fonction objectif $f(x)$ (coût, surface, volume, bénéfice, temps...). Préciser les unités.
  \item \textbf{Domaine réel} : Déterminer l'intervalle $I \subset D_f$ qui a du sens physiquement ($x \geq 0$, $x \leq$ capacité max, etc.).
  \item \textbf{Étude mathématique} : Étudier $f$ sur $I$ (dérivée, variations, extremums).
  \item \textbf{Conclusion} : Identifier l'optimum (max ou min) sur $I$. Vérifier s'il est atteint à l'intérieur (racine de $f'$) ou au bord de $I$.
  \item \textbf{Réponse} : Phrase finale : « Le coût minimal est de ... FCFA pour une production de ... unités. »
\end{enumerate}
\end{methode}

\begin{exoresolu}[Optimisation d'un colis postal (Type Bac)]
\textbf{Énoncé :} La CAMPOST accepte les colis rectangulaires dont la somme de la longueur $L$ et du périmètre de la section carrée de côté $x$ ne dépasse pas $1,50$ m. On veut maximiser le volume $V$ du colis.
\begin{enumerate}
  \item Exprimer $V$ en fonction de $x$. Déterminer l'intervalle de valeurs possibles pour $x$.
  \item Étudier la fonction $V$ sur cet intervalle.
  \item Quelles dimensions donnent le volume maximal ? Quelle est cette valeur ?
\end{enumerate}
\tcblower
\textbf{Solution :}
\begin{enumerate}
  \item Contrainte : $L + 4x \leq 1,50 \Rightarrow L = 1,50 - 4x$.
  Volume $V(x) = x^2 L = x^2(1,50 - 4x) = 1,50 x^2 - 4x^3$.
  Conditions physiques : $x > 0$ et $L > 0 \Rightarrow 1,50 - 4x > 0 \Rightarrow x < 0,375$.
  Donc $x \in ]0, 0,375[$. On étudie sur $[0, 0,375]$ (continuité).
  \item $V(x) = -4x^3 + 1,5x^2$. $V'(x) = -12x^2 + 3x = 3x(1 - 4x)$.
  $V'(x)=0 \Rightarrow x=0$ (bord) ou $x=0,25 \in ]0, 0,375[$.
  Signe de $V'$ : $+$ sur $]0, 0,25[$, $-$ sur $]0,25, 0,375[$.
  \begin{center}
  \begin{tabular}{|c|ccccc|}\hline
  $x$ & $0$ & & $0,25$ & & $0,375$ \\ \hline
  $V'(x)$ & & $+$ & $0$ & $-$ & \\ \hline
  $V(x)$ & $0$ & $\nearrow$ & $V(0,25)$ & $\searrow$ & $0$ \\ \hline
  \end{tabular}
  \end{center}
  $V(0,25) = 1,5(0,0625) - 4(0,015625) = 0,09375 - 0,0625 = 0,03125\ \text{m}^3 = 31,25\ \text{L}$.
  \item Volume maximal pour $x = 0,25\ \text{m} = 25\ \text{cm}$ et $L = 1,50 - 1 = 0,50\ \text{m} = 50\ \text{cm}$.
  Volume max $= 31,25\ \text{L}$.
\end{enumerate}
\end{exoresolu}

\begin{exoresolu}[Modélisation exponentielle : Concentration d'un médicament]
\textbf{Énoncé :} Après une injection, la concentration $C(t)$ (en mg/L) dans le sang au temps $t$ (en heures) est modélisée par $C(t) = 2(e^{-0,2t} - e^{-0,5t})$ pour $t \geq 0$.
\begin{enumerate}
  \item Étudier les variations de $C$ sur $[0, +\infty[$. Donner la concentration maximale et l'instant où elle est atteinte.
  \item Déterminer à partir de quand la concentration devient inférieure à $0,1$ mg/L (seuil d'efficacité).
\end{enumerate}
\tcblower
\textbf{Solution :}
\begin{enumerate}
  \item $C(t) = 2e^{-0,2t} - 2e^{-0,5t}$, définie et dérivable sur $[0, +\infty[$.
  Limites : $C(0)=0$ et $\lim_{t\to+\infty} C(t) = 0$ (asymptote horizontale $y=0$). Pour $t>0$, $e^{-0,2t} > e^{-0,5t}$, donc $C(t)>0$.

  $C'(t) = -0,4e^{-0,2t} + e^{-0,5t}$. Comme $e^{-0,2t} = e^{-0,5t}\times e^{0,3t}$, on factorise :
  \[ C'(t) = e^{-0,5t}\left(1 - 0,4\,e^{0,3t}\right). \]
  \textbf{Attention au piège} : en factorisant par $e^{-0,5t}$, c'est le terme en $e^{-0,2t}$ qui reçoit le facteur $e^{0,3t}$.

  Comme $e^{-0,5t}>0$, $C'(t)$ a le signe de $1 - 0,4\,e^{0,3t}$ :
  \[ 1 - 0,4\,e^{0,3t} > 0 \iff e^{0,3t} < 2,5 \iff t < t_0 = \frac{\ln 2,5}{0,3} \approx 3,05. \]
  \begin{center}
  \begin{tabular}{|c|ccccc|}
  \hline
  $t$ & $0$ & & $t_0$ & & $+\infty$ \\ \hline
  $C'(t)$ & & $+$ & $0$ & $-$ & \\ \hline
  $C(t)$ & $0$ & $\nearrow$ & $C(t_0)$ & $\searrow$ & $0$ \\ \hline
  \end{tabular}
  \end{center}
  Au maximum, $e^{-0,3t_0} = 0,4$, donc $C(t_0) = 2e^{-0,2t_0}\left(1 - e^{-0,3t_0}\right) = 1,2\,e^{-0,2t_0} \approx 1,2\times 0,543 \approx 0,65$.

  \textbf{Conclusion :} la concentration est maximale environ $3$ h après l'injection ($t_0\approx 3,05$ h) et vaut environ $0,65$ mg/L.
  \item Sur $[0\,;t_0]$, $C$ est croissante : elle dépasse $0,1$ mg/L dès $t\approx 0,18$ h (environ $11$ min). Sur $[t_0\,;+\infty[$, $C$ est continue et strictement décroissante de $0,65$ vers $0$ : d'après le théorème des valeurs intermédiaires (bijection), l'équation $C(t)=0,1$ y admet une unique solution $t_1$.
  La calculatrice donne $C(14,9)\approx 0,1004 > 0,1$ et $C(15)\approx 0,0985 < 0,1$, d'où $t_1 \approx 14,9$ h.

  \textbf{Conclusion :} la concentration devient inférieure à $0,1$ mg/L environ $15$ heures après l'injection.
\end{enumerate}
\end{exoresolu}

\begin{exercice}[Entraînement Bac — Polynôme et Rationnelle]
\textbf{Exercice 1.} Soit $f(x) = \frac{x^3 - 2x^2 + x}{x-1}$ définie sur $\mathbb{R}^*_+$.
\begin{enumerate}
  \item Simplifier $f(x)$ pour $x \neq 1$. En déduire l'extension par continuité en $1$.
  \item Étudier les variations de $f$ sur $]0, +\infty[$.
  \item Tracer la courbe.
\end{enumerate}

\textbf{Exercice 2.} Soit $g(x) = \ln(x^2 - 4x + 5)$.
\begin{enumerate}
  \item Déterminer $D_g$. Montrer que $g$ est paire.
  \item Étudier les variations de $g$ sur $[0, +\infty[$.
  \item Résoudre $g(x) = \ln 2$.
\end{enumerate}
\end{exercice}

\begin{corrige}[Exercice 1]
\begin{enumerate}
  \item $x^3-2x^2+x = x(x^2-2x+1) = x(x-1)^2$.
  Pour $x \neq 1$, $f(x) = \frac{x(x-1)^2}{x-1} = x(x-1) = x^2 - x$.
  $\lim_{x\to 1} f(x) = 0$. Extension par continuité $\tilde{f}(1)=0$.
  \item Sur $]0, +\infty[$, $f(x) = x^2 - x$ (polynôme). $f'(x) = 2x-1$.
  $f'(x)=0 \iff x=0,5$.
  Variations : $\searrow$ sur $]0, 0,5]$, $\nearrow$ sur $[0,5, +\infty[$.
  Minimum $f(0,5) = -0,25$.
  \item Courbe parabole $y=x^2-x$ avec un trou en $(1, 0)$ (pointillé ouvert).
\end{enumerate}
\end{corrige}

\begin{corrige}[Exercice 2]
\begin{enumerate}
  \item $x^2-4x+5 = (x-2)^2+1 \geq 1 > 0$. $D_g = \mathbb{R}$.
  $g(-x) = \ln(x^2+4x+5) \neq g(x)$. \textbf{Attention : $g$ n'est pas paire.}
  $x^2-4x+5$ n'est pas pair. $g(2+x) = \ln(x^2+1) = g(2-x)$. Symétrie axiale $x=2$.
  \item Sur $[0, +\infty[$, $u(x)=x^2-4x+5$. $u'(x)=2x-4$.
  $g'(x) = \frac{2x-4}{x^2-4x+5}$. Signe de $g'$ = signe de $2x-4$.
  $g'$ négatif sur $[0,2]$, nul en $2$, positif sur $[2, +\infty[$.
  Minimum $g(2) = \ln 1 = 0$.
  \item $\ln(x^2-4x+5) = \ln 2 \iff x^2-4x+5 = 2 \iff x^2-4x+3=0 \iff (x-1)(x-3)=0$.
  Solutions : $x=1$ et $x=3$ (symétriques par rapport à $2$).
\end{enumerate}
\end{corrige}

\begin{exercice}[Approfondissement — Optimisation géométrique]
Un câble de longueur $L = 10$ m relie le sommet d'un pylône vertical (hauteur $H=6$ m) à un point d'ancrage au sol, à distance $x$ du pied du pylône. Le câble passe par une poulie fixée au sommet du pylône et soutient une charge au bout libre. On note $y$ la longueur du câble du pylône à la charge ($y = L - \sqrt{H^2+x^2}$). La tension dans le câble est proportionnelle à $\frac{1}{y}$. On veut minimiser la tension, donc maximiser $y$.
\begin{enumerate}
  \item Exprimer $y$ en fonction de $x$. Déterminer $D_y$ physique.
  \item Étudier les variations de $y(x)$. Pour quelle distance $x$ la tension est-elle minimale ?
  \item Calculer la tension minimale (valeur approchée) si la constante de proportionnalité est $k=500$ N$\cdot$m.
\end{enumerate}
\end{exercice}

\begin{corrige}[Exercice Approfondissement]
\begin{enumerate}
  \item $y(x) = 10 - \sqrt{36 + x^2}$.
  Contraintes : $x \geq 0$. Longueur câble pylône-charge $> 0 \Rightarrow \sqrt{36+x^2} < 10 \Rightarrow 36+x^2 < 100 \Rightarrow x^2 < 64 \Rightarrow x < 8$.
  $D_y = [0, 8[$.
  \item $y'(x) = -\frac{2x}{2\sqrt{36+x^2}} = -\frac{x}{\sqrt{36+x^2}}$.
  Sur $[0, 8[$, $y'(x) \leq 0$ (nul seulement en $0$). $y$ est strictement décroissante.
  Maximum de $y$ (donc tension min) atteint pour $x$ minimal, soit $x=0$.
  $y_{max} = y(0) = 10 - 6 = 4$ m.
  \item Tension $T = \frac{k}{y}$. $T_{min} = \frac{500}{4} = 125$ N.
  \textbf{Interprétation :} Pour minimiser la tension, il faut ancrer le câble au pied du pylône ($x=0$), le câble est alors vertical.
\end{enumerate}
\end{corrige}

\findecours

\end{document}
